% !TeX program = xelatex
% !TeX encoding = UTF-8
% 《晶体与 D-模》中文译本 —— Dennis Gaitsgory, Nick Rozenblyum 著
% 底本：arXiv:1111.2087v4 [math.AG], 2 Oct 2014
% 编译：xelatex -interaction=nonstopmode main.tex（三遍，生成目录与交叉引用）
\documentclass[UTF8, a4paper, 11pt, openany]{ctexbook}

% ---------- 宏包 ----------
\usepackage{amsmath, amssymb, amsthm, mathtools}
\usepackage{fontspec}
\usepackage{geometry}
\usepackage{fancyhdr}
\usepackage{tikz}
\usepackage{tikz-cd}
\usepackage{enumitem}
\usepackage{xcolor}
\usepackage{hyperref}

% ---------- 字体 ----------
% 本机只装了 Noto Serif SC / Noto Sans SC 的可变字体，没有独立 Bold/Italic 字面，
% 故用 AutoFakeBold / AutoFakeSlant 合成粗斜体。
\setCJKmainfont{Noto Serif SC}[AutoFakeBold=1.5, AutoFakeSlant=0.2]
\setCJKsansfont{Noto Sans SC}[AutoFakeBold=1.5]
\setCJKmonofont{Noto Sans SC}
\setmonofont{Consolas}[Scale=0.9]

% ---------- 版式 ----------
\geometry{left=2.8cm, right=2.6cm, top=2.8cm, bottom=2.6cm,
  headsep=0.6cm, footskip=0.9cm}
\setlength{\headheight}{14pt}
\pagestyle{fancy}
\fancyhf{}
\renewcommand{\headrulewidth}{0.4pt}
\fancyhead[LE]{\nouppercase{\leftmark}}
\fancyhead[RO]{\nouppercase{\rightmark}}
\fancyfoot[C]{\thepage}
\fancypagestyle{plain}{%
  \fancyhf{}\renewcommand{\headrulewidth}{0pt}\fancyfoot[C]{\thepage}}

\tolerance=1500
\emergencystretch=3em
\hfuzz=2pt

% 节标题：原书为左对齐加粗（中文书籍默认居中，此处改回左对齐）
\ctexset{section/format=\Large\bfseries\raggedright}

% ---------- 原书编号 ----------
% 原书的编号段与定理共用一套「节.序号」编号（如 2.3.1、2.3.2……），
% 且编号是显式印出的，不随排版浮动。故这里用显式编号宏：
%   \ob{2.3.1}{lab:2.3.1}          —— 编号段
%   \begin{othm}{定理}{2.4.4}{lab} —— 定理类环境
% 两者都写 \label，交叉引用仍走 \ref。
\makeatletter
\newcounter{oblk}
\newcommand{\ob}[2]{%
  \par\smallskip\refstepcounter{oblk}%
  \def\@currentlabel{#1}\label{#2}%
  \noindent\textbf{#1.}\ \ignorespaces}
\newenvironment{othm}[3]{%
  \par\smallskip\refstepcounter{oblk}%
  \def\@currentlabel{#2}\label{#3}%
  \noindent\textbf{#1 #2.}\ \ignorespaces}{\par\smallskip}
\makeatother

\renewcommand{\proofname}{证明}
\renewcommand{\qedsymbol}{$\square$}

% ---------- 校注（译者注） ----------
% 三类：易错预警、含糊点拨、背景直觉。每条须附来源等级。
\newenvironment{annote}{%
  \par\vspace{0.3em}\begingroup\small
  \noindent\textcolor{black!50}{\rule[-0.25em]{2.2pt}{1.1em}}\hspace{0.5em}%
  \textsf{\textbf{译注}}\hspace{0.4em}%
}{\par\endgroup\vspace{0.3em}}

% ---------- 常用记号 ----------
\newcommand{\dR}{\mathrm{dR}}
\newcommand{\PreStk}{\mathrm{PreStk}}
\newcommand{\DGSch}{\mathrm{DGSch}}
\newcommand{\DGIndSch}{\mathrm{DGIndSch}}
\newcommand{\DGSchaff}{\mathrm{DGSch}_{\mathrm{aff}}}
\newcommand{\Sch}{\mathrm{Sch}}
\newcommand{\Schaff}{\mathrm{Sch}_{\mathrm{aff}}}
\newcommand{\Crys}{\mathrm{Crys}}
\newcommand{\QCoh}{\mathrm{QCoh}}
\newcommand{\IndCoh}{\mathrm{IndCoh}}
\newcommand{\Coh}{\mathrm{Coh}}
\newcommand{\DGCat}{\mathrm{DGCat}}
\newcommand{\DGCatcont}{\mathrm{DGCat}_{\mathrm{cont}}}
\newcommand{\Funct}{\mathrm{Funct}}
\newcommand{\Maps}{\mathrm{Maps}}
\newcommand{\Hom}{\mathrm{Hom}}
\newcommand{\Spec}{\mathrm{Spec}}
\newcommand{\Grpd}{\mathrm{Grpd}}
\newcommand{\Vect}{\mathrm{Vect}}
\newcommand{\Tw}{\mathrm{Tw}}
\newcommand{\Ge}{\mathrm{Ge}}
\newcommand{\oblv}{\mathrm{oblv}}
\newcommand{\indf}{\mathrm{ind}}
\newcommand{\Tot}{\mathrm{Tot}}
\newcommand{\Ker}{\mathrm{Ker}}
\newcommand{\Cone}{\mathrm{Cone}}
\newcommand{\clred}{^{\mathrm{cl,red}}}
\newcommand{\laft}{\mathrm{laft}}
\newcommand{\aft}{\mathrm{aft}}
\newcommand{\onered}{\mathbf{1}}
\newcommand{\Gr}{\mathrm{Gr}}
\newcommand{\Pic}{\mathrm{Pic}}
\newcommand{\Aut}{\mathrm{Aut}}
\newcommand{\End}{\mathrm{End}}
\newcommand{\Ext}{\mathrm{Ext}}
\newcommand{\Diff}{\mathrm{Diff}}
\newcommand{\Sym}{\mathrm{Sym}}
\newcommand{\dmod}{\text{-}\mathrm{mod}}

\theoremstyle{plain}
\newtheorem*{plainthm}{定理}

% ---------- 超链接 ----------
\hypersetup{
  colorlinks=true, linkcolor=blue!70!black, citecolor=green!50!black,
  urlcolor=blue!70!black, bookmarks=true, bookmarksnumbered=true,
  pdftitle={晶体与 D-模（中文译本）},
  pdfauthor={Dennis Gaitsgory, Nick Rozenblyum}
}

\setcounter{tocdepth}{1}

\begin{document}

\frontmatter
\thispagestyle{empty}
\begin{center}
\vspace*{3cm}

{\Huge\bfseries 晶体与 D-模}

\vspace{1.2cm}

{\Large Dennis Gaitsgory\quad Nick Rozenblyum\ 著}

\vspace{0.6cm}

{\large 数译\ 译}

\vspace{2.4cm}

{\large\itshape 献给 A. Todorov 的纪念}

\vfill

{\small 底本：Dennis Gaitsgory, Nick Rozenblyum, \emph{Crystals and D-modules},\\
arXiv:1111.2087v4 [math.AG], 2014 年 10 月 2 日}

\vspace{0.6cm}
\end{center}
\clearpage

\chapter*{译者说明}
\addcontentsline{toc}{chapter}{译者说明}
\markboth{译者说明}{译者说明}

\textbf{底本。}本书译自 Dennis Gaitsgory 与 Nick Rozenblyum 的 \emph{Crystals and D-modules}（arXiv:1111.2087v4 [math.AG]，2014 年 10 月 2 日，共 65 页）。这是一个尚在展开的研究纲领的第一部分：作者在导出代数几何的框架里定义晶体，并把它与 D-模这一经典对象联系起来。第二部分（文中记作 [GR2]）当时尚未成文。

\textbf{编排。}译文保留原书的全部编号：章的编号、节的编号（如 2.3），以及节内逐条编号的段落与定理（如 2.3.1、2.3.2）。原书中编号段与定理共用一套编号，本书照此办理。公式编号（如 (0.1)、(2.5)）亦照原书。交叉引用一律指向这些原书编号，读者可与英文原文逐条对照。原书目录中的页码是 arXiv 版的页码，本书重排后不再保留；正文中「见 §2.3」一类指向的是节号，不受影响。

\textbf{术语。}术语的选定以原书用法为准，并尽量与本译丛已有的 D-模译法保持一致：holonomic 作「完整」，direct image 作「直像」，inverse image 作「逆像」，characteristic variety 作「特征簇」，good filtration 作「好滤过」。几个关键译名在此说明：crystal 作「晶体」，prestack 作「预栈」，gerbe 与 torsor 保留英文（首现处加注），twisting 作「扭化」，eventually coconnective 作「最终余连接的」，almost of finite type 作「几乎有限型的」。全书术语见书末附表。

\textbf{译注。}正文中以灰色小字排出的段落是译者所加，共三类：易错处的预警、原文含混处的点拨、背景与直觉的补充。译注不代替原文，也不改写原文的论证；凡涉及数学内容的补充，均注明依据。来源分三级：A 级为原书、作者本人论文、经典教材与原始论文；B 级为 nLab、Stacks Project、ProofWiki 等白名单资源（须与 A 级交叉印证）；I 级为本译丛既有的译法。凡无外部依据、仅由上下文推断者，一律标明「译者按（上下文推断，无外部依据）」。原文的排印错误（如拼写错误）径行改正，不另出注；凡可能影响数学理解的笔误，则保留原样并以译注说明。

\textbf{公式。}译文中的公式逐符号对照原书页面图像核校，而不依赖文本提取。原书的交换图以 Ti\emph{k}Z 重绘，位置与形状尽量与原书一致。

\vspace{1em}
\noindent 数译\\
2026 年 11 月
\clearpage

\tableofcontents

\mainmatter
\setcounter{chapter}{0}
\chapter*{引言}
\addcontentsline{toc}{chapter}{引言}
\markboth{引言}{引言}

\section{平坦连接、D-模与晶体}

\ob{0.1.1}{n:0.1.1}设 $M$ 是光滑流形，$V$ 是其上的向量丛。回忆：$V$ 上的平坦连接是满足 Leibniz 法则、且曲率 $[\nabla,\nabla]=0$ 的映射
\[
\nabla:V\to V\otimes\Omega^1_M .
\]
把连接映射对偶化，得到映射
\[
T_M\otimes V\to V .
\]
连接的平坦性保证 $V$ 由此成为向量场 Lie 代数上的模。等价地，$V$ 成为 $M$ 上微分算子代数 $\Diff_M$ 上的模。

这一概念立刻推广到特征零的光滑代数簇。在这样的簇上，D-模定义为微分算子层上的模，且作为 $\mathcal{O}$-模是拟凝聚的。于是代数簇上 D-模的概念推广了带平坦连接的向量丛，并编码了系数为多项式的线性微分方程组的资料。光滑代数簇上 D-模的研究是一个内容极其丰富的理论，在表示论等诸多领域有应用；带平坦连接的向量丛的微分几何中的许多想法，都可以搬到这一框架里。

然而，上述 D-模的进路有几处困难。例如，除光滑概形外，还需要考虑奇异概形上带平坦连接的层。微分算子代数在奇异簇上仍有定义，但它的模范畴并不是我们想要的那个范畴（比如，这个代数一般不是 Noether 的）。另一方面，即使对光滑代数簇，如何对非线性对象——例如层的范畴——定义连接，也并不清楚。

\ob{0.1.2}{n:0.1.2}\textbf{平行移动。}更好的定义来自微分几何中平坦连接的另一种解读，即平行移动。

设 $M$ 是光滑流形，其上带平坦连接的向量丛 $V$，并取道路 $\gamma:[0,1]\to M$。于是得到纤维之间的同构
\[
\Pi_\gamma:V_{\gamma(0)}\simeq V_{\gamma(1)},
\]
它只依赖道路的同伦类。这一构造可以换个说法。取 $M$ 内的小球 $B$。平行移动同构只依赖道路的同伦类，而 $B$ 可缩，所以沿 $B$ 内的道路，得到纤维的一个相容等同
\[
V_x\simeq V_y\qquad (x,y\in B).
\]
粗略地说，连接的数据给出了流形上「邻近」点处纤维的等同。

沿着这一想法，Grothendieck [Gr] 给出了平行移动这一概念的纯代数类比，并且本质地用到概形理论（而不只是簇）：他引入了概形 $X$ 的 $R$-点之间的无穷小接近关系。具体地说，两个 $R$-点 $x,y:\Spec(R)\to X$ 称为无穷小接近的，如果它们限制到 $\Spec(\mathrm{red}R)$ 上相同，其中 $\mathrm{red}R$ 是 $R$ 模去幂零根所得的商环。

概形 $X$ 上的晶体，按定义就是 $X$ 上关于无穷小接近关系等变的拟凝聚层。更确切地说，$X$ 上的晶体是一个拟凝聚层 $\mathcal{F}$，附带同构
\[
x^*(F)\simeq y^*(F)
\]
的数据，对任意两个无穷小接近的点 $x,y:\Spec(R)\to X$ 给出，并要求满足余循环条件。

Grothendieck 证明：在光滑代数簇上，晶体的 Abel 范畴等价于微分算子环上左模范畴。就此而言，晶体给出了带平坦连接的层的一个更基本的定义。

晶体范畴的一个突出特点是，Kashiwara 引理已经内置于它的定义之中：对概形的闭浸入 $i:Z\to X$，$Z$ 上晶体的范畴等价于 $X$ 上作为集合支撑在 $Z$ 上的晶体的范畴。这一观察使我们能把概形上晶体的研究化归到光滑概形的情形：把给定的概形（局部地）浸入一个光滑概形即可。

\ob{0.1.3}{n:0.1.3}本文在导出代数几何的框架里发展晶体的理论。那里考虑的不是通常的环，而是导出环，即 $E_\infty$ 环谱。由于在特征零的域 $k$ 上工作，我们取连接的交换 DG $k$-代数作为导出环的模型（相应地，使用术语「DG 概形」而不是「导出概形」）。关键的想法是：应当把导出环的高阶同伦群看作幂零元概念的推广。

于是，仿照 Simpson [Si]，对 DG 概形 $X$，定义它的 de Rham 预叠 $X_{\dR}$ 为导出环范畴上的函子
\[
X_{\dR}:R\mapsto X(\clred R),
\]
其中
\[
\clred R:=\mathrm{red}\big(\pi_0(R)\big)
\]
是与 $R$ 的底层经典环相对应的既约环。换言之，按 [GL:Stacks] 的术语，$X_{\dR}$ 是一个预叠。

我们定义 $X$ 上的晶体为预叠 $X_{\dR}$ 上的拟凝聚层。预叠上拟凝聚层的理论见 [Lu1, Sect.~2]，简要的回顾见 [GL:QCoh, Sect.~1.1]。

上述定义与前面提到的 Grothendieck 的定义并不一致：后者把映射 $\Spec(R)\to X$ 指定到等价类为止，前者则只指定映射 $\Spec(\clred R)\to X$。不过我们将证明：当 $X$ 最终余连接，即其结构环只有有限多个非零同伦群时，晶体的这两种定义等价。\footnote{当 $X$ 不是最终余连接时，这两种概念不同，正确的一种是经由 $X_{\dR}$ 定义的那一种。}

\ob{0.1.4}{n:0.1.4}晶体范畴虽然与 D-模范畴等价，却提供了一个更灵活的发展理论的框架。定义立刻推广到非光滑概形，相应的范畴性质良好（例如，任意概形上晶体的范畴是局部 Noether 的）。

\begin{annote}
为什么需要导出环？经典概形的 de Rham 预叠只看见既约部分，单看点集并不需要高阶同伦群。导出环的用处在于：沿非平坦映射取纤维积、取形式完备化、把晶体沿任意态射拉回时，结构环会出现幂零元；在导出代数几何中，这些幂零元自然由高阶同伦群承载。于是 connective DG 代数既保留了经典情形，又在取极限时封闭。\hfill【来源：原书 §0.1.3 及脚注 1；[GR1, Sect.~6.1]（A）】
\end{annote}


设 $f:X\to Y$ 是 DG 概形的映射。我们将构造自然的拉回函子
\[
f^\dagger:\Crys(Y)\to\Crys(X).
\]
事实上，我们将把 $X\mapsto\Crys(X)$ 延拓为从 $\DGSch^{\mathrm{op}}$ 到稳定 $\infty$-范畴的函子。后者使我们不只在 DG 概形上，而且在任意预叠上定义晶体。

此外，晶体的概念立刻推广到非线性与范畴化的情形：我们同样可以定义 $X$ 上的概形晶体或范畴晶体。

\section{左晶体与右晶体}

\ob{0.2.1}{n:0.2.1}回忆：在光滑代数簇 $X$ 上，除通常的（即左）D-模外，还可以考虑右 D-模的范畴。这两个范畴等价，相应的函子是沿函数环与对偶化线丛 $\omega_X$ 作张量积。然而这个等价不保持到拟凝聚层的遗忘函子。因此，我们可以考虑一个抽象的 D-模范畴，它带有两个不同的实现函子，分别到拟凝聚层。在左实现中，D-模的拉回函子成为拟凝聚层上的 $*$-拉回函子；在右实现中，它成为 $!$-拉回函子。

事实证明，「右」实现比「左」实现有几个优点。主要的优点也许是：「右」实现给 D-模范畴赋予了一个函子性非常好的 $t$-结构。特别地，在 Riemann--Hilbert 对应下，这个 $t$-结构成为反常 $t$-结构。

\ob{0.2.2}{n:0.2.2}于是可以问：在任意 DG 概形上，是否同样有「左」「右」晶体？答案是：两个范畴确实都可以非常一般地定义。

左晶体就是 §0.1.3 中定义的对象。但为了定义右晶体，需要把通常的拟凝聚层范畴换成它的重整化版本，即 [IndCoh] 中引入的归纳凝聚层范畴。

归纳凝聚层范畴 $\IndCoh(X)$ 对（几乎）局部有限型的（导出）概形性质良好，所以右晶体只在这样的 DG 概形上定义，进而只在此性质的预叠上定义。

回忆 [IndCoh, Sect.~5]：对 DG 概形之间的映射 $f:X\to Y$，有 $!$-拉回函子
\[
f^!:\IndCoh(Y)\to\IndCoh(X).
\]
附着 $X\mapsto\IndCoh(X)$ 是从 $\DGSch^{\mathrm{op}}$ 到稳定 $\infty$-范畴的函子，因而可以延拓为从预叠范畴出发的函子。

对 DG 概形 $X$，定义右晶体的范畴 $\Crys^r(X)$ 为 $\IndCoh(X_{\dR})$。也可以按 Grothendieck 的方式重新表述这个定义：$X$ 上的右晶体是对象 $F\in\IndCoh(X)$，附带对每对无穷小接近的点 $x,y:\Spec(R)\to X$ 的同构
\begin{equation}
x^!(F)\simeq y^!(F),\tag{0.1}
\end{equation}
并要求满足（$\infty$-范畴版本的）余循环条件。可以证明，与左晶体的情形不同，这确实给出了右晶体的等价定义，而不需要任何余连接假设。

\begin{annote}
为什么右晶体必须用归纳凝聚层？因为条件 (0.1) 要求对任意 $x,y$ 比较 $x^!(F)$ 与 $y^!(F)$，而 $!$-拉回一般只在归纳凝聚层上定义；$\QCoh$ 上的 $*$-拉回无法表达「支撑在无穷小邻域上」的信息。另一方面，$\Upsilon_X:\QCoh(X)\to\IndCoh(X)$ 在 $X$ 局部几乎有限型时把两种实现联系起来，故左右晶体最终仍然等价（§2.4）。\hfill【来源：原书 §0.2.2--0.2.3；[IndCoh, Sect.~4--5]（A）】
\end{annote}


\ob{0.2.3}{n:0.2.3}右晶体的范畴既已定义，就可以问它是否与左晶体的范畴等价。答案同样是肯定的。具体地说，对几乎有限型的任意 DG 概形 $X$，与对偶化复形 $\omega_X$ 作张量积给出函子
\[
\Upsilon_X:\QCoh(X)\to\IndCoh(X),
\]
它把拟凝聚层上的 $*$-拉回与归纳凝聚层上的 $!$-拉回联系起来。

虽然对单个 $S$，函子 $\Upsilon_S$ 一般不是等价（除非 $S$ 光滑），但对所有映入 $X$ 的 de Rham 预叠的 DG 概形取这些映射的整体，得到的却是左晶体与右晶体之间的等价。

于是，正如光滑簇的情形，可以认为对每个 DG 概形 $X$ 都附着一个范畴 $\Crys(X)$，它带两个「实现」函子
\[
\begin{tikzcd}
 & \Crys(X) \arrow[dl, "\oblv^l_X"'] \arrow[dr, "\oblv^r_X"] & \\
 \QCoh(X) \arrow[rr, "\Upsilon_X"'] & & \IndCoh(X)
\end{tikzcd}
\]
然而在非光滑概形的情形，「右」实现所对应的 $\Crys(X)$ 上 $t$-结构的优点更加明显。

\ob{0.2.4}{n:0.2.4}\textbf{历史注。}据我们所知，经由右晶体研究 D-模的进路最早由 A. Beilinson 在 90 年代初提出，当时是在 Abel 范畴的层面。

此后一段时间里，人们误以为不能用左晶体来定义 D-模，理由是 $t$-结构不相容。然而 J. Lurie 说明：只要放弃 $t$-结构、直接定义相应的稳定 $\infty$-范畴，左晶体同样适用。

\section{晶体/D-模的理论}

让我们说明本文及其续篇 [GR2] 所发展的理论的整体结构。

\ob{0.3.1}{n:0.3.1}对每个（局部几乎有限型的）预叠 $Y$，我们赋予一个稳定 $\infty$-范畴
\[
Y\rightsquigarrow\Crys(Y).
\]
这个范畴有两个实现函子：到 $\QCoh(Y)$ 的左实现函子，与到 $\IndCoh(Y)$ 的右实现函子。二者由如下交换图联系：
\[
\begin{tikzcd}
 & \Crys(Y) \arrow[dl, "\oblv^l_Y"'] \arrow[dr, "\oblv^r_Y"] & \\
 \QCoh(Y) \arrow[rr, "\Upsilon_Y"'] & & \IndCoh(Y)
\end{tikzcd}
\]
其中 $\Upsilon_Y$ 是与对偶化复形 $\omega_Y$ 作张量积给出的函子 $\QCoh(Y)\to\IndCoh(Y)$。

\ob{0.3.2}{n:0.3.2}把 $\Crys(Y)$ 附着到 $Y$ 上的做法在多种意义下具有函子性。对映射 $f:Y_1\to Y_2$，有拉回函子
\[
f^\dagger:\Crys(Y_2)\to\Crys(Y_1),
\]
它对 $f$ 具有函子性；也就是说，这一附着给出函子
\[
\Crys^\dagger_{\PreStk}:(\PreStk)^{\mathrm{op}}\to\DGCat_{\mathrm{cont}}.
\]
D-模上的拉回函子与实现函子相容。具体地说，有交换图
\[
\begin{tikzcd}
 \Crys(Y_1) & \Crys(Y_2) \arrow[l, "f^\dagger"'] \\
 \QCoh(Y_1) \arrow[u, "\oblv^l_{Y_1}"] & \QCoh(Y_2) \arrow[l, "f^*"] \arrow[u, "\oblv^l_{Y_2}"']
\end{tikzcd}
\]
以及
\[
\begin{tikzcd}
 \Crys(Y_1) & \Crys(Y_2) \arrow[l, "f^\dagger"'] \\
 \IndCoh(Y_1) \arrow[u, "\oblv^r_{Y_1}"] & \IndCoh(Y_2) \arrow[l, "f^!"] \arrow[u, "\oblv^r_{Y_2}"']
\end{tikzcd}
\]
此外，这种相容性本身对 $f$ 具有函子性；也就是说，有一个自然交换的函子图
\[
\begin{tikzcd}
 & \Crys^\dagger_{\PreStk} \arrow[dl, "\oblv^l"'] \arrow[dr, "\oblv^r"] & \\
 \QCoh_{\PreStk} \arrow[rr, "\Upsilon"'] & & \IndCoh^!_{\PreStk}
\end{tikzcd}
\]

\ob{0.3.3}{n:0.3.3}上述这部分理论在本文中构造。也就是说，本文关心的是附着
\[
Y\rightsquigarrow\Crys(Y)
\]
以及拉回运算。因此，本文发展的是晶体/D-模的局部理论。

不过，除函子 $f^\dagger$ 之外，我们还期待有推前函子 $f_{\dR,*}$，并且两者须满足各种相容关系。后者将在 [GR2] 中完成。这里先指出这套合并理论的主要构件。

\ob{0.3.4}{n:0.3.4}对预叠之间概形性的拟紧映射 $f:Y_1\to Y_2$，有 de Rham 推前函子
\[
f_{\dR,*}:\Crys(Y_1)\to\Crys(Y_2),
\]
它对 $f$ 具有函子性。这一附着给出另一个函子
\[
\big(\Crys_{\dR,*}\big)_{\PreStk^{\mathrm{sch\text{-}qc}}}:\PreStk^{\mathrm{sch\text{-}qc}}\to\DGCat_{\mathrm{cont}},
\]
其中 $\PreStk^{\mathrm{sch\text{-}qc}}$ 是 $\PreStk$ 的非全子范畴，其 1-态射限制为概形性的拟紧映射。

设 $Y=X$ 是 DG 概形。\footnote{更一般地，可以取 $Y$ 为具有形变理论的预栈。}此时遗忘函子
\[
\oblv^r_Y:\Crys(Y)\to\IndCoh(Y)
\]
有左伴随，记作
\[
\indf^r_Y:\IndCoh(Y)\to\Crys(Y),
\]
称为归纳函子。

归纳函子与 de Rham 推前相容。具体地说，有交换图
\[
\begin{tikzcd}
 \IndCoh(Y_1) \arrow[r, "f^{\IndCoh}_*"] \arrow[d, "\indf^r_{Y_1}"'] & \IndCoh(Y_2) \arrow[d, "\indf^r_{Y_2}"] \\
 \Crys(Y_1) \arrow[r, "f_{\dR,*}"'] & \Crys(Y_2)
\end{tikzcd}
\]
这种相容性本身也具有函子性；也就是说，有函子的自然变换
\[
\indf^r:(\IndCoh_*)_{\PreStk^{\mathrm{sch\text{-}qc}}}\to\big(\Crys_{\dR,*}\big)_{\PreStk^{\mathrm{sch\text{-}qc}}}.
\]

\ob{0.3.5}{n:0.3.5}当 $f$ 固有时，函子 $(f_{\dR,*},f^\dagger)$ 构成伴随对；当 $f$ 光滑时，若 $n$ 是 $f$ 的相对维数，则 $(f^\dagger[-2n],f_{\dR,*})$ 构成伴随对。

一般地，这两个函子并不伴随，但满足基变换公式。正如 [IndCoh, Sect.~5.1] 所解释的，编码基变换公式的函子性的一种方式是考虑对应范畴。具体地说，令 $(\PreStk)^{\mathrm{corr:all,sch\text{-}qc}}$ 是这样的 $\infty$-范畴：其对象是局部有限型的预叠，从 $Y_1$ 到 $Y_2$ 的态射由对应
\[
\begin{tikzcd}
 & Z \arrow[dl, "f"'] \arrow[dr, "g"] & \\
 Y_2 & & Y_1
\end{tikzcd}
\]
给出，其中 $f$ 概形性且拟紧，$g$ 任意。这个范畴中的复合由对应的纤维积给出。于是，函子 $\Crys^\dagger$ 与 $\Crys_{\dR,*}$ 的相容基变换公式是一个函子
\[
\Crys_{(\PreStk)^{\mathrm{corr:all,sch\text{-}qc}}}:(\PreStk)^{\mathrm{corr:all,sch\text{-}qc}}\to\DGCat_{\mathrm{cont}},
\]
连同它在 $(\PreStk)^{\mathrm{op}}$ 上的限制与 $\Crys^\dagger_{\PreStk}$ 的等同，以及在 $\PreStk^{\mathrm{sch\text{-}qc}}$ 上的限制与 $(\Crys_{\dR,*})_{\PreStk^{\mathrm{sch\text{-}qc}}}$ 的等同。

\section{扭化}

\ob{0.4.1}{n:0.4.1}除 D-模外，考虑扭 D-模往往也很重要。例如在表示论中，Beilinson 与 Bernstein 的局部化定理把具有固定中心特征 $\chi$ 的约化 Lie 代数 $\mathfrak{g}$ 的表示范畴，等同于旗簇 $G/B$ 上扭 D-模的范畴，其扭化由 $\chi$ 决定。

在光滑簇的情形，扭化与扭 D-模的理论由 Beilinson 与 Bernstein [BB] 引入。扭化的重要例子来自线丛的复张量幂。对光滑簇 $X$，扭化构成一个 Picard 群胚，它可以这样描述。令 $T$ 为次数 1 与 2 上的层复形
\[
T:=\Omega^1\to\Omega^{2,\mathrm{cl}},
\]
其中 $\Omega^1$ 是 $X$ 上 1-形式层，$\Omega^{2,\mathrm{cl}}$ 是闭 2-形式层，映射是 de Rham 微分。于是扭化的 Picard 群胚的对象空间由 $H^2(X,T)$ 给出；对一个给定对象，同构空间是 $H^1(X,T)$。

\ob{0.4.2}{n:0.4.2}本文最后两节在导出（因而特别是非光滑）的框架里发展扭化与扭晶体的理论。我们给出扭化概念的几种等价表述，并证明它们在光滑簇的情形与 [BB] 中的定义一致。

对（几乎局部有限型的）预叠 $Y$，定义扭化为 de Rham 预叠 $Y_{\dR}$ 上的 $\mathbb{G}_m$-gerbe，并附带其到 $Y$ 的拉回的平凡化。$Y$ 上的线丛 $L$ 给出一个扭化：$Y_{\dR}$ 上的 gerbe 平凡，而 $Y$ 上的平凡化由 $L$ 给出。

给定扭化 $T$，$Y$ 上 $T$-扭晶体的范畴定义为 $Y_{\dR}$ 上被 $T$ 所给的 $\mathbb{G}_m$-gerbe 扭化的（归纳凝聚或拟凝聚）层范畴。

\section{内容概览}

下面逐节说明本文的内容。

\ob{0.5.1}{n:0.5.1}在 §1 中，对预叠 $Y$ 定义 de Rham 预叠 $Y_{\dR}$，并建立它的一些基本性质。最重要的是证明：若 $Y$ 局部几乎有限型，则 $Y_{\dR}$ 也局部几乎有限型。

\ob{0.5.2}{n:0.5.2}在 §2 中，把左晶体定义为 de Rham 预叠上的拟凝聚层；在局部几乎有限型的情形，把右晶体定义为 de Rham 预叠上的归纳凝聚层。后者之所以有意义，是因为 §1 证明了：局部几乎有限型的预叠，其 de Rham 预叠也局部几乎有限型。在这一情形，我们证明左晶体与右晶体的范畴等价。此外还证明这一框架下的 Kashiwara 引理。

\ob{0.5.3}{n:0.5.3}在 §3 中，证明晶体范畴满足 h-下降（特别地，fppf 下降）。我们还引入预叠 $Y$ 的无穷小群胚，即自然映射 $Y\to Y_{\dR}$ 的 Čech 神经。具体地说，$Y$ 的无穷小群胚由
\[
(Y\times Y)^\wedge_Y\rightrightarrows Y
\]
给出，其中 $(Y\times Y)^\wedge_Y$ 是 $Y\times Y$ 沿对角线的形式完备化。

§3 的大部分内容专门讨论 $Y$ 是 ind-概形的情形。无穷小群胚上的层，就是 $Y$ 上关于无穷小接近等价关系等变的层。在归纳凝聚层的情形，这个范畴等价于右晶体。然而一般说来，无穷小群胚上的拟凝聚层并不等价于左晶体。我们证明：当 $Y$ 是最终余连接的 DG 概形或经典形式光滑的预叠时，$Y$ 的无穷小群胚上的拟凝聚层等价于左晶体。特别地，对经典概形，这个等价成立。

我们还定义从 $\QCoh(Y)$ 与 $\IndCoh(Y)$ 到 $Y$ 上晶体的归纳函子。在归纳凝聚层的情形，归纳函子是遗忘函子的左伴随，且右晶体的范畴等价于相应单子上的模范畴。对 $\QCoh$ 与左晶体，当 $Y$ 是最终余连接的 DG 概形时，类似结果成立。

\ob{0.5.4}{n:0.5.4}在 §4 中，证明晶体范畴有两个自然的 $t$-结构：一个与到 $\QCoh$ 的左实现相容，另一个与到 $\IndCoh$ 的右实现相容。对拟紧 DG 概形，这两个 $t$-结构相差一个有界幅度。

我们还证明：对仿射 DG 概形，晶体范畴等价于其关于右 $t$-结构的心上的导出范畴。

\ob{0.5.5}{n:0.5.5}在 §5 中，把 DG 概形上作用在 $\IndCoh$（相应地，$\QCoh$）上、负责右（相应地，左）晶体范畴的单子，与微分算子层联系起来。

作为推论，我们把晶体范畴与 D-模的导出范畴联系起来。

\ob{0.5.6}{n:0.5.6}在 §6 中，定义预叠 $Y$ 上扭化的 Picard 群胚为 de Rham 预叠 $Y_{\dR}$ 上、在 $Y$ 上平凡化的 $\mathbb{G}_m$-gerbe 的 Picard 群胚。然后给出这个定义的几种等价表述。例如，利用某一版本的指数映射，证明扭化的 Picard 群胚等价于 de Rham 预叠 $Y_{\dR}$ 上、在 $Y$ 上平凡化的 $\mathbb{G}_a$-gerbe 的 Picard 群胚。特别地，这使扭化自然成为一个 $k$-线性 Picard 群胚。

此外，利用扭化的 $\mathbb{G}_a$-gerbe 描述，我们把扭化的 $\infty$-群胚等同于
\[
\tau^{\le 2}\Big(H_{\dR}(Y)\times_{H(Y)}\{*\}\Big)[2],
\]
其中 $H_{\dR}(Y)$ 是 $Y$ 的 de Rham 上同调，$H(Y)$ 是 $Y$ 的凝聚上同调。特别地，对光滑经典概形，这表明这里定义的扭化概念与 [BB] 中的定义一致。

最后，我们证明：局部有限型的 DG（ind-）概形 $X$ 上扭化的范畴，等同于其无穷小群胚的中心扩张的范畴。

\ob{0.5.7}{n:0.5.7}在 §7 中，定义扭晶体的范畴并建立其基本性质。特别地，我们证明关于晶体的大部分结果都可以搬到扭化的情形。

\section{约定与记号}

我们的约定紧跟 [GR1]。这里回顾其中最必要的部分。

\ob{0.6.1}{n:0.6.1}\textbf{基域。}全文在固定的特征零基域 $k$ 上工作。

\ob{0.6.2}{n:0.6.2}\textbf{$\infty$-范畴。}所谓 $\infty$-范畴，总是指 $(\infty,1)$-范畴。我们有时把 $\infty$-范畴简称为「范畴」，这是轻微的语言滥用。我们对 $\infty$-范畴的使用不依赖模型，但心目中取的是它们作为准范畴的实现。准范畴意义下 $\infty$-范畴的基本参考书是 [Lu0]。

记 $\infty\text{-}\Grpd$ 为 $\infty$-群胚的 $\infty$-范畴，它等同于 [Lu0] 记号下的空间范畴 $\mathcal{S}$。

对 $\infty$-范畴 $\mathcal{C}$ 与 $x,y\in\mathcal{C}$，记 $\Maps_{\mathcal{C}}(x,y)\in\infty\text{-}\Grpd$ 为相应的映射空间。记 $\Hom_{\mathcal{C}}(x,y)$ 为集合 $\pi_0(\Maps_{\mathcal{C}}(x,y))$，即 [Lu0] 中记作 $\Hom_{h\mathcal{C}}(x,y)$ 的东西。

稳定 $\infty$-范畴 $\mathcal{C}$ 自然地充实于谱。此时对 $x,y\in\mathcal{C}$，记 $\Maps_{\mathcal{C}}(x,y)$ 为 $x$ 到 $y$ 的映射谱。特别地，映射空间等于该映射谱的 $\Omega^\infty$。

在固定的 $\infty$-范畴 $\mathcal{C}$ 中工作时，对两个对象 $x,y\in\mathcal{C}$，我们把 $\Maps_{\mathcal{C}}(x,y)$ 的点称为同构，而 [Lu0] 中称之为等价。也就是说，同构是有同伦逆的映射。我们保留「等价」一词，专指 $\infty$-范畴之间的（同伦）等价。

\ob{0.6.3}{n:0.6.3}\textbf{DG 范畴。}我们关于 DG 范畴的约定遵循 [IndCoh, Sect.~0.6.4]。所谓 DG 范畴，指的是 $k$ 上可表现的 DG 范畴；特别地，我们的 DG 范畴都假设余完备。除另有说明，我们只考虑 DG 范畴之间的连续函子（即与直和交换、等价地与所有余极限交换的精确函子）。换言之，我们在 [GL:DG] 记号下的范畴 $\DGCat_{\mathrm{cont}}$ 中工作。\footnote{可以把 $\DGCat_{\mathrm{cont}}$ 换成（等价的）稳定可表现 $\infty$-范畴的 $(\infty,1)$-范畴，这些范畴张量于 $\Vect$ 上，态射为保余极限的函子。}

记 $\Vect$ 为 $k$-向量空间复形的 DG 范畴。范畴 $\DGCat_{\mathrm{cont}}$ 有自然的对称单幺结构，其单位是 $\Vect$。

对带 $t$-结构的 DG 范畴 $\mathcal{C}$，记 $\mathcal{C}^{\le n}$（相应地，$\mathcal{C}^{\ge m}$、$\mathcal{C}^{\le n,\ge m}$）为 $\mathcal{C}$ 中由满足 $H^i(x)=0$（$i>n$；相应地，$i<m$；$(i>n)\wedge(i<m)$）的对象 $x$ 张成的全子范畴。包含 $\mathcal{C}^{\le n}\hookrightarrow\mathcal{C}$ 有右伴随，记作 $\tau^{\le n}$；其余类似。

有一个从 $\DGCat_{\mathrm{cont}}$ 到稳定 $\infty$-范畴与连续精确函子范畴的全忠实函子。如此得到的稳定 $\infty$-范畴充实于 $\Vect$。因此我们常把映射谱 $\Maps_{\mathcal{C}}(x,y)$ 看作 $\Vect$ 的对象；前者由后者经 Dold--Kan 对应得到。

\ob{0.6.4}{n:0.6.4}\textbf{（预）叠与 DG 概形。}我们关于（预）叠与 DG 概形的约定遵循 [GL:Stacks]：

记 $\DGSch_{\mathrm{aff}}$ 为 $k$ 上连接交换 DG 代数范畴的反范畴。

预叠的范畴 $\PreStk$ 按定义是全体函子
\[
(\DGSch_{\mathrm{aff}})^{\mathrm{op}}\to\infty\text{-}\Grpd
\]
构成的范畴。

记 $^{<\infty}\DGSch_{\mathrm{aff}}$ 为 $\DGSch_{\mathrm{aff}}$ 中由最终余连接对象张成的全子范畴。

回忆：最终余连接的仿射 DG 概形 $S=\Spec(A)$ 称为几乎有限型的，如果
\begin{itemize}[nosep]
\item $H^0(A)$ 在 $k$ 上有限型；
\item 每个 $H^i(A)$ 作为 $H^0(A)$-模有限生成。
\end{itemize}
记 $^{<\infty}\DGSch_{\mathrm{aff}}^{\mathrm{aft}}$ 为 $^{<\infty}\DGSch_{\mathrm{aff}}$ 中由几乎有限型的概形张成的全子范畴，并记 $\PreStk_{\mathrm{laft}}$ 为全体函子
\[
\big({}^{<\infty}\DGSch_{\mathrm{aff}}^{\mathrm{aft}}\big)^{\mathrm{op}}\to\infty\text{-}\Grpd
\]
构成的范畴。正如 [GL:Stacks, Sect.~1.3.11] 所解释的，$\PreStk_{\mathrm{laft}}$ 经由适当的 Kan 延拓自然地成为 $\PreStk$ 的子范畴。

为了应用 [IndCoh] 中发展的归纳凝聚层形式体系，本文大部分内容假设所考虑的预叠局部几乎有限型。遇到不如此假设的地方，我们会明确指出。

\ob{0.6.5}{n:0.6.5}\textbf{既约环。}记 $({}^{\mathrm{red}}\Sch_{\mathrm{aff}})^{\mathrm{op}}\subset(\DGSch_{\mathrm{aff}})^{\mathrm{op}}$ 为既约离散环的范畴。包含函子有自然的左伴随
\[
(-)^{\mathrm{cl,red}}:(\DGSch_{\mathrm{aff}})^{\mathrm{op}}\to({}^{\mathrm{red}}\Sch_{\mathrm{aff}})^{\mathrm{op}},
\]
由
\[
S\mapsto H^0(S)/\mathrm{nilp}(H^0(S))
\]
给出，其中 $\mathrm{nilp}(H^0(S))$ 是 $H^0(S)$ 中幂零元构成的理想。

\section{致谢}

我们感谢 Jacob Lurie 的多次有益讨论。他的想法对本论文影响如此之深，以至于很难一一指出哪些具体陈述是直接向他学到的。

D.G. 的研究由 NSF 基金 DMS-1063470 资助。

\chapter{de Rham 预叠}
\markboth{第 1 章\ de Rham 预叠}{第 1 章\ de Rham 预叠}

对预叠 $Y$，晶体定义为 $Y$ 的 de Rham 预叠 $Y_{\dR}$ 上的层（拟凝聚或归纳凝聚）。本节定义函子 $Y\mapsto Y_{\dR}$，并建立它的一些基本性质。

最重要的是，我们将证明：若 $Y$ 局部几乎有限型，则 $Y_{\dR}$ 也局部几乎有限型。此时我们还将证明 $Y_{\dR}$ 是经典的，即可以完全在「经典」代数几何的范围内研究它，无须涉及导出环。

读者可能觉得本节格外抽象。一个可行的策略是初读时跳过本节，等到本节建立的论断用于晶体时再回头查阅。

\section{定义与基本性质}

\ob{1.1.1}{n:1.1.1}设 $Y$ 是 $\PreStk$ 的对象。定义 $Y$ 的 de Rham 预叠 $Y_{\dR}\in\PreStk$ 为
\begin{equation}
Y_{\dR}(S):=Y(\clred S),\qquad S\in\DGSch_{\mathrm{aff}}.\tag{1.1}
\end{equation}

\ob{1.1.2}{n:1.1.2}更抽象地，可以改写为
\[
Y_{\dR}:=\mathrm{RKE}_{{}^{\mathrm{red}}\Sch_{\mathrm{aff}}\hookrightarrow\DGSch_{\mathrm{aff}}}(\clred Y),
\]
其中 $\clred Y:=Y|_{{}^{\mathrm{red}}\Sch_{\mathrm{aff}}}$ 是 $Y$ 到既约经典仿射概形的限制，而
\[
\mathrm{RKE}_{{}^{\mathrm{red}}\Sch_{\mathrm{aff}}\hookrightarrow\DGSch_{\mathrm{aff}}}
\]
是函子沿包含 ${}^{\mathrm{red}}\Sch_{\mathrm{aff}}\hookrightarrow\DGSch_{\mathrm{aff}}$ 的右 Kan 延拓。

\ob{1.1.3}{n:1.1.3}下面这个（显然的）观察在后续会用到。

\begin{othm}{引理}{1.1.4}{n:1.1.4}
函子 $\dR:\PreStk\to\PreStk$ 与极限、余极限交换。
\end{othm}

\begin{proof}
由 $\PreStk=\Funct\big((\DGSch_{\mathrm{aff}})^{\mathrm{op}},\infty\text{-}\Grpd\big)$ 中的极限与余极限逐对象计算可知。
\end{proof}

作为推论，得到：

\begin{othm}{推论}{1.1.5}{n:1.1.5}
函子 $\dR:\PreStk\to\PreStk$ 是函子
\[
\dR|_{\DGSch_{\mathrm{aff}}}:\DGSch_{\mathrm{aff}}\to\PreStk
\]
沿 $\DGSch_{\mathrm{aff}}\hookrightarrow\PreStk$ 的左 Kan 延拓。
\end{othm}

\begin{proof}
对从 $\PreStk$ 到任意 $\infty$-范畴的保余极限函子，这都是对的。
\end{proof}

\ob{1.1.6}{n:1.1.6}此外还有：

\begin{othm}{引理}{1.1.7}{n:1.1.7}
函子 $\dR|_{\DGSch_{\mathrm{aff}}}:\DGSch_{\mathrm{aff}}\to\PreStk$ 同构于函子
\[
\dR|_{{}^{\mathrm{red}}\Sch_{\mathrm{aff}}}:{}^{\mathrm{red}}\Sch_{\mathrm{aff}}\to\PreStk
\]
沿 ${}^{\mathrm{red}}\Sch_{\mathrm{aff}}\hookrightarrow\DGSch_{\mathrm{aff}}$ 的左 Kan 延拓。
\end{othm}

\begin{proof}
对任意目标范畴 $D$ 与任意函子 $\Phi:\DGSch_{\mathrm{aff}}\to D$，映射
\[
\mathrm{LKE}_{{}^{\mathrm{red}}\Sch_{\mathrm{aff}}\hookrightarrow\DGSch_{\mathrm{aff}}}(\Phi|_{{}^{\mathrm{red}}\Sch_{\mathrm{aff}}})\to\Phi
\]
是同构，当且仅当自然变换
\[
\Phi(\clred S)\to\Phi(S),\qquad S\in\DGSch_{\mathrm{aff}}
\]
是同构。取 $D=\PreStk$、$\Phi$ 为函子 $S\mapsto S_{\dR}$ 时，后者按定义成立。
\end{proof}

\ob{1.1.8}{n:1.1.8}设 $\mathcal{C}_1\subset\mathcal{C}_2$ 是 $\PreStk$ 的下列全子范畴中的一对：
\[
{}^{\mathrm{red}}\Sch_{\mathrm{aff}},\ \Sch_{\mathrm{aff}},\ \DGSch_{\mathrm{aff}},\ \DGSch_{\mathrm{qs\text{-}qc}},\ \DGSch,\ \PreStk
\]
（下标「qs-qc」意为「拟分离且拟紧」）。

由引理 1.1.7 与推论 1.1.5 得到：

\begin{othm}{推论}{1.1.9}{n:1.1.9}
由 $\dR|_{\mathcal{C}_2}$ 给出的函子 $\mathcal{C}_2\to\PreStk$，同构于函子 $\dR|_{\mathcal{C}_1}:\mathcal{C}_1\to\PreStk$ 沿 $\mathcal{C}_1\hookrightarrow\mathcal{C}_2$ 的左 Kan 延拓。
\end{othm}

\section{$Y$ 与 $Y_{\dR}$ 的关系}

\ob{1.2.1}{n:1.2.1}函子 $\dR:\PreStk\to\PreStk$ 附带自然变换
\[
p_{\dR}:\mathrm{Id}\to\dR,
\]
即对每个 $Y\in\PreStk$ 有典范映射
\[
p_{\dR,Y}:Y\to Y_{\dR}.
\]

\ob{1.2.2}{n:1.2.2}记 $Y_\bullet/Y_{\dR}$ 为 $p_{\dR,Y}$ 的 Čech 神经，看作 $\PreStk$ 的单纯对象。它以 $Y_{\dR}$ 为增广。

注意每个 $Y_i/Y_{\dR}$ 是 $Y_i$ 沿主对角线的形式完备化。（关于形式完备化的约定，参见 [GR1, Sect.~6.1.1]。）

有典范映射
\begin{equation}
|Y_\bullet/Y_{\dR}|\to Y_{\dR}.\tag{1.2}
\end{equation}

\begin{annote}
直观地说，$Y_i/Y_{\dR}$ 是 $Y^i$ 中「$i$ 个点彼此无穷小接近」所成的形式邻域，而 (1.2) 说的是把这些形式邻域粘起来可以还原出 $Y_{\dR}$。当 $Y$ 经典形式光滑时，这一粘合确实是同构（引理 1.2.4）；一般情形则只是映射。\hfill【来源：[GR1, Sect.~6.1.1]（A）】
\end{annote}

\ob{1.2.3}{n:1.2.3}\textbf{经典形式光滑的预叠。}若对 $S\in\DGSch_{\mathrm{aff}}$，映射
\[
\Maps(S,Y)\to\Maps(\clred S,Y)
\]
在 $\pi_0$ 上诱导满射，则称预叠 $Y$ 经典形式光滑。

下面的结果直接来自定义：

\begin{othm}{引理}{1.2.4}{n:1.2.4}
若 $Y$ 经典形式光滑，则映射
\[
|Y_\bullet/Y_{\dR}|\to Y_{\dR}
\]
是 $\PreStk$ 中的同构。
\end{othm}

\section{局部几乎有限型的情形}

\ob{1.3.1}{n:1.3.1}回忆 $\PreStk$ 含全子范畴 $\PreStk_{\mathrm{laft}}$，由局部几乎有限型的预叠组成，见 [GL:Stacks, Sect.~1.3.9]。按定义，对象 $Y\in\PreStk$ 属于 $\PreStk_{\mathrm{laft}}$，如果：

\begin{itemize}[leftmargin=2.2em]
\item $Y$ 收敛，即对 $S\in\DGSch_{\mathrm{aff}}$，自然映射
\[
\Maps(S,Y)\to\varprojlim_{n\ge 0}\Maps({}^{\le n}S,Y)
\]
是同构，其中 ${}^{\le n}S$ 表示 $S$ 的 $n$-余连接截断；

\item 对每个 $n$，限制 ${}^{\le n}Y:=Y|_{{}^{\le n}\DGSch_{\mathrm{aff}}}$ 属于 ${}^{\le n}\PreStk_{\mathrm{lft}}$，即函子
\[
S\mapsto\Maps(S,Y),\qquad ({}^{\le n}\DGSch_{\mathrm{aff}})^{\mathrm{op}}\to\infty\text{-}\Grpd
\]
与滤过余极限交换（等价地，它是从全子范畴 ${}^{\le n}\DGSch_{\mathrm{aff}}^{\mathrm{ft}}\hookrightarrow{}^{\le n}\DGSch_{\mathrm{aff}}$ 出发的左 Kan 延拓）。
\end{itemize}

\begin{annote}
「收敛」一词在这里是技术性的：它说的是一个预叠的取值由它在各阶截断上的取值决定，正如形式概形由它的无穷小邻域决定。凡本文中出现的 $\PreStk_{\mathrm{laft}}$，都同时要求这两个条件。\hfill【来源：[GL:Stacks, Sect.~1.3.9]（A）】
\end{annote}

\ob{1.3.2}{n:1.3.2}下面的观察在本文中起重要作用。

\begin{othm}{命题}{1.3.3}{n:1.3.3}
设 $Y\in\PreStk_{\mathrm{laft}}$。则：
\begin{enumerate}[label=(\alph*),leftmargin=2.2em,nosep]
\item $Y_{\dR}\in\PreStk_{\mathrm{laft}}$；
\item $Y_{\dR}$ 是经典的，即属于全子范畴 ${}^{\mathrm{cl}}\PreStk\subset\PreStk$。
\end{enumerate}
\end{othm}

\ob{1.3.4}{n:1.3.4}\textbf{(a) 的证明。}

需要验证两条性质：
\begin{enumerate}[label=(\roman*),leftmargin=2.2em,nosep]
\item $Y_{\dR}$ 收敛；
\item 每个截断 ${}^{\le n}(Y_{\dR})$ 局部有限型。
\end{enumerate}
性质 (i) 是同义反复，对任意 $Y\in\PreStk$ 都成立。为证 (ii)，须说明函子 $Y_{\dR}$ 把 ${}^{\le n}\DGSch_{\mathrm{aff}}$ 中的滤过极限送到 $\infty\text{-}\Grpd$ 中的余极限。既然 $Y$ 本身具有这一性质，只需说明函子
\[
S\mapsto\clred S:\DGSch_{\mathrm{aff}}\to\DGSch_{\mathrm{aff}}
\]
保持滤过极限，而这是显然的。

\ob{1.3.5}{n:1.3.5}\textbf{(b) 的证明。}

由推论 1.1.9，须证明余极限
\[
\varinjlim_{S\in(\Sch_{\mathrm{aff}})_{/Y}}S_{\dR}\in\PreStk
\]
是经典的。由 (a)，函子
\[
(\Sch_{\mathrm{aff}}^{\mathrm{ft}})_{/Y}\to(\Sch_{\mathrm{aff}})_{/Y}
\]
共尾；因此
\[
\varinjlim_{S\in(\Sch_{\mathrm{aff}}^{\mathrm{ft}})_{/Y}}S_{\dR}\to\varinjlim_{S\in(\Sch_{\mathrm{aff}})_{/Y}}S_{\dR}
\]
是同构。

于是，由于全子范畴 ${}^{\mathrm{cl}}\PreStk\subset\PreStk$ 对余极限封闭，不妨设 $Y$ 是有限型的经典仿射概形。

更一般地，我们将证明：对 $X\in\DGSch_{\mathrm{aff}}^{\mathrm{aft}}$，预叠 $X_{\dR}$ 是经典的。取闭浸入 $i:X\hookrightarrow Z$，其中 $Z$ 是有限型的光滑经典仿射概形。

记 $Y$ 为 $Z$ 沿 $X$ 的形式完备化 $\widehat{Z}_X$（见 [GR1, Sects.~6.1.1 或 6.5]）。映射 $X\to Y$ 诱导同构 $X_{\dR}\to Y_{\dR}$。因此只需说明 $Y_{\dR}$ 是经典的。

考虑 $Y_\bullet/Y_{\dR}$（见上文 §1.2.2）。注意 $Y$ 形式光滑，因为 $Z$ 形式光滑（见 [GR1, Sect.~8.1]）；特别地，$Y$ 经典形式光滑。由于子范畴 ${}^{\mathrm{cl}}\PreStk\subset\PreStk$ 对余极限封闭，再由引理 1.2.4，只需说明每一项 $Y_i/Y_{\dR}$ 作为预叠是经典的。

注意 $Y_i/Y_{\dR}$ 同构于 $Z^i$ 沿对角嵌入的 $X$ 的形式完备化。于是由 [GR1, Proposition 6.8.2]，$Y_i/Y_{\dR}$ 是经典的。

\begin{annote}
这里的关键一步是把 $X$ 换成它在光滑 $Z$ 中的形式完备化 $Y=\widehat{Z}_X$：$X_{\dR}$ 与 $Y_{\dR}$ 同构，而 $Y$ 形式光滑，于是每一阶无穷小邻域 $Y_i/Y_{\dR}$ 都落在经典概形的范围内。这是 §2.5 中 Kashiwara 引理证明的同一个技巧的原型。\hfill【来源：[GR1, Sect.~6.8]（A）】
\end{annote}

\ob{1.3.6}{n:1.3.6}由命题 1.3.3 得到：

\begin{othm}{推论}{1.3.7}{n:1.3.7}
设 $\mathcal{C}_1\subset\mathcal{C}_2$ 是 $\DGSch_{\mathrm{aff}}$ 的下列全子范畴中的任意一对：
\[
\Sch_{\mathrm{aff}}^{\mathrm{ft}},\quad {}^{<\infty}\DGSch_{\mathrm{aff}}^{\mathrm{aft}},\quad \DGSch_{\mathrm{aff}}^{\mathrm{aft}},\quad \Sch_{\mathrm{aff}},\quad \DGSch_{\mathrm{aff}}.
\]
则对 $Y\in\PreStk_{\mathrm{laft}}$，函子
\[
(\mathcal{C}_1)_{/Y_{\dR}}\to(\mathcal{C}_2)_{/Y_{\dR}}
\]
共尾。
\end{othm}

\begin{proof}
只需对包含
\[
\Sch_{\mathrm{aff}}^{\mathrm{ft}}\hookrightarrow\Sch_{\mathrm{aff}}\hookrightarrow\DGSch_{\mathrm{aff}}
\]
证明论断。对右箭头，断言由命题 1.3.3(b) 得到；对左箭头，由命题 1.3.3(a) 得到。
\end{proof}

\ob{1.3.8}{n:1.3.8}现在考虑下列全子范畴
\begin{equation}
{}^{\mathrm{red}}\Sch_{\mathrm{aff}}^{\mathrm{ft}},\quad \Sch_{\mathrm{aff}}^{\mathrm{ft}},\quad \DGSch_{\mathrm{aff}}^{\mathrm{ft}},\quad \DGSch_{\mathrm{aff}}^{\mathrm{aft}},\quad \DGSch_{\mathrm{laft}},\quad \PreStk_{\mathrm{laft}},\tag{1.3}
\end{equation}
它们是 §1.1.8 中出现的那些范畴的子范畴。

\begin{othm}{推论}{1.3.9}{n:1.3.9}
函子 $\dR$ 到 $\PreStk_{\mathrm{laft}}$ 上的限制，同构于它到 $\mathcal{C}$ 上的左 Kan 延拓，其中 $\mathcal{C}$ 是 (1.3) 中的任一子范畴。
\end{othm}

\begin{proof}
只需对 $\mathcal{C}={}^{\mathrm{red}}\Sch_{\mathrm{aff}}^{\mathrm{ft}}$ 证明。由推论 1.1.9，只需说明对 $Y\in\PreStk_{\mathrm{laft}}$，函子
\[
({}^{\mathrm{red}}\Sch_{\mathrm{aff}}^{\mathrm{ft}})_{/Y}\to(\Sch_{\mathrm{aff}})_{/Y}
\]
共尾。

由命题 1.3.3(a)，函子
\[
(\Sch_{\mathrm{aff}}^{\mathrm{ft}})_{/Y}\to(\Sch_{\mathrm{aff}})_{/Y}
\]
共尾。再由包含 ${}^{\mathrm{red}}\Sch_{\mathrm{aff}}^{\mathrm{ft}}\hookrightarrow\Sch_{\mathrm{aff}}^{\mathrm{ft}}$ 有右伴随这一事实，即得断言。
\end{proof}

\chapter{晶体的定义}
\markboth{第 2 章\ 晶体的定义}{第 2 章\ 晶体的定义}

本节定义左晶体（对 $\PreStk$ 的任意对象），以及对 $\PreStk_{\mathrm{laft}}$ 的对象定义右晶体。我们将证明：在后一情形，两套理论等价。

\section{左晶体}

\ob{2.1.1}{n:2.1.1}对 $Y\in\PreStk$，定义
\[
\Crys^l(Y):=\QCoh(Y_{\dR}).
\]
也就是说，
\[
\Crys^l(Y)=\varprojlim_{S\in(\DGSch_{\mathrm{aff}}/Y_{\dR})^{\mathrm{op}}}\QCoh(S).
\]
粗略地说，对象 $M\in\Crys^l(Y)$ 是如下资料：对每个带映射 $\clred S\to Y$ 的仿射 DG 概形 $S\in\DGSch_{\mathrm{aff}}$，指定一个拟凝聚层 $F_S\in\QCoh(S)$；并且对每个仿射 DG 概形的态射 $f:S'\to S$，指定同构
\[
f^*(F_S)\simeq F_{S'}\in\QCoh(S').
\]

\ob{2.1.2}{n:2.1.2}更函子化地说，记 $\Crys^l_{\PreStk}$ 为函子 $(\PreStk)^{\mathrm{op}}\to\DGCat_{\mathrm{cont}}$，定义为
\[
\Crys^l_{\PreStk}:=\QCoh^*_{\PreStk}\circ\dR,
\]
其中
\[
\QCoh^*_{\PreStk}:(\PreStk)^{\mathrm{op}}\to\DGCat_{\mathrm{cont}}
\]
是把预栈对应到其拟凝聚层范畴的函子 [GL:QCoh, Sect.~1.1.5]。

对 $\PreStk$ 中的映射 $f:Y_1\to Y_2$，记 $f^{\dagger,l}$ 为相应的拉回函子
\[
\Crys^l(Y_2)\to\Crys^l(Y_1).
\]
按构造，若 $f$ 诱导底层既约经典预栈的同构 $\clred Y_1\to\clred Y_2$，则它诱导 de Rham 预栈的同构 $Y_{1,\dR}\to Y_{2,\dR}$；特别地，$f^{\dagger,l}$ 是等价。

\ob{2.1.3}{n:2.1.3}回忆函子 $\QCoh^*_{\PreStk}:(\PreStk)^{\mathrm{op}}\to\DGCat_{\mathrm{cont}}$ 按定义是函子
\[
\QCoh^*_{\DGSch_{\mathrm{aff}}}:(\DGSch_{\mathrm{aff}})^{\mathrm{op}}\to\DGCat_{\mathrm{cont}}
\]
沿 $(\DGSch_{\mathrm{aff}})^{\mathrm{op}}\hookrightarrow(\PreStk)^{\mathrm{op}}$ 的右 Kan 延拓。

特别地，它把 $\PreStk$ 中的余极限送到 $\DGCat_{\mathrm{cont}}$ 中的极限。因此由推论 1.1.9，对 $Y\in\PreStk$ 得到：

\begin{othm}{推论}{2.1.4}{n:2.1.4}
设 $\mathcal{C}$ 是 §1.1.8 列表中的任一范畴。则对 $Y\in\PreStk$，函子
\[
\Crys^l(Y)\to\varprojlim_{X\in(\mathcal{C}_{/Y})^{\mathrm{op}}}\Crys^l(X)
\]
是等价。
\end{othm}

粗略地说，这个推论表明：对象 $M\in\Crys^l(Y)$ 的数据，等价于对每个 $S\in\mathcal{C}_{/Y}$ 给定 $M_S\in\Crys^l(S)$，并对每个 $f:S'\to S$ 给定同构
\[
f^{\dagger,l}(M_S)\simeq M_{S'}\in\Crys^l(S').
\]

\ob{2.1.5}{n:2.1.5}回忆自然变换 $p_{\dR}:\mathrm{Id}\to\dR$。它诱导自然变换
\[
\oblv^l:\Crys^l_{\PreStk}\to\QCoh^*_{\PreStk}.
\]
也就是说，对每个 $Y\in\PreStk$ 有函子
\begin{equation}
\oblv^l_Y:\Crys^l(Y)\to\QCoh(Y),\tag{2.1}
\end{equation}
且对每个态射 $f:Y_1\to Y_2$ 有交换图
\begin{equation}
\begin{tikzcd}[column sep=large, row sep=large]
 \Crys^l(Y_1) \arrow[r, "\oblv^l_{Y_1}"] \arrow[d, "f^{\dagger,l}"'] & \QCoh(Y_1) \arrow[d, "f^*"] \\
 \Crys^l(Y_2) \arrow[r, "\oblv^l_{Y_2}"'] & \QCoh(Y_2).
\end{tikzcd}\tag{2.2}
\end{equation}

\ob{2.1.6}{n:2.1.6}回忆 §1.2.2 的单纯对象 $Y_\bullet/Y_{\dR}$。

由引理 1.2.4 得到：

\begin{othm}{引理}{2.1.7}{n:2.1.7}
若 $Y$ 经典形式光滑，则函子
\[
\Crys^l(Y)\to\Tot\big(\QCoh(Y_\bullet/Y_{\dR})\big)
\]
是等价。
\end{othm}

\begin{othm}{注}{2.1.8}{n:2.1.8}
本书对 $Y$ 上左晶体的定义，就是 Grothendieck 术语中 $Y$ 的无穷小 site 上的拟凝聚层。范畴 $\Tot(\QCoh(Y_\bullet/Y_{\dR}))$ 则是 Grothendieck 术语中 $Y$ 的分层 site 上的拟凝聚层。于是引理 2.1.7 说：对经典形式光滑的预栈，两者等价。我们将在 §3.4 看到，当 $Y$ 是局部几乎有限型的最终余连接 DG 概形时，同样的结论也成立。不过，对非最终余连接的 DG 概形（即使是局部几乎有限型的），这个等价不成立。
\end{othm}

\section{局部几乎有限型预栈上的左晶体}

在本节余下部分，除另有说明，我们把讨论限制在子范畴 $\PreStk_{\mathrm{laft}}\subset\PreStk$。

因此，除明确说明外，所谓预栈/DG 概形/仿射 DG 概形，都指局部几乎有限型的。

记 $\Crys^l_{\PreStk_{\mathrm{laft}}}$ 为 $\Crys^l_{\PreStk}$ 到 $\PreStk_{\mathrm{laft}}\subset\PreStk$ 的限制。

\ob{2.2.1}{n:2.2.1}下面的推论说明：为了对 $Y\in\PreStk_{\mathrm{laft}}$ 定义 $\Crys^l(Y)$，我们「不需要知道」无限型概形或导出代数几何。换言之，要在局部几乎有限型的预栈上定义晶体，可以留在有限型经典仿射概形的世界里。

事实上，由推论 1.3.7 得到：

\begin{othm}{推论}{2.2.2}{n:2.2.2}
设 $\mathcal{C}$ 是 $\DGSch_{\mathrm{aff}}$ 的下列全子范畴之一：
\[
\Sch_{\mathrm{aff}}^{\mathrm{ft}},\quad {}^{<\infty}\DGSch_{\mathrm{aff}}^{\mathrm{aft}},\quad \DGSch_{\mathrm{aff}}^{\mathrm{aft}},\quad \Sch_{\mathrm{aff}}.
\]
则对 $Y\in\PreStk_{\mathrm{laft}}$，自然函子
\[
\Crys^l(Y)\to\varprojlim_{S\in(\mathcal{C}_{/Y_{\dR}})^{\mathrm{op}}}\QCoh(S)
\]
是等价。
\end{othm}

\ob{2.2.3}{n:2.2.3}回忆：按推论 2.1.4，范畴 $\Crys^l(Y)$ 可以由函子
\[
\Crys^l:\mathcal{C}_{/Y}\to\DGCat_{\mathrm{cont}}
\]
还原，其中 $\mathcal{C}$ 是下列范畴之一：
\[
{}^{\mathrm{red}}\Sch_{\mathrm{aff}},\ \Sch_{\mathrm{aff}},\ \DGSch_{\mathrm{aff}},\ \DGSch_{\mathrm{qs\text{-}qc}},\ \DGSch\subset\PreStk.
\]
现在断言：上述范畴也可以换成它们在 (1.3) 中的全子范畴：
\[
{}^{\mathrm{red}}\Sch_{\mathrm{aff}}^{\mathrm{ft}},\quad \Sch_{\mathrm{aff}}^{\mathrm{ft}},\quad \DGSch_{\mathrm{aff}}^{\mathrm{ft}},\quad \DGSch_{\mathrm{aff}}^{\mathrm{aft}},\quad \DGSch_{\mathrm{laft}},\quad \PreStk_{\mathrm{laft}}.
\]

\begin{othm}{推论}{2.2.4}{n:2.2.4}
设 $Y\in\PreStk_{\mathrm{laft}}$，$\mathcal{C}$ 是 (1.3) 中的范畴之一。则函子
\[
\Crys^l(Y)\to\varprojlim_{X\in(\mathcal{C}_{/Y})^{\mathrm{op}}}\Crys^l(X)
\]
是等价。
\end{othm}

\begin{proof}
由推论 1.3.9 得到。
\end{proof}

粗略地说，上述推论表明：对象 $M\in\Crys^l(Y)$ 可以由如下资料还原：对每个 $S\in\mathcal{C}_{/Y}$ 给定 $M_S\in\Crys^l(S)$，并对每个 $f:S'\to S$ 给定同构
\[
f^{\dagger,l}(M_S)\simeq M_{S'}\in\Crys^l(S').
\]

\ob{2.2.5}{n:2.2.5}再次考虑 (2.1) 中的函子
\[
\oblv^l_Y:\Crys^l(Y)\to\QCoh(Y).
\]
有：

\begin{othm}{引理}{2.2.6}{n:2.2.6}
对 $Y\in\PreStk_{\mathrm{laft}}$，函子 $\oblv^l_Y$ 保守。
\end{othm}

其证明推迟到 §2.4.7。

\section{右晶体}

\ob{2.3.1}{n:2.3.1}回忆 $\PreStk_{\mathrm{laft}}$ 也可以看作全体函子
\[
({}^{<\infty}\DGSch_{\mathrm{aff}}^{\mathrm{aft}})^{\mathrm{op}}\to\infty\text{-}\Grpd
\]
构成的范畴，见 [GL:Stacks, Sect.~1.3.11]。

此外，有 [IndCoh, Sect.~10.1.2] 中的函子
\[
\IndCoh^!_{\PreStk_{\mathrm{laft}}}:(\PreStk_{\mathrm{laft}})^{\mathrm{op}}\to\DGCat_{\mathrm{cont}},
\]
它定义为相应函子
\[
\IndCoh^!_{{}^{<\infty}\DGSch_{\mathrm{aff}}^{\mathrm{aft}}}:({}^{<\infty}\DGSch_{\mathrm{aff}}^{\mathrm{aft}})^{\mathrm{op}}\to\DGCat_{\mathrm{cont}}
\]
沿
\[
({}^{<\infty}\DGSch_{\mathrm{aff}}^{\mathrm{aft}})^{\mathrm{op}}\to(\PreStk_{\mathrm{laft}})^{\mathrm{op}}
\]
的右 Kan 延拓。

特别地，函子 $\IndCoh^!_{\PreStk_{\mathrm{laft}}}$ 把 $\PreStk_{\mathrm{laft}}$ 中的余极限送到 $\DGCat_{\mathrm{cont}}$ 中的极限。

\ob{2.3.2}{n:2.3.2}定义函子
\[
\Crys^r_{\PreStk_{\mathrm{laft}}}:(\PreStk_{\mathrm{laft}})^{\mathrm{op}}\to\DGCat_{\mathrm{cont}}
\]
为复合
\[
\Crys^r_{\PreStk_{\mathrm{laft}}}:=\IndCoh^!_{\PreStk_{\mathrm{laft}}}\circ\dR.
\]
上式中用到命题 1.3.3(a)，以确保 $\dR$ 是函子 $\PreStk_{\mathrm{laft}}\to\PreStk_{\mathrm{laft}}$。

\begin{othm}{注}{2.3.3}{n:2.3.3}
在定义 $\Crys^r_{\PreStk_{\mathrm{laft}}}$ 时，我们「不需要知道」无限型概形：自函子 $\dR:\PreStk_{\mathrm{laft}}\to\PreStk_{\mathrm{laft}}$ 可以直接由公式
\[
\Maps(S,Y_{\dR})=\Maps(\clred S,Y)
\]
定义，其中 $S\in{}^{<\infty}\DGSch_{\mathrm{aff}}^{\mathrm{aft}}$。
\end{othm}

\ob{2.3.4}{n:2.3.4}对 $\PreStk_{\mathrm{laft}}$ 中的映射 $f:Y_1\to Y_2$，记 $f^{\dagger,r}$ 为相应的函子 $\Crys^r(Y_2)\to\Crys^r(Y_1)$。

若 $f$ 诱导等价 $\clred Y_1\to\clred Y_2$，则映射 $Y_{1,\dR}\to Y_{2,\dR}$ 是等价；特别地，$f^{\dagger,r}$ 也是等价。

\ob{2.3.5}{n:2.3.5}按定义，对 $Y\in\PreStk_{\mathrm{laft}}$ 有
\[
\Crys^r(Y)=\varprojlim_{S\in(({}^{<\infty}\DGSch_{\mathrm{aff}}^{\mathrm{aft}})_{/Y_{\dR}})^{\mathrm{op}}}\IndCoh(S).
\]
粗略地说，对象 $M\in\Crys^r(Y)$ 是如下资料：对每个 $S\in{}^{<\infty}\DGSch_{\mathrm{aff}}^{\mathrm{aft}}$ 及映射 $\clred S\to Y$，指定一个对象 $F_S\in\IndCoh(S)$；并且对每个 $f:S'\to S$，指定同构
\[
f^!(F_S)\simeq F_{S'}\in\IndCoh(S').
\]

\ob{2.3.6}{n:2.3.6}如同 §2.2.1，为了还原 $\Crys^r(Y)$，我们「不需要知道」DG 概形：

\begin{othm}{推论}{2.3.7}{n:2.3.7}
对 $Y\in\PreStk_{\mathrm{laft}}$，函子
\[
\Crys^r(Y)\to\varprojlim_{S\in((\Sch_{\mathrm{aff}}^{\mathrm{ft}})_{/Y_{\dR}})^{\mathrm{op}}}\IndCoh(S)
\]
是等价。
\end{othm}

\begin{proof}
容易由推论 1.3.7 推出。
\end{proof}

粗略地说，上述推论表明：$M\in\Crys^r(Y)$ 可以由如下资料还原：对每个 $S\in\Sch_{\mathrm{aff}}^{\mathrm{ft}}$ 及映射 $\clred S\to Y$，指定一个对象 $F_S\in\IndCoh(S)$；并且对每个 $f:S'\to S$，指定同构
\[
f^!(F_S)\simeq F_{S'}\in\IndCoh(S').
\]

\ob{2.3.8}{n:2.3.8}此外，推论 2.2.4 的类比对右晶体也成立：

\begin{othm}{推论}{2.3.9}{n:2.3.9}
设 $\mathcal{C}$ 是 (1.3) 中的任一范畴。则函子
\[
\Crys^r(Y)\to\varprojlim_{X\in(\mathcal{C}_{/Y})^{\mathrm{op}}}\Crys^r(X)
\]
是等价。
\end{othm}

\begin{proof}
由推论 1.3.9 得到。
\end{proof}

粗略地说，这个推论表明：对象 $M\in\Crys^r(Y)$ 可以由如下资料还原：对每个 $S\in\mathcal{C}_{/Y}$ 给定 $M_S\in\Crys^r(S)$，并对每个 $f:S'\to S$ 给定同构
\[
f^{\dagger,r}(M_S)\simeq M_{S'}\in\Crys^r(S').
\]

\ob{2.3.10}{n:2.3.10}自然变换 $p_{\dR}:\mathrm{Id}\to\dR$ 诱导自然变换
\[
\oblv^r:\Crys^r_{\PreStk_{\mathrm{laft}}}\to\IndCoh_{\PreStk_{\mathrm{laft}}}.
\]
也就是说，对每个 $Y\in\PreStk_{\mathrm{laft}}$ 有函子
\[
\oblv^r_Y:\Crys^r(Y)\to\IndCoh(Y),
\]
且对每个态射 $f:Y_1\to Y_2$ 有交换图
\begin{equation}
\begin{tikzcd}[column sep=large, row sep=large]
 \Crys^r(Y_1) \arrow[r, "\oblv^r_{Y_1}"] \arrow[d, "f^{\dagger,r}"'] & \IndCoh(Y_1) \arrow[d, "f^!"] \\
 \Crys^r(Y_2) \arrow[r, "\oblv^r_{Y_2}"'] & \IndCoh(Y_2).
\end{tikzcd}\tag{2.3}
\end{equation}
有：

\begin{othm}{引理}{2.3.11}{n:2.3.11}
若 $Y$ 经典形式光滑，则函子
\[
\Crys^r(Y)\to\Tot\big(\IndCoh(Y_\bullet/Y_{\dR})\big)
\]
是等价。
\end{othm}

\begin{proof}
与引理 2.1.7 的证明相同，即由引理 1.2.4 推出。
\end{proof}

\begin{othm}{引理}{2.3.12}{n:2.3.12}
对任意 $Y$，函子 $\oblv^r_Y$ 保守。
\end{othm}

\begin{proof}
由推论 2.3.9 与 (2.3) 的交换性，不妨设 $Y=X$ 是局部几乎有限型的仿射 DG 概形。取闭浸入 $i:X\to Z$，其中 $Z$ 是光滑经典有限型概形，并记 $Y$ 为 $Z$ 沿 $X$ 的形式完备化，$i':X\to Y$ 为相应的映射。

考虑交换图
\[
\begin{tikzcd}[column sep=large, row sep=large]
 \Crys^r(Y) \arrow[r, "\oblv^r_Y"] \arrow[d, "i'^{\dagger,r}"'] & \IndCoh(Y) \arrow[d, "i'^!"] \\
 \Crys^r(X) \arrow[r, "\oblv^r_X"'] & \IndCoh(X).
\end{tikzcd}
\]
此图中左竖直箭头是等价，因为 $i'_{\dR}:X_{\dR}\to Y_{\dR}$ 是同构。上方水平箭头保守，因为 $Y$ 形式光滑（特别地，经典形式光滑），可用引理 2.3.11。

于是只需说明函子 $i'^!$ 保守。这可以由 [GR1, Proposition 7.4.5] 与 [IndCoh, Proposition 4.1.7(a)] 合并得到。
\end{proof}

\section{左晶体与右晶体的比较}

请读者回想：我们假定所有预栈与 DG 概形都局部几乎有限型。

\ob{2.4.1}{n:2.4.1}回忆（见 [IndCoh, Sect.~5.7.5]）：对任意 $S\in\DGSch_{\mathrm{aff}}^{\mathrm{aft}}$，有典范定义的函子
\[
\Upsilon_S:\QCoh(S)\to\IndCoh(S),
\]
它由与对偶化层 $\omega_S\in\IndCoh(S)$ 作张量给出，使得对 $f:S_1\to S_2$，图表
\[
\begin{tikzcd}[column sep=large, row sep=large]
 \QCoh(S_1) \arrow[r, "\Upsilon_{S_1}"] \arrow[d, "f^*"'] & \IndCoh(S_1) \arrow[d, "f^!"] \\
 \QCoh(S_2) \arrow[r, "\Upsilon_{S_2}"'] & \IndCoh(S_2)
\end{tikzcd}
\]
典范交换。事实上，上述数据提升为函子的自然变换
\[
\Upsilon_{\DGSch_{\mathrm{aff}}^{\mathrm{aft}}}:\QCoh^*_{\DGSch_{\mathrm{aff}}^{\mathrm{aft}}}\to\IndCoh^!_{\DGSch_{\mathrm{aff}}^{\mathrm{aft}}},
\]
从而给出自然变换
\[
\Upsilon_{\PreStk_{\mathrm{laft}}}:\QCoh^*_{\PreStk_{\mathrm{laft}}}\to\IndCoh^!_{\PreStk_{\mathrm{laft}}},
\]
见 [IndCoh, Sect.~10.3.3]。

对单个对象 $Y\in\PreStk_{\mathrm{laft}}$，我们得到函子
\[
\Upsilon_Y:\QCoh(Y)\to\IndCoh(Y).
\]

\ob{2.4.2}{n:2.4.2}对 $Y\in\PreStk_{\mathrm{laft}}$，把 $\Upsilon$ 作用于 $Y_{\dR}$，得到典范定义的函子
\begin{equation}
\Upsilon_{Y_{\dR}}:\Crys^l(Y)\to\Crys^r(Y),\tag{2.4}
\end{equation}
使得图表
\begin{equation}
\begin{tikzcd}[column sep=large, row sep=large]
 \Crys^l(Y) \arrow[r, "\Upsilon_{Y_{\dR}}"] \arrow[d, "\oblv^l_Y"'] & \Crys^r(Y) \arrow[d, "\oblv^r_Y"] \\
 \QCoh(Y) \arrow[r, "\Upsilon_Y"'] & \IndCoh(Y)
\end{tikzcd}\tag{2.5}
\end{equation}
交换。

事实上，我们得到自然变换
\[
\Upsilon_{\PreStk_{\mathrm{laft}}}\circ\dR:\Crys^l_{\PreStk_{\mathrm{laft}}}\to\Crys^r_{\PreStk_{\mathrm{laft}}}.
\]
特别地，对 $f:Y_1\to Y_2$，图表
\[
\begin{tikzcd}[column sep=large, row sep=large]
 \Crys^l(Y_1) \arrow[r, "\Upsilon_{Y_{1,\dR}}"] \arrow[d, "f^{\dagger,l}"'] & \Crys^r(Y_1) \arrow[d, "f^{\dagger,r}"] \\
 \Crys^l(Y_2) \arrow[r, "\Upsilon_{Y_{2,\dR}}"'] & \Crys^r(Y_2)
\end{tikzcd}
\]
交换。

\ob{2.4.3}{n:2.4.3}我们断言：

\begin{othm}{命题}{2.4.4}{n:2.4.4}
对 $Y\in\PreStk_{\mathrm{laft}}$，函子 (2.4) 是等价。
\end{othm}

\begin{proof}
由推论 2.2.4 与 2.3.9，命题化归为：对几乎有限型的仿射 DG 概形 $X$，函子
\[
\Upsilon_{X_{\dR}}:\Crys^l(X)\to\Crys^r(X)
\]
是等价。

取闭浸入 $i:X\hookrightarrow Z$，其中 $Z$ 是光滑经典概形，并记 $Y$ 为 $Z$ 沿 $X$ 的形式完备化。由于 $X_{\dR}\to Y_{\dR}$ 是同构，函子
\[
f^{\dagger,l}:\Crys^l(Y)\to\Crys^l(X),\qquad f^{\dagger,r}:\Crys^r(Y)\to\Crys^r(X)
\]
都是等价。因此只需对 $Y$ 证明断言。

记 $Y_\bullet/Y_{\dR}$ 为 $\PreStk_{\mathrm{laft}}$ 中对应于映射
\[
p_{\dR,Y}:Y\to Y_{\dR}
\]
的 Čech 神经。

考虑交换图
\[
\begin{tikzcd}[column sep=large, row sep=large]
 \Crys^l(Y) \arrow[r, "\Upsilon_{Y_{\dR}}"] \arrow[d] & \Crys^r(Y) \arrow[d] \\
 \Tot\big(\QCoh(Y_\bullet/Y_{\dR})\big) \arrow[r, "\Tot(\Upsilon_{Y_\bullet/Y_{\dR}})"'] & \Tot\big(\IndCoh(Y_\bullet/Y_{\dR})\big).
\end{tikzcd}
\]
由引理 2.1.7 与 2.3.11，图中竖直箭头是等价。因此只需证明：对每个 $i$，
\[
\Upsilon_{Y^i/Y_{\dR}}:\QCoh(Y^i/Y_{\dR})\to\IndCoh(Y^i/Y_{\dR})
\]
是等价。

回忆（命题 1.3.3 的证明中也有）：$Y^i/Y_{\dR}$ 是光滑经典概形 $Z^i$ 沿 $X$ 的对角线副本的形式完备化。记 $U_i\subset Z^i$ 为补开子预栈。

由 [GR1, Propositions 7.1.3 and 7.4.5 and Diagram (7.16)]，我们得到 DG 范畴「短正合列」的映射
\[
\begin{tikzcd}[column sep=small, row sep=large]
 0 \arrow[r] & \QCoh(Y^i/Y_{\dR}) \arrow[r] \arrow[d, "\Upsilon_{Y^i/Y_{\dR}}"'] & \QCoh(Z^i) \arrow[r] \arrow[d, "\Upsilon_{Z^i}"] & \QCoh(U_i) \arrow[r] \arrow[d, "\Upsilon_{U_i}"] & 0 \\
 0 \arrow[r] & \IndCoh(Y^i/Y_{\dR}) \arrow[r] & \IndCoh(Z^i) \arrow[r] & \IndCoh(U_i) \arrow[r] & 0.
\end{tikzcd}
\]
现在，函子
\[
\Upsilon_{Z^i}:\QCoh(Z^i)\to\IndCoh(Z^i)\quad\text{与}\quad \Upsilon_{U_i}:\QCoh(U_i)\to\IndCoh(U_i)
\]
都是等价，因为 $Z^i$ 与 $U_i$ 光滑。

事实上，由 [IndCoh, Proposition 9.3.3]，对任意 $S\in\DGSch_{\mathrm{aff}}^{\mathrm{aft}}$，函子 $\Upsilon_S$ 是 $\Psi_S:\IndCoh(S)\to\QCoh(S)$ 的对偶；而后者在 $S$ 光滑时是等价，见 [IndCoh, Lemma 1.1.6]。
\end{proof}

\ob{2.4.5}{n:2.4.5}命题 2.4.4 使我们可以对 $Y\in\PreStk_{\mathrm{laft}}$ 认同左晶体与右晶体。

换言之，我们可以把范畴 $\Crys(Y)$ 看作带两个实现：「左」实现与「右」实现；它们分别体现为从 $\Crys(Y)$ 到 $\QCoh(Y)$ 与 $\IndCoh(Y)$ 的遗忘函子 $\oblv^l_Y$ 与 $\oblv^r_Y$。

两个遗忘函子由交换图
\begin{equation}
\begin{tikzcd}[column sep=large, row sep=large]
 & \Crys(Y) \arrow[dl, "\oblv^l_Y"'] \arrow[dr, "\oblv^r_Y"] & \\
 \QCoh(Y) \arrow[rr, "\Upsilon_Y"'] & & \IndCoh(Y)
\end{tikzcd}\tag{2.6}
\end{equation}
相联系。

对态射 $f:Y_1\to Y_2$，有自然定义的函子
\[
f^{\dagger}:\Crys(Y_2)\to\Crys(Y_1),
\]
它使下列图表交换：
\[
\begin{tikzcd}[column sep=large, row sep=large]
 \Crys(Y_1) & \Crys(Y_2) \arrow[l, "f^{\dagger}"'] \\
 \QCoh(Y_1) \arrow[u, "\oblv^l_{Y_1}"] & \QCoh(Y_2) \arrow[l, "f^*"] \arrow[u, "\oblv^l_{Y_2}"']
\end{tikzcd}
\]
以及
\[
\begin{tikzcd}[column sep=large, row sep=large]
 \Crys(Y_1) & \Crys(Y_2) \arrow[l, "f^{\dagger}"'] \\
 \IndCoh(Y_1) \arrow[u, "\oblv^r_{Y_1}"] & \IndCoh(Y_2) \arrow[l, "f^!"] \arrow[u, "\oblv^r_{Y_2}"'].
\end{tikzcd}
\]

\ob{2.4.6}{n:2.4.6}在下文中，我们将混用记号 $\Crys(Y)$、$\Crys^r(Y)$ 与 $\Crys^l(Y)$：前者强调命题不依赖于所选实现（左或右），后两者则在实现的选取重要时使用。

\ob{2.4.7}{n:2.4.7}引理 2.2.6 的证明。
\begin{proof}
由引理 2.3.12 与命题 2.4.4 合并得到。
\end{proof}

\section{Kashiwara 引理}

映射 $Y\mapsto\Crys(Y)$ 的一个特点是：Kashiwara 引理几乎成为同义反复。

我们将在晶体的右晶体实现中陈述并证明它。由命题 2.4.4，这也蕴含左晶体的相应断言。不过，用左晶体的语言写出同样的证明也很容易。

\ob{2.5.1}{n:2.5.1}回忆：$\PreStk$ 中的映射 $i:X\to Z$ 称为闭浸入，如果它在底层经典预栈的层面是闭浸入。也就是说，对每个 $S\in(\Sch_{\mathrm{aff}})_{/Z}$，基变换映射
\[
{}^{cl}(S\times_Z X)\to S
\]
是闭浸入；特别地，${}^{cl}(S\times_Z X)$ 是经典仿射概形。

若 $X,Z\in\PreStk_{\mathrm{laft}}$，则只需对 $S\in(\Sch_{\mathrm{aff}}^{\mathrm{ft}})_{/Z}$ 验证上述条件。

\ob{2.5.2}{n:2.5.2}对 $\PreStk_{\mathrm{laft}}$ 中对象的闭浸入 $i:X\hookrightarrow Z$，记 $j:Z^{\circ}\hookrightarrow Z$ 为补开浸入。诱导映射
\[
j_{\dR}:Z^{\circ}_{\dR}\to Z_{\dR}
\]
也是预栈的开浸入。考虑限制函子
\[
j^{\dagger,r}:\Crys^r(Z)\to\Crys^r(Z^{\circ}).
\]
由 [IndCoh, Lemma 4.1.1]，上述函子有全忠实的右伴随，记为 $j_{\dR,*}$，使得对每个 $S\in(\DGSch_{\mathrm{aff}}^{\mathrm{ft}})_{/Z_{\dR}}$ 及
\[
S^{\circ}:=S\times_{Z_{\dR}}Z^{\circ}_{\dR}\hookrightarrow S,
\]
由图
\[
\begin{tikzcd}[column sep=large, row sep=large]
 \IndCoh(S) & \Crys^r(Z) \arrow[l, "\oblv^r_Z"'] \\
 \IndCoh(S^{\circ}) \arrow[u, "(j_S)_*^{\IndCoh}"] & \Crys^r(Z^{\circ}) \arrow[l, "\oblv^r_{Z^{\circ}}"'] \arrow[u, "j_{\dR,*}"']
\end{tikzcd}
\]
经伴随性得到的图
\[
\begin{tikzcd}[column sep=large, row sep=large]
 \IndCoh(S) & \Crys^r(Z) \arrow[l, "\oblv^r_Z"'] \\
 \IndCoh(S^{\circ}) \arrow[u, "(j_S)^!"] & \Crys^r(Z^{\circ}) \arrow[l, "\oblv^r_{Z^{\circ}}"'] \arrow[u, "j^{\dagger,r}"']
\end{tikzcd}
\]
中的自然变换是同构。

特别地，自然变换
\[
\oblv^r_Z\circ j_{\dR,*}\to j_*^{\IndCoh}\circ\oblv^r_{Z^{\circ}}
\]
是同构。

\ob{2.5.3}{n:2.5.3}记 $\Crys^r(Z)_X$ 为 $\Crys^r(Z)$ 中等于 $\ker(j^{\dagger,r})$ 的全子范畴。

显然，对象 $M\in\Crys^r(Z)$ 属于 $\Crys^r(Z)_X$，当且仅当：对每个带映射 $\clred S\to Z$ 的 $S\in\DGSch_{\mathrm{aff}}^{\mathrm{ft}}$，相应的对象 $F_S\in\IndCoh(S)$ 落在
\[
\IndCoh(S)_{S-S^{\circ}}:=\ker\big(j_S^!:\IndCoh(S)\to\IndCoh(S^{\circ})\big)
\]
中。

\ob{2.5.4}{n:2.5.4}函子 $\Crys^r(Z)_X\hookrightarrow\Crys^r(Z)$ 有右伴随，由
\[
M\mapsto\Cone\big(M\to j_{\dR,*}\circ j^{\dagger,r}(M)\big)[-1]
\]
给出。因此，可以把 $\Crys^r(Z)_X$ 看作 $\Crys^r(Z)$ 的余局部化。

\ob{2.5.5}{n:2.5.5}由于复合 $i^{\dagger,r}\circ j_{\dR,*}$ 为零，函子 $i^{\dagger,r}:\Crys^r(Z)\to\Crys^r(X)$ 穿过上述余局部化：
\[
\Crys^r(Z)\to\Crys^r(Z)_X\xrightarrow{i'^{\dagger,r}}\Crys^r(X).
\]
Kashiwara 引理断言：

\begin{othm}{命题}{2.5.6}{n:2.5.6}
上述函子
\[
i'^{\dagger,r}:\Crys^r(Z)_X\to\Crys^r(X)
\]
是等价。
\end{othm}

\begin{proof}
注意在 $\PreStk_{\mathrm{laft}}$ 中有同构
\[
\varinjlim_{S\in(\DGSch_{\mathrm{aff}}^{\mathrm{ft}})_{/Z_{\dR}}}S\times_{Z_{\dR}}X_{\dR}\simeq X_{\dR}.
\]
此外，$S^{\wedge}:=S\times_{Z_{\dR}}X_{\dR}$ 等同于 $S$ 沿
\[
\clred S\times_{\clred Z}\clred X
\]
的形式完备化。

于是范畴 $\Crys^r(X)$ 可以描述为
\[
\varprojlim_{S\in((\DGSch_{\mathrm{aff}}^{\mathrm{ft}})_{/Z_{\dR}})^{\mathrm{op}}}\IndCoh(S^{\wedge}).
\]
按定义，范畴 $\Crys^r(Z)_X$ 由
\[
\varprojlim_{S\in((\DGSch_{\mathrm{aff}}^{\mathrm{ft}})_{/Z_{\dR}})^{\mathrm{op}}}
\ker\big(j_S^!:\IndCoh(S)\to\IndCoh(S^{\circ})\big)
\]
给出。

现在，[GR1, Proposition 7.4.5] 说：对上述任意 $S$，$!$-拉回给出等价
\[
\ker\big(j_S^!:\IndCoh(S)\to\IndCoh(S^{\circ})\big)\to\IndCoh(S^{\wedge}),
\]
这正是所要的。
\end{proof}

\begin{othm}{注}{2.5.7}{n:2.5.7}
若把上述证明改用左晶体而非右晶体的语言表述，我们就要用 [GR1, Proposition 7.1.3] 代替 [GR1, Proposition 7.4.5]。
\end{othm}

\chapter{晶体的下降性质}
\markboth{第 3 章\ 晶体的下降性质}{第 3 章\ 晶体的下降性质}

本节中，所有预栈（包括 DG 概形与 DG ind-概形）除明确说明外都假定为局部几乎有限型。

本节的目标是建立一系列关于晶体在 DG 概形与 DG ind-概形上行为的性质。这些性质包括：用无穷小群胚解释（左、右）晶体；h-下降；晶体范畴的单子描述；右晶体与左晶体的归纳函子。

\section{无穷小群胚}

本小节中，$X$ 是局部几乎有限型的 DG ind-概形，见 [GR1, Sect.~1.7.1]。

\ob{3.1.1}{n:3.1.1}考虑单纯预栈 $X_\bullet/X_{\dR}$，即映射 $X\to X_{\dR}$ 的 Čech 神经。前面已经指出，每个 $X^i/X_{\dR}$ 都是 $X^i$ 沿主对角线的形式完备化。特别地，所有 $X^i/X_{\dR}$ 也都属于 $\DGIndSch_{\mathrm{laft}}$。

我们把
\[
X\times_{X_{\dR}}X\rightrightarrows X
\]
称为 $X$ 的无穷小群胚。

\ob{3.1.2}{n:3.1.2}考虑余单纯范畴 $\IndCoh(X_\bullet/X_{\dR})$。

\begin{othm}{命题}{3.1.3}{n:3.1.3}
由增广定义的函子
\[
\Crys^r(X)\to\Tot\big(\IndCoh(X_\bullet/X_{\dR})\big)
\]
是等价。
\end{othm}

\begin{othm}{注}{3.1.4}{n:3.1.4}
注意：由引理 2.3.11，把 $X$ 换成任意经典形式光滑的对象 $Y\in\PreStk_{\mathrm{laft}}$，命题的断言仍成立。
\end{othm}

\begin{proof}
只需证明：对任意 $S\in\DGSch_{\mathrm{aff}}^{\mathrm{ft}}$ 及映射 $S\to X_{\dR}$，函子
\[
\IndCoh(S)\to\Tot\Big(\IndCoh\big(S\times_{X_{\dR}}X^\bullet/X_{\dR}\big)\Big)
\]
是等价。

注意，单纯预栈 $S\times_{X_{\dR}}(X_\bullet/X_{\dR})$ 是映射
\begin{equation}
S\times_{X_{\dR}}X\to S\tag{3.1}
\end{equation}
的 Čech 神经。

注意 $S\times_{X_{\dR}}X$ 等同于 $S\times X$ 沿映射 $\clred S\to S\times X$ 的形式完备化，其中 $\clred S\to X$ 是对应于 $S\to X_{\dR}$ 的映射。特别地，我们得到：(3.1) 中的映射是 ind-固有的（ind-固有性的概念见 [GR1, Sect.~2.7.4]）且满。

于是，断言由 [GR1, Lemma 2.10.3] 推出。
\end{proof}

\section{晶体的 fppf 下降与 h-下降}

\ob{3.2.1}{n:3.2.1}回忆范畴 $\DGSch_{\mathrm{aff}}^{\mathrm{aft}}$ 上的 h-拓扑 [IndCoh, Sect.~8.2]。它由 Zariski 覆盖与固有满覆盖生成。

考虑函子
\[
\Crys^r_{\DGSch_{\mathrm{aff}}^{\mathrm{aft}}}:=\Crys^r_{\PreStk_{\mathrm{laft}}}|_{\DGSch_{\mathrm{aff}}^{\mathrm{aft}}}:(\DGSch_{\mathrm{aff}}^{\mathrm{aft}})^{\mathrm{op}}\to\DGCat.
\]
我们将证明：

\begin{othm}{命题}{3.2.2}{n:3.2.2}
函子 $\Crys^r_{\DGSch_{\mathrm{aff}}^{\mathrm{aft}}}$ 满足 h-下降。
\end{othm}

\begin{proof}
我们将证明 $\Crys^r_{\DGSch_{\mathrm{aff}}^{\mathrm{aft}}}$ 满足平展下降与固有满下降。

平展下降是清楚的：若 $S'\to S$ 是 $\DGSch_{\mathrm{aff}}^{\mathrm{aft}}$ 中的平展覆盖，则相应的映射 $S'_{\dR}\to S_{\dR}$ 是 $\PreStk_{\mathrm{laft}}$ 中概形的、平展的、满的映射。特别地，它是 fppf 拓扑的覆盖，于是断言由 $\IndCoh$ 的 fppf 下降推出，见 [IndCoh, Corollary 10.4.5]。

于是，设 $S'\to S$ 是固有满映射。考虑 $\PreStk_{\mathrm{laft}}$ 中的双单纯对象
\[
(S'^\bullet/S)_\star/(S'^{\bullet}_{\dR}/S_{\dR}),
\]
即 $S'\to S$ 的 Čech 神经的逐项无穷小群胚。具体地说，它是这样的双单纯对象：其 $(p,q)$-单形由映射
\[
S'^p/S\to S'^p_{\dR}/S_{\dR}
\]
的 Čech 神经的 $q$-单形给出；因此 $\star$ 标记指标 $q$，$\bullet$ 标记指标 $p$。

由命题 3.1.3，只需证明复合函子
\begin{equation}
\Crys^r(S):=\IndCoh(S_{\dR})\to\Tot\big(\IndCoh((S'^\bullet/S)_{\dR})\big)\to
\Tot\big(\IndCoh((S'^\bullet/S)_\star/(S'^\bullet/S)_{\dR})\big)\tag{3.2}
\end{equation}
是等价。

然而注意，$\PreStk_{\mathrm{laft}}$ 的双单纯对象有典范同构
\[
(S'^\bullet/S)_\star/(S'^{\bullet}_{\dR}/S_{\dR})\simeq(S'^\star/S_{\dR})_\bullet/(S'^\star/S_{\dR}),
\]
后者是余单纯对象
\[
(S'^\star/S_{\dR})\to(S^\star/S_{\dR})
\]
的逐项 Čech 神经。

(3.2) 中的映射可以改写为
\[
\Crys^r(S):=\IndCoh(S_{\dR})\to\Tot\big(\IndCoh(S^\bullet/S_{\dR})\big)\to\Tot\big(\IndCoh((S'^\star/S_{\dR})_\bullet/(S^\star/S_{\dR}))\big).
\]
再次应用命题 3.1.3，我们得到：只需证明对每个 $i$，映射
\[
\IndCoh(S^i/S_{\dR})\to\Tot\Big(\IndCoh\big((S'^i/S'^i_{\dR})_\bullet/(S^i/S_{\dR})\big)\Big)
\]
是等价。

然而注意，映射
\[
S'^i/S_{\dR}\to S^i/S_{\dR}
\]
是 ind-固有的且满。于是断言由 [GR1, Lemma 2.10.3] 推出。
\end{proof}

\ob{3.2.3}{n:3.2.3}考虑范畴 $\DGSch_{\mathrm{aff}}^{\mathrm{aft}}$ 上的 fppf 拓扑，它由 $\DGSch_{\mathrm{aff}}$ 上的 fppf 拓扑诱导（见 [GL:Stacks, Sect.~2.2]）。注意每个 fppf 覆盖特别是 h-覆盖。因此得到：

\begin{othm}{推论}{3.2.4}{n:3.2.4}
函子 $\Crys^r_{\DGSch_{\mathrm{aff}}^{\mathrm{aft}}}$ 满足 fppf 下降。
\end{othm}

如同 [IndCoh, Theorem 8.3.2]，fppf 下降是 Nisnevich 下降与有限平坦下降的结合\footnote{这个观察由 J. Lurie 向我们解释。}。特别地，我们在命题 3.2.2 的证明中建立 fppf 下降时，并没有用到「每个 fppf 覆盖也是 h-覆盖」这一事实。

\ob{3.2.5}{n:3.2.5}$\DGSch_{\mathrm{aff}}^{\mathrm{aft}}$ 上的 fppf（相应地，h-）拓扑诱导全子范畴
\[
{}^{<\infty}\DGSch_{\mathrm{aff}}^{\mathrm{aft}}\subset\DGSch_{\mathrm{aff}}^{\mathrm{aft}}
\]
上的 fppf（相应地，h-）拓扑。

命题 3.2.2 蕴含：

\begin{othm}{推论}{3.2.6}{n:3.2.6}
函子
\[
\Crys^r_{{}^{<\infty}\DGSch_{\mathrm{aff}}^{\mathrm{aft}}}:=\Crys^r_{\PreStk_{\mathrm{laft}}}|_{{}^{<\infty}\DGSch_{\mathrm{aff}}^{\mathrm{aft}}}:({}^{<\infty}\DGSch_{\mathrm{aff}}^{\mathrm{aft}})^{\mathrm{op}}\to\DGCat
\]
在 ${}^{<\infty}\DGSch_{\mathrm{aff}}^{\mathrm{aft}}$ 上满足 h-下降，特别地满足 fppf 下降。
\end{othm}

于是由 [Lu0, Corollary 6.2.3.5] 得到：

\begin{othm}{推论}{3.2.7}{n:3.2.7}
设 $Y_1\to Y_2$ 是 $\PreStk_{\mathrm{laft}}$ 中的映射，且在 h-拓扑下是满射。则自然映射
\[
\Crys(Y_2)\to\Tot\big(\Crys(Y^1_\bullet/Y_2)\big)
\]
是等价。
\end{othm}

\section{右晶体的归纳函子}

\ob{3.3.1}{n:3.3.1}记 $p_s,p_t$ 为两个投影
\[
X\times_{X_{\dR}}X\rightrightarrows X.
\]
注意映射 $p_i$（$i=s,t$）是 ind-固有的。因此函子 $p_i^!$ 有左伴随 $(p_i)_*^{\IndCoh}$，见 [GR1, Corollary 2.8.3]。

\begin{othm}{命题}{3.3.2}{n:3.3.2}
(a) 遗忘函子
\[
\oblv^r_X:\Crys(X)\to\IndCoh(X)
\]
有左伴随，记为 $\indf^r_X$。

(b) 有函子的典范同构
\[
\oblv^r_X\circ\indf^r_X\simeq(p_t)_*^{\IndCoh}\circ(p_s)^!.
\]

(c) 伴随对
\[
\indf^r_X:\IndCoh(X)\rightleftarrows\Crys^r(X):\oblv^r_X
\]
是单子的；也就是说，从 $\Crys^r(X)$ 到 $\IndCoh(X)$ 上关于单子 $\oblv^r_X\circ\indf^r_X$ 的模范畴的自然函子是等价。
\end{othm}

\begin{proof}
由命题 3.1.3 与 [Lu2, Theorem 6.2.4.2]，只需证明余单纯范畴
\[
\IndCoh(X_\bullet/X_{\dR})
\]
满足 Beck--Chevalley 条件，即：对每个 $n$，余面映射
\[
d^0:\IndCoh(X^n/X_{\dR})\to\IndCoh(X^{n+1}/X_{\dR})
\]
有左伴随，记为 $t^0$；并且对 $\Delta$ 中每个映射 $[m]\to[n]$，图表
\[
\begin{tikzcd}[column sep=large, row sep=large]
 \IndCoh(X^m/X_{\dR}) & \IndCoh(X^{m+1}/X_{\dR}) \arrow[l, "t^0"'] \\
 \IndCoh(X^n/X_{\dR}) \arrow[u] & \IndCoh(X^{n+1}/X_{\dR}) \arrow[l, "t^0"'] \arrow[u]
\end{tikzcd}
\]
（它先验地只交换到自然变换）实际上交换。

在这一情形，Beck--Chevalley 条件归结为 DG ind-概形间 ind-固有态射的 $*$-推前与 $!$-拉回之间的伴随性与基变换，由 [GR1, Proposition 2.9.2] 给出。
\end{proof}

\begin{othm}{推论}{3.3.3}{n:3.3.3}
范畴 $\Crys^r(X)$ 是紧生成的。
\end{othm}

\begin{proof}
紧生成元的集合由把 $\indf^r_X$ 作用于 $\IndCoh(X)$ 的紧生成元 $\Coh(X)$ 得到（见 [GR1, Corollary 2.4.4]）。
\end{proof}

\section{左晶体的归纳函子与无穷小群胚}

\ob{3.4.1}{n:3.4.1}由引理 1.2.4，对光滑经典概形 $X$，命题 3.1.3 的类比对左晶体成立，即函子
\begin{equation}
\Crys^l(X)=\QCoh(X_{\dR})\to\Tot\big(\QCoh(X^\bullet/X_{\dR})\big)\tag{3.3}
\end{equation}
是等价。

由命题 3.1.3，右晶体的类似断言对任意 DG 概形（乃至 DG ind-概形）$X$ 成立。然而，左晶体并非如此。

\ob{3.4.2}{n:3.4.2}我们断言：

\begin{othm}{命题}{3.4.3}{n:3.4.3}
若 DG 概形 $X$ 最终余连接，则函子 (3.3) 是等价。
\end{othm}

\begin{othm}{注}{3.4.4}{n:3.4.4}
可以证明，命题的断言对任意局部几乎有限型的 DG 概形 $X$ 都成立。但证明更为复杂。此外，由引理 2.1.7，该断言对任意经典形式光滑的预栈也成立。
\end{othm}

\begin{othm}{例}{3.4.5}{n:3.4.5}
考虑 DG 概形 $X=\Spec(k[\alpha])$，其中 $\alpha$ 位于次数 $-2$。这是构造涉及 $\Crys^l(X)$ 的断言的反例时值得记住的好例子。
\end{othm}

\begin{annote}
这个例子值得记住的原因：$\alpha$ 落在次数 $-2$，使 $k[\alpha]$ 在负次数有无界同伦群，从而 $X$ 不最终余连接。直观上，此时 $X$ 沿自身主对角线的形式完备化不再是 $X$ 本身，$\Tot(\QCoh(X^\bullet/X_{\dR}))$ 会比 $\Crys^l(X)$ 多出对象；这正是命题 3.4.3 需要最终余连接假设的原因。\hfill【译者按（上下文推断，无外部依据）】
\end{annote}


\begin{proof}[命题 3.4.3 的证明]
有函子的交换图
\[
\begin{tikzcd}[column sep=large, row sep=large]
 \Crys^l(X) \arrow[r] \arrow[d, "\Upsilon_X"'] & \Tot\big(\QCoh(X^\bullet/X_{\dR})\big) \arrow[d, "\Tot(\Upsilon_{X^\bullet/X_{\dR}})"] \\
 \Crys^r(X) \arrow[r] & \Tot\big(\IndCoh(X^\bullet/X_{\dR})\big),
\end{tikzcd}
\]
其中左竖直映射与下水平映射是等价。于是得到：$\Crys^l(X)$ 是 $\Tot(\QCoh(X^\bullet/X_{\dR}))$ 的收缩。

回忆：若 $Z$ 是最终余连接的 DG 概形，则函子
\[
\Upsilon_Z:\QCoh(Z)\to\IndCoh(Z)
\]
是全忠实的（见 [IndCoh, Corollary 9.6.3]）。于是由 [GR1, Propositions 7.1.3 and 7.4.5]，对最终余连接 DG 概形沿 Zariski 闭子集的形式完备化，同样的结论也成立。因此函子
\[
\Upsilon_{X^i/X_{\dR}}:\QCoh(X^i/X_{\dR})\to\IndCoh(X^i/X_{\dR})
\]
全忠实。于是上述交换图中的函子 $\Tot(\Upsilon_{X^\bullet/X_{\dR}})$ 也全忠实。但恒等函子是它的收缩，故它本质满。
\end{proof}

\ob{3.4.6}{n:3.4.6}对 DG 概形 $X$，定义函子
\[
\indf^l_X:\QCoh(X)\to\Crys^l(X)
\]
为
\[
\indf^l_X:=(\Upsilon_{X_{\dR}})^{-1}\circ\indf^r_X\circ\Upsilon_X.
\]
我们断言：

\begin{othm}{引理}{3.4.7}{n:3.4.7}
若 $X$ 是最终余连接的 DG 概形，则函子 $(\indf^l_X,\oblv^l_X)$ 互为伴随。
\end{othm}

\begin{othm}{注}{3.4.8}{n:3.4.8}
若去掉 $X$ 最终余连接的假设，引理的断言不成立。事实上，此时函子 $\indf^l_X$ 不保持紧性。
\end{othm}

\begin{proof}[引理 3.4.7 的证明]
回忆（见 [IndCoh, Sect.~9.6.6]）：对最终余连接的 $X$，函子 $\Upsilon_X$ 有右伴随，记为 $\Xi^{\vee}_X$；此外，$\Upsilon_X$ 本身全忠实。

由此得到 $\indf^l_X$ 的右伴随为
\[
\Xi^{\vee}_X\circ\oblv^r_X\circ\Upsilon_{X_{\dR}}\simeq\Xi^{\vee}_X\circ\Upsilon_X\circ\oblv^l_X\simeq\oblv^l_X,
\]
这正是所要的。
\end{proof}

在引理 3.4.7 的证明过程中，我们还看到：

\begin{othm}{引理}{3.4.9}{n:3.4.9}
函子 $\oblv^l_X$ 典范同构于
\[
\Xi^{\vee}_X\circ\oblv^r_X\circ\Upsilon_{X_{\dR}}.
\]
\end{othm}

\ob{3.4.10}{n:3.4.10}现在断言：

\begin{othm}{命题}{3.4.11}{n:3.4.11}
设 $X$ 是最终余连接的 DG 概形。则伴随对
\[
\indf^l_X:\QCoh(X)\rightleftarrows\Crys^l(X):\oblv^l_X
\]
是单子的；也就是说，从 $\Crys^l(X)$ 到 $\QCoh(X)$ 上关于单子 $\oblv^l_X\circ\indf^l_X$ 的模范畴的自然函子是等价。
\end{othm}

\begin{proof}
需要说明 Barr--Beck--Lurie 定理的条件成立。函子 $\oblv^l_X$ 连续，因此与所有余极限交换。$\oblv^l_X$ 保守这一事实由引理 2.2.6 给出。
\end{proof}

\chapter{晶体上的 t-结构}
\markboth{第 4 章\ 晶体上的 t-结构}{第 4 章\ 晶体上的 t-结构}

晶体范畴有两个自然的 t-结构，分别与左实现和右实现相容。右实现的主要优点之一，是与它相容的 t-结构性质好得多。

本节将定义这两个 t-结构，并证明它们的一些基本性质。这些性质包括：归纳/遗忘函子的左/右 t-正合性与上同调幅度的有界性；DG 概形的 $\Crys$ 的左完备性；与 t-结构心的导出范畴的关系。

\section{左 t-结构}

本小节中，我们不作预栈局部几乎有限型的假设。

\ob{4.1.1}{n:4.1.1}回忆 [GL:QCoh, Sec.~1.2.3]：对任意预栈 $Z$，范畴 $\QCoh(Z)$ 有典范 t-结构，其特征如下：对象 $F\in\QCoh(Z)$ 属于 $\QCoh(Z)^{\leq 0}$，当且仅当对每个 $S\in\DGSch_{\mathrm{aff}}$ 及映射 $\varphi:S\to Z$，有
\[
\varphi^*(F)\in\QCoh(S)^{\leq 0}.
\]
特别地，取 $Z=Y_{\dR}$，我们得到 $\Crys^l(Y)$ 上的典范 t-结构，称之为「左 t-结构」。

按定义，函子
\[
\oblv^l:\Crys^l(X)\to\QCoh(X)
\]
对左 t-结构是右 t-正合的。

\ob{4.1.2}{n:4.1.2}一般说来，左 t-结构的性质相当差。不过，我们有如下断言：

\begin{othm}{命题}{4.1.3}{n:4.1.3}
设 $Y$ 是经典形式光滑的预栈。则
\[
M\in\Crys^l(Y)^{\leq 0}\iff \oblv^l_Y(M)\in\QCoh(Y)^{\leq 0}.
\]
\end{othm}

\begin{proof}
需要证明：若 $M\in\Crys^l(Y)$ 满足 $\oblv^l_Y(M)\in\QCoh(Y)^{\leq 0}$，则 $M\in\Crys^l(Y)^{\leq 0}$。也就是说，需要证明对每个 $S\in\DGSch_{\mathrm{aff}}$ 及 $\varphi:S\to Y_{\dR}$，有 $\varphi^*(M)\in\QCoh(S)^{\leq 0}$。

记 $Y_\bullet/Y_{\dR}$ 为映射 $p_{\dR,Y}:Y\to Y_{\dR}$ 的 Čech 神经。由引理 1.2.4，存在映射 $\varphi':S\to Y$ 及同构 $\varphi\simeq p_{\dR,Y}\circ\varphi'$。断言现在由 $\varphi'^*$ 右 t-正合推出。
\end{proof}

\section{右 t-结构}

从此刻到本节末，我们恢复所有预栈局部几乎有限型的假设（除明确说明外）。

本小节专门讨论 DG 概形的情形。

\ob{4.2.1}{n:4.2.1}设 $X$ 是 DG 概形。回忆范畴 $\IndCoh(X)$ 有自然的、与滤过余极限相容的 t-结构，见 [IndCoh, Sect.~1.2]。

它的特征性质是：$\IndCoh(X)$ 中的对象是连接的（即落在 $\IndCoh(X)^{\leq 0}$ 中），当且仅当它在函子 $\Psi_X:\IndCoh(X)\to\QCoh(X)$ 下的像连接。

\ob{4.2.2}{n:4.2.2}我们通过宣布
\[
M\in\Crys^r(X)^{\geq 0}\iff \oblv^r_X(M)\in\IndCoh(X)^{\geq 0}
\]
来定义 $\Crys^r(X)$ 上的右 t-结构。

下文中，我们把 $\Crys^r(X)$ 上的右 t-结构称为晶体上的「那个」t-结构。换言之，默认情况下我们考虑的 t-结构是右 t-结构。按构造，这个 t-结构也与滤过余极限相容，因为 $\oblv^r_X$ 连续。

\begin{annote}
「那个」t-结构这个叫法来自 §0.2.1：在光滑经典概形上，右实现给出的 t-结构在 Riemann--Hilbert 对应下就是右 D-模的反常 t-结构，而左 t-结构只反映拟凝聚层的通常 t-结构。因此凡不加限定时，本文默认使用右 t-结构。\hfill【来源：原书 §0.2.1、§4.2.2（A）】
\end{annote}


\ob{4.2.3}{n:4.2.3}我们断言 $\Crys^r(X)$ 上的右 t-结构是 Zariski 局部的，即一个对象连接/余连接，当且仅当它限制到某个 Zariski 覆盖后具有该性质。事实上，这由 $\IndCoh(X)$ 上 t-结构的相应性质推出，见 [IndCoh, Corollary 4.2.3]。

\ob{4.2.4}{n:4.2.4}右 t-结构与 Kashiwara 引理。设 $i:X\to Z$ 是 DG 概形的闭浸入。记 $i_{\dR,*}$ 为函子 $\Crys^r(X)\to\Crys^r(Z)$，它等于复合
\[
\Crys^r(X)\xrightarrow{(i'^{\dagger,r})^{-1}}\Crys^r(Z)_X\hookrightarrow\Crys^r(Z),
\]
按构造，它是 $i^{\dagger,r}$ 的左伴随。

有：

\begin{othm}{命题}{4.2.5}{n:4.2.5}
函子 $i_{\dR,*}:\Crys^r(X)\to\Crys^r(Z)$ 是 t-正合的。
\end{othm}

\begin{proof}
注意全子范畴
\[
\Crys^r(Z)_X\subset\Crys^r(Z)
\]
与 t-结构相容，因为它是函子 $j^{\dagger,r}$ 的核，而后者 t-正合（这里 $j$ 表示开浸入 $Z-X\hookrightarrow Z$）。

于是只需证明函子
\[
i'^{\dagger,r}:\Crys^r(Z)_X\to\Crys^r(X)
\]
t-正合。

因此需要证明：对 $M\in\Crys^r(Z)_X$，有
\[
\oblv^r_Z(M)\in\IndCoh(Z)^{>0}\iff \oblv^r_X(i^{\dagger,r}(M))\in\IndCoh(X)^{>0}.
\]
回忆记号
\[
\IndCoh(Z)_X:=\ker\big(j^!:\IndCoh(Z)\to\IndCoh(Z-X)\big).
\]
只需证明下述引理：
\end{proof}

\begin{othm}{引理}{4.2.6}{n:4.2.6}
设 $i:X\to Z$ 是闭浸入。则对 $F\in\IndCoh(Z)_X$ 有
\[
F\in\IndCoh(Z)^{>0}\iff i^!(F)\in\IndCoh(X)^{>0}.
\]
\end{othm}

\begin{proof}[引理 4.2.6 的证明]
「$\Rightarrow$」由 $i^!$ 左 t-正合推出（它是 t-正合函子 $i_*^{\IndCoh}$ 的右伴随）。

对反向蕴含，注意全子范畴
\[
\IndCoh(Z)_X\subset\IndCoh(Z)
\]
也与 t-结构相容，因为它是 t-正合函子 $j^!$ 的核。此外，由 [IndCoh, Proposition 4.1.7(b)]，$\IndCoh(Z)_X$ 上的 t-结构由
\[
\Coh(Z)_X:=\ker\big(j^*:\Coh(Z)\to\Coh(Z-X)\big)
\]
上的 t-结构生成。

设 $F\in\IndCoh(Z)_X$ 满足 $i^!(F)\in\IndCoh(X)^{>0}$。需要证明 $F$ 与 $(\Coh(Z)_X)^{\leq 0}$ 右正交。按假设，$F$ 与 $\Coh(X)^{\leq 0}$ 在
\[
\Coh(X)\xrightarrow{i_*}\Coh(Z)_X\to\IndCoh(Z)_X
\]
下的本质像右正交。然而容易看出，$(\Coh(Z)_X)^{\leq 0}$ 的每个对象都可以由 $\Coh(X)^{\leq 0}$ 的本质像中的对象经过有限次逐步扩张得到，这蕴含所需断言。
\end{proof}

\begin{othm}{推论}{4.2.7}{n:4.2.7}
若 DG 概形的映射 $X_1\to X_2$ 诱导同构
\[
\clred X_1\to\clred X_2,
\]
则 $\Crys^r(X_1)\simeq\Crys^r(X_2)$ 上相应的 t-结构重合。
\end{othm}

\ob{4.2.8}{n:4.2.8}设 $X$ 是 DG 概形。按构造，遗忘函子 $\oblv^r_X$ 左 t-正合。于是由伴随性，函子 $\indf^r_X$ 右 t-正合。

现在断言：

\begin{othm}{命题}{4.2.9}{n:4.2.9}
函子 $\indf^r_X$ 是 t-正合的。
\end{othm}

\begin{proof}
只需证明复合 $\oblv^r_X\circ\indf^r_X$ 左 t-正合。我们由命题 3.3.2(b) 推出这一点：

函子 $p_s^!$ 左 t-正合（例如对 $\Delta_X:X\to X\times X$ 应用引理 4.2.6）。函子 $(p_t)_*^{\IndCoh}$ 由 [GR1, Lemma 2.7.11] 左 t-正合（事实上 t-正合）。
\end{proof}

\ob{4.2.10}{n:4.2.10}现在断言：

\begin{othm}{命题}{4.2.11}{n:4.2.11}
(a) 若 $X$ 是光滑经典概形，则 $\oblv^r_X$ t-正合。

(b) 对拟紧 DG 概形 $X$，函子 $\oblv^r_X$ 有界上同调幅度。
\end{othm}

\begin{proof}
设 $X$ 是光滑经典概形。按 $\Crys^r(X)$ 上 t-结构的定义，$\IndCoh(X)^{\leq 0}$ 在 $\indf^r_X$ 下的本质像通过取余极限生成 $\Crys^r(X)^{\leq 0}$。因此，为证明 $\oblv^r_X$ 右 t-正合，只需对函子 $\oblv^r_X\circ\indf^r_X$ 证明同样的事情。我们将由命题 3.3.2(b) 推出这一点：

可以写
\[
(X\times X)^{\wedge}_X\simeq\varinjlim_n X_n,
\]
其中 $X_n\to X\times X$ 是对角线的第 $n$ 个无穷小邻域。于是由 [GR1, Equation (2.2)]，
\[
(p_t)_*^{\IndCoh}\circ p_s^!\simeq\varinjlim_n (p_t\circ i_n)_*^{\IndCoh}\circ(p_s\circ i_n)^!.
\]
现在，每个函子 $(p_t\circ i_n)_*^{\IndCoh}$ 由 [GR1, Lemma 2.7.11] t-正合；每个函子 $(p_s\circ i_n)^!$ t-正合，因为 $p_s\circ i_n:X_n\to X$ 有限平坦。

现在设 $X$ 是拟紧 DG 概形，来证明 $\oblv^r_X$ 有界上同调幅度。问题容易化归到 $X$ 仿射的情形；取闭浸入 $i:X\hookrightarrow Z$，其中 $Z$ 光滑。由命题 4.2.5 与 (a)，只需证明函子 $i^!:\IndCoh(Z)\to\IndCoh(X)$ 有界上同调幅度，而后者容易由 $Z$ 正则推出。
\end{proof}

\section{ind-概形上晶体的右 t-结构}

\ob{4.3.1}{n:4.3.1}设 $X$ 是 DG ind-概形。固定 $X$ 的一个表示
\begin{equation}
X=\varinjlim_\alpha X_\alpha\tag{4.1}
\end{equation}
如 [GR1, Prop.~1.7.6]。对每个 $\alpha$，记 $i_\alpha$ 为相应的闭浸入 $X_\alpha\to X$；对每个 $\alpha_1\to\alpha_2$，记 $i_{\alpha_1,\alpha_2}$ 为闭浸入 $X_{\alpha_1}\to X_{\alpha_2}$。

有：
\[
\Crys^r(X)\simeq\varprojlim_\alpha\Crys^r(X_\alpha),
\]
其中对 $\alpha_1\to\alpha_2$，函子 $\Crys^r(X_{\alpha_2})\to\Crys^r(X_{\alpha_1})$ 由 $i^{\dagger,r}_{\alpha_1,\alpha_2}$ 给出。

于是由 [GL:DG, Sect.~1.3.3]，
\[
\Crys^r(X)\simeq\varinjlim_\alpha\Crys^r(X_\alpha),
\]
其中对 $\alpha_1\to\alpha_2$，函子 $\Crys^r(X_{\alpha_1})\to\Crys^r(X_{\alpha_2})$ 由 $(i_{\alpha_1,\alpha_2})_{\dR,*}$ 给出。

特别地，对每个 $\alpha$ 得到一对伴随函子
\[
(i_\alpha)_{\dR,*}:\Crys^r(X_\alpha)\rightleftarrows\Crys^r(X):i^{\dagger,r}_\alpha.
\]

\ob{4.3.2}{n:4.3.2}回忆 [GR1, Sect.~2.5]：$\IndCoh(X)$ 有自然的、与滤过余极限相容的 t-结构。

利用 $\IndCoh(X)$ 上的这个 t-结构，可以定义 $\Crys^r(X)$ 上的右 t-结构。即
\[
M\in\Crys^r(X)^{\geq 0}\iff \oblv^r_X(M)\in\IndCoh(X)^{\geq 0}.
\]
由于 $\oblv^r_X$ 保持余极限，这个 t-结构与滤过余极限相容。利用表示 (4.1)，可以更具体地描述这个 t-结构，方式类似于 [GR1, Lemma 2.5.3] 中对 $\IndCoh(X)$ 上 t-结构的描述：

\begin{othm}{引理}{4.3.3}{n:4.3.3}
在上述情形下有：

(a) 对象 $F\in\Crys^r(X)$ 属于 $\Crys^r(X)^{\geq 0}$，当且仅当对每个 $\alpha$，对象 $i^{\dagger,r}_\alpha(F)\in\Crys^r(X_\alpha)$ 属于 $\Crys^r(X_\alpha)^{\geq 0}$。

(b) 范畴 $\Crys^r(X)^{\leq 0}$ 由函子 $(i_\alpha)_{\dR,*}\big(\Crys^r(X_\alpha)^{\leq 0}\big)$ 的本质像在余极限下生成。
\end{othm}

\begin{proof}
(a) 由定义与 [GR1, Lemma 2.5.3(a)] 推出。(b) 由 (a) 形式地推出。
\end{proof}

\ob{4.3.4}{n:4.3.4}设 $i:X\to X$ 是把 DG 概形闭浸入 DG ind-概形。由与 [GR1, Lemma 2.5.5] 完全相同的论证，有：

\begin{othm}{引理}{4.3.5}{n:4.3.5}
函子 $i_{\dR,*}$ 是 t-正合的。
\end{othm}
\ob{4.3.6}{n:4.3.6}作为上述 DG ind-概形上右晶体 t-结构行为的一个说明，考虑如下情形。设 $i:X\to Z$ 是拟紧 DG 概形的闭浸入。记 $Y$ 为 $X$ 在 $Z$ 中的形式完备化，看作 $\DGIndSch_{\mathrm{laft}}$ 的对象；记 $i':X\to Y$ 为相应的映射。

我们断言：

\begin{othm}{引理}{4.3.7}{n:4.3.7}
由同构 $i'_{\dR}:X_{\dR}\to Y_{\dR}$ 诱导的等价 $\Crys^r(X)\simeq\Crys^r(Y)$ 与 t-结构相容。
\end{othm}

\begin{proof}[证明一]
由命题 4.2.5 以及 [GR1, Proposition 7.4.5] 中的等价
\[
\IndCoh(Y)\simeq\IndCoh(Z)_Y
\]
与 t-结构相容（见 [GR1, Lemma 7.4.8]）推出。
\end{proof}

\begin{proof}[证明二]
由交换图
\[
\begin{tikzcd}[column sep=large, row sep=large]
 \Crys^r(X) & \Crys^r(Y) \arrow[l, "i'^{\dagger,r}"'] \\
 \IndCoh(X) \arrow[u, "\oblv^r_X"] & \IndCoh(Y) \arrow[l, "i'^!"] \arrow[u, "\oblv^r_Y"']
\end{tikzcd}
\]
得到：只需证明对 $F\in\IndCoh(Y)$ 有
\[
F\in\IndCoh(Y)^{>0}\iff i'^!(F)\in\IndCoh(X)^{>0},
\]
这由 [GR1, Lemma 7.4.8] 与引理 4.2.6 形式地推出（也可以直接证明）。
\end{proof}

\section{左 t-结构的进一步性质}

\ob{4.4.1}{n:4.4.1}首先描述光滑经典概形情形下晶体上左、右 t-结构之间的关系。

\begin{othm}{命题}{4.4.2}{n:4.4.2}
设 $X$ 是 $n$ 维光滑经典概形。则
\[
F\in\Crys^l(X)^{\leq 0}\iff F\in\Crys^r(X)^{\leq -n}.
\]
也就是说，左 t-结构与右 t-结构相差 $X$ 的维数这样一个平移。
\end{othm}

\begin{proof}
回忆两个遗忘函子由交换图
\[
\begin{tikzcd}[column sep=large, row sep=large]
 & \Crys(X) \arrow[dl, "\oblv^l_X"'] \arrow[dr, "\oblv^r_X"] & \\
 \QCoh(X) \arrow[rr, "\Upsilon_X"'] & & \IndCoh(X)
\end{tikzcd}
\]
相联系。

当 $X$ 是 $n$ 维光滑经典概形时，函子 $\Upsilon_X$ 是等价，并把 $\QCoh(X)^{\leq 0}$ 同构地送到 $\IndCoh(X)^{\leq -n}$。断言现在由命题 4.1.3 与命题 4.2.11(a)，结合 $\oblv^r_X$ 保守这一事实推出。
\end{proof}

\ob{4.4.3}{n:4.4.3}下一个命题对任意 DG 概形 $X$ 比较 $\Crys(X)$ 上的「左」t-结构与「右」t-结构。

\begin{othm}{命题}{4.4.4}{n:4.4.4}
设 $X$ 拟紧。则恒等函子
\[
\Crys^l(X)\to\Crys^r(X)
\]
有界幅度，即左、右 t-结构之间的差别有界。
\end{othm}

\begin{proof}
不妨设 $X$ 仿射。设 $Z$ 是 $n$ 维光滑经典概形，$i:X\hookrightarrow Z$ 是闭浸入。我们断言：对 $M^l\in\Crys^l(X)$ 及相应的对象 $M^r\in\Crys^r(X)$，有
\begin{equation}
M^l\in(\Crys^l(X))^{\leq 0}\Rightarrow M^r\in(\Crys^r(X))^{\leq 0},\qquad
M^r\in(\Crys^r(X))^{\leq 0}\Rightarrow M^l\in(\Crys^l(X))^{\leq n}.\tag{4.2}
\end{equation}
记 $j:U\hookrightarrow Z$ 为补开浸入。记 $Y$ 为 $X$ 在 $Z$ 中的形式完备化；记 $\widehat{i}:Y\to Z$。

映射 $X\to Y$ 定义同构 $X_{\dR}\to Y_{\dR}$，从而可以认同 $\Crys^l(X)\simeq\Crys^l(Y)$。应用命题 4.1.3，得到
\begin{equation}
M^l\in(\Crys^l(X))^{\leq 0}\iff \oblv^l_Y(M^l)\in\QCoh(Y)^{\leq 0},\tag{4.3}
\end{equation}
其中 $\QCoh(Y)$ 上的 t-结构是 §4.1.1 中的那个。

考虑子范畴 $\QCoh(Z)_X\subset\QCoh(Z)$，按定义它等于
\[
\ker\big(j^*:\QCoh(Z)\to\QCoh(U)\big).
\]
这个子范畴与 $\QCoh(Z)$ 上的 t-结构相容，因为函子 $j^*$ t-正合。

回忆（见 [GR1, Proposition 7.1.3]）：函子 $\widehat{i}_*$ 定义等价
\[
\QCoh(Z)_X\to\QCoh(Y).
\]
设 $F$ 是 $\QCoh(Z)_X$ 中对应于 $\oblv^l_Y(M^l)\in\QCoh(Y)$ 的对象。有
\[
\Upsilon_Z(F)\simeq\oblv^r_Z(i_{\dR,*}(M^r)).
\]
由于函子 $i_{\dR,*}$ t-正合（命题 4.2.5），且 $\Upsilon_Z$ 把上同调次数平移 $[-n]$，得到
\begin{equation}
M^r\in(\Crys^r(X))^{\leq 0}\iff F\in(\QCoh(Z)_X)^{\leq n}.\tag{4.4}
\end{equation}
结合 (4.3) 与 (4.4)，(4.2) 中的蕴含由下述断言推出：
\end{proof}

\begin{othm}{引理}{4.4.5}{n:4.4.5}
等价 $\widehat{i}_*:\QCoh(Z)_X\simeq\QCoh(Y)$ 关于 $\QCoh(Z)_X$ 上从 $\QCoh(Z)$ 继承的 t-结构与 §4.1.1 中 $\QCoh(Y)$ 上的 t-结构，具有如下性质：

(a) 若 $F\in(\QCoh(Z)_X)^{\leq 0}$，则 $\widehat{i}_*(F)\in\QCoh(Y)^{\leq 0}$。

(b) 若 $\widehat{i}_*(F)\in\QCoh(Y)^{\leq 0}$，则 $F\in(\QCoh(Z)_X)^{\leq n}$。
\end{othm}

\begin{proof}[引理 4.4.5 的证明]
(a) 由函子 $\widehat{i}^*$ 右 t-正合推出。

为证明 (b)，注意范畴 $\QCoh(Y)^{\leq 0}$ 在取余极限下由 $\QCoh(Z)^{\leq 0}$ 在函子 $\widehat{i}_*$ 下的本质像生成，见 [GR1, Proposition 7.3.5]。因此只需证明函子
\[
\widehat{i}_*:\QCoh(Z)\to\QCoh(Y)\simeq\QCoh(Z)_X
\]
的上同调幅度被 $n$ 界住。然而上述函子是嵌入
\[
\QCoh(Z)_X\hookrightarrow\QCoh(Z)
\]
的右伴随，由
\[
F'\mapsto\Cone\big(F'\to j_*j^*(F')\big)[-1]
\]
给出。现在 $j^*$ t-正合，$j_*$ 的上同调幅度被 $n-1$ 界住。这蕴含所需断言。
\end{proof}

\ob{4.4.6}{n:4.4.6}设 $X$ 是任意拟紧 DG 概形。有：

\begin{othm}{命题}{4.4.7}{n:4.4.7}
(a) 函子 $\oblv^l_X:\Crys(X)\to\QCoh(X)$ 有界上同调幅度。

(b) 若 $X$ 最终余连接，则函子 $\indf^l_X:\QCoh(X)\to\Crys(X)$ 的上同调幅度上有界。
\end{othm}

\begin{proof}
对 (a)，不妨设 $X$ 仿射，并取闭浸入 $i:X\hookrightarrow Z$，其中 $Z$ 是光滑经典概形。此时断言由命题 4.2.5 与函子
\[
i^*:\QCoh(Z)\to\QCoh(X)
\]
有界上同调幅度推出。

(b) 由 (a) 与 $(\indf^l_X,\oblv^l_X)$-伴随性推出。
\end{proof}

\begin{othm}{注}{4.4.8}{n:4.4.8}
(b) 中 $X$ 最终余连接的假设是本质的；否则可以用例 3.4.5 中的 DG 概形给出反例。此外，容易证明：$\indf^l_X$ 的上同调幅度下有界，当且仅当 $X$ 是 Gorenstein 的（见引理 4.6.12）。
\end{othm}

\section{左完备性}

\ob{4.5.1}{n:4.5.1}设 $X$ 是仿射光滑经典概形。我们注意到，此时范畴 $\Crys^r(X)$ 含有一个典范对象
\[
\indf^r_X(\mathcal{O}_X),
\]
它落在 t-结构的心中（见命题 4.2.9），并且是投射的，即
\[
H^0(N)=0\Rightarrow\Hom_{\Crys^r(X)}\big(\indf^r_X(\mathcal{O}_X),N\big)=0.
\]
此外，$\indf^r_X(\mathcal{O}_X)$ 是 $\Crys^r(X)$ 的紧生成元。这蕴含：

\begin{othm}{推论}{4.5.2}{n:4.5.2}
设 $X$ 是仿射光滑经典概形。则范畴 $\Crys^r(X)$ 在其 t-结构下左完备。
\end{othm}

\ob{4.5.3}{n:4.5.3}上述推论对任意 DG 概形 $X$ 蕴含左完备性：

\begin{othm}{推论}{4.5.4}{n:4.5.4}
对任意 DG 概形 $X$，范畴 $\Crys^r(X)$ 在「右」t-结构下左完备。
\end{othm}

\begin{proof}
首先注意，左完备性是 Zariski 局部的（用与 [GL:QCoh, Proposition 5.2.4] 相同的论证证明）。于是不妨设 $X$ 仿射。取闭浸入 $i:X\hookrightarrow Z$，其中 $Z$ 是光滑经典概形。此时断言由以下事实形式地推出：函子 $i_{\dR,*}$ 连续、全忠实（由命题 2.5.6）、t-正合（由命题 4.2.5），且 $\Crys^r(Z)$ 左完备（由上一推论）。

这里还有一个替代论证：

由推论 4.2.7，可以假设 $X$ 最终余连接。此时函子 $\oblv^l_X$ 与极限交换，因为它有左伴随。此外，由引理 2.2.6，$\oblv^l_X$ 保守；由命题 4.4.7，它有界上同调幅度。因此，$\QCoh(X)$ 在其 t-结构下左完备这一事实蕴含 $\Crys^r(X)$ 的相应事实。
\end{proof}

\begin{othm}{注}{4.5.5}{n:4.5.5}
右完备性不成问题：由于我们的 t-结构与滤过余极限相容，右完备性等价于 t-结构在余连接子范畴上可分，而这是显然的，因为 $\oblv^r_X$ 左 t-正合且保守，并且 $\IndCoh(X)^+\simeq\QCoh(X)^+$ 上的 t-结构具有该性质。
\end{othm}

\ob{4.5.6}{n:4.5.6}把推论 4.5.4 与命题 4.4.4 结合，得到：

\begin{othm}{推论}{4.5.7}{n:4.5.7}
对拟紧 DG 概形 $X$，范畴 $\Crys(X)$ 在「左」t-结构下也左完备。
\end{othm}

\section{「粗糙的」归纳函子与遗忘函子}

\ob{4.6.1}{n:4.6.1}设 $X$ 是 DG 概形。回忆函子 $\Psi_X$ 把范畴 $\QCoh(X)$ 与 $\IndCoh(X)$ 的左完备化等同起来（见 [IndCoh, Proposition 1.3.4]）。

由于范畴 $\Crys^r(X)$ 在其 t-结构下左完备，且函子 $\indf^r_X$ t-正合，由左完备化的泛性质得到：

\begin{othm}{推论}{4.6.2}{n:4.6.2}
函子 $\indf^r_X$ 典范地分解为
\[
\IndCoh(X)\xrightarrow{\Psi_X}\QCoh(X)\xrightarrow{{}'\indf^r_X}\Crys^r(X).
\]
\end{othm}

\ob{4.6.3}{n:4.6.3}也可以考虑函子
\[
{}'\oblv^r_X:\Crys^r(X)\to\QCoh(X),
\]
它由 $\Psi_X\circ\oblv^r_X$ 给出，其中 $\Psi_X:\IndCoh(X)\to\QCoh(X)$ 是 [IndCoh, Sect.~1.1.5] 中的函子。

显然函子 ${}'\oblv^r_X$ 有有限上同调幅度。事实上，这由 $\oblv^r_X$ 的相应事实与 $\Psi_X$ t-正合（见 [IndCoh, Lemma 1.2.2]）推出。

\begin{othm}{命题}{4.6.4}{n:4.6.4}
函子 ${}'\oblv^r_X$ 保守。
\end{othm}

\begin{proof}
断言是 Zariski 局部的，于是不妨设 $X$ 仿射。取闭浸入 $i:X\to Z$，其中 $Z$ 是光滑经典仿射概形。

记 $i^{\QCoh,!}:\QCoh(Z)\to\QCoh(X)$ 为 $i_*:\QCoh(X)\to\QCoh(Z)$ 的右伴随\footnote{尽管这与我们无关，我们仍指出 $i^{\QCoh,!}$ 连续。这是因为函子 $i_*:\QCoh(X)\to\QCoh(Z)$ 把紧对象送到紧对象（由于 $Z$ 正则，其上任一凝聚层都是完美的）。}。容易看出有函子的典范同构
\[
\Psi_X\circ i^!\simeq i^{\QCoh,!}\circ\Psi_Z.
\]
于是对 $M\in\Crys^r(Z)$，有
\[
{}'\oblv^r_X(i^{\dagger,r}(M))\simeq i^{\QCoh,!}({}'\oblv^r_Z(M)).
\]
应用 Kashiwara 引理，命题的断言由下述引理推出：
\end{proof}

\begin{othm}{引理}{4.6.5}{n:4.6.5}
函子 $i^{\QCoh,!}:\QCoh(Z)\to\QCoh(X)$ 限制到 $\QCoh(Z)_X$ 上时保守。
\end{othm}

\begin{proof}[引理 4.6.5 的证明]
需要证明函子 $i_*:\QCoh(X)\to\QCoh(Z)$ 的本质像生成 $\QCoh(Z)_X$。

首先断言 $\QCoh(Z)_X$ 由有界对象的子范畴 $(\QCoh(Z)_X)^b$ 生成。这由 $\QCoh(Z)$ 的相应事实以及嵌入 $\QCoh(Z)_X\hookrightarrow\QCoh(Z)$ 有有界上同调幅度的右伴随推出。由 dévissage得到：$\QCoh(Z)_X$ 由 $(\QCoh(Z)_X)^{\heartsuit}$ 生成，进而由 $(\QCoh(Z)_X)^{\heartsuit}\cap\Coh(Z)$ 生成。

然而显然 $(\QCoh(Z)_X)^{\heartsuit}\cap\Coh(Z)$ 的每个对象都是 $\Coh(X)^{\heartsuit}$ 的本质像中对象的有限次扩张。
\end{proof}

\begin{othm}{注}{4.6.6}{n:4.6.6}
当 $X$ 最终余连接时，我们将在下面给出命题 4.6.4 的一个更干净的证明。
\end{othm}

\ob{4.6.7}{n:4.6.7}现在假设 $X$ 最终余连接。回忆此时函子 $\Psi_X$ 有全忠实的左伴随 $\Xi_X$（见 [IndCoh, Proposition 1.5.3]）。

我们注意到：

\begin{othm}{引理}{4.6.8}{n:4.6.8}
存在典范同构 ${}'\indf^r_X\simeq\indf^r_X\circ\Xi_X$。
\end{othm}

\begin{proof}
由同构 $\indf^r_X\simeq{}'\indf^r_X\circ\Psi_X$ 与 $\Psi_X\circ\Xi_X\simeq\mathrm{Id}_{\QCoh(X)}$ 推出。
\end{proof}

\begin{othm}{推论}{4.6.9}{n:4.6.9}
函子 $({}'\indf^r_X,{}'\oblv^r_X)$ 构成伴随对。
\end{othm}

\begin{proof}
由引理 4.6.8 通过伴随性形式地推出。
\end{proof}

\begin{othm}{注}{4.6.10}{n:4.6.10}
除非 $X$ 最终余连接，函子 $({}'\indf^r_X,{}'\oblv^r_X)$ 不是伴随对。事实上，若 $X$ 不最终余连接，则函子 ${}'\indf^r_X$ 不保持紧对象：它把 $\mathcal{O}_X\in\QCoh(X)$ 送到 $\Crys^r(X)$ 的非紧对象。
\end{othm}

\begin{proof}[命题 4.6.4 的替代证明]
由推论 4.6.9，命题 4.6.4 的断言（在最终余连接情形）等价于：函子 ${}'\indf^r_X$ 的本质像生成 $\Crys^r(X)$。然而后者由 $\indf^r_X$ 的相应事实显然成立。
\end{proof}

\ob{4.6.11}{n:4.6.11}设 $X$ 是最终余连接的 DG 概形，考虑伴随函子对
\[
\Xi_X:\QCoh(X)\rightleftarrows\IndCoh(X):\Psi_X,
\]
其中 $\Xi_X$ 全忠实（见 [IndCoh, Sect.~1.4]）。

我们已经看到，函子 $\indf^r_X$ 通过余局部化函子 $\Psi_X$ 分解。然而，函子 $\oblv^r_X$ 一般并不通过 $\Xi_X$ 分解，即不取值于 $\QCoh(X)$（经 $\Xi_X$ 看作 $\IndCoh(X)$ 的全子范畴）中。

事实上，有下述结论：

\begin{othm}{引理}{4.6.12}{n:4.6.12}
（Drinfeld）函子 $\oblv^r_X$ 通过 $\QCoh(X)$ 在 $\Xi_X$ 下的本质像分解，当且仅当 $X$ 是 Gorenstein 的。
\end{othm}

回忆：DG 概形 $X$ 称为 Gorenstein 的，如果

(a) $\omega_X\in\Coh(X)$（这等价于 $X$ 最终余连接，见 [IndCoh, Proposition 9.6.11]）；

(b) 看作凝聚层时，$\omega_X$ 是分次线丛（这等价于 $\omega_X\in\QCoh(X)^{\mathrm{perf}}$，见 [IndCoh, Corollary 7.4.3]）。

\begin{annote}
条件 (b) 中的「分次线丛」是导出意义下的说法：$\omega_X$ 在同伦意义下可逆，即属于 $\QCoh(X)^{\mathrm{perf}}$。光滑概形上 $\omega_X$ 是线丛，条件自动满足；因此 Gorenstein 条件刻画的是「对偶化复形不比线丛更坏」的奇异性。\hfill【来源：原书 §4.6.11--4.6.12；[IndCoh, Sect.~9.6]（A）】
\end{annote}


\begin{proof}[引理 4.6.12 的证明]
假设 $\oblv^r_X$ 通过 $\QCoh(X)$ 分解。特别地，得到 $\omega_X\in\Coh(X)$ 落在 $\Xi_X$ 的本质像中。现在断言由 [IndCoh, Lemma 1.5.8] 推出。

对反向蕴含，把 $\oblv^r_X(M)$ 写成
\[
\Upsilon_X(M)=\oblv^l_X(M)\otimes\omega_X,
\]
其中张量积按 $\QCoh(X)$ 对 $\IndCoh(X)$ 的作用理解，见 [IndCoh, Sect.~1.4]。还要回忆：函子 $\Xi_X$ 按定义与上述 $\QCoh(X)$-作用相容。因此，若 $\omega_X$ 作为完美对象落在 $\Xi_X$ 的本质像中，则 $\oblv^l_X(M)\otimes\omega_X$ 也落在其中。
\end{proof}

\section{与 Abel 范畴的关系}

本小节中，设 $X$ 是仿射 DG 概形。我们将把范畴 $\Crys^r(X)$ 与一个更熟悉的对象联系起来。

\ob{4.7.1}{n:4.7.1}由于 $\Crys^r(X)$ 上的 t-结构与滤过余极限相容，得到 $\Crys^r(X)^{\heartsuit}$ 是 Grothendieck Abel 范畴。

利用 $\Crys^r(X)$ 在其 t-结构下右完备这一事实，由 [Lu2, Theorem 1.3.2.2] 中的箭头反转，得到典范定义的 t-正合函子
\begin{equation}
D\big(\Crys^r(X)^{\heartsuit}\big)^{+}\to\Crys^r(X)^{+},\tag{4.5}
\end{equation}
其中 $D(-)^{+}$ 表示 Abel 范畴的导出范畴的最终余连接部分。

\ob{4.7.2}{n:4.7.2}我们将证明：

\begin{othm}{命题}{4.7.3}{n:4.7.3}
函子 (4.5) 唯一地延拓为范畴的等价
\[
D\big(\Crys^r(X)^{\heartsuit}\big)\to\Crys^r(X).
\]
\end{othm}

本小节余下部分专门证明命题 4.7.3。不妨设 $X$ 是经典的。

\ob{4.7.4}{n:4.7.4}第 1 步。先假设 $X$ 是光滑经典概形。此时断言由下述事实显然：
\[
\indf^r_X(\mathcal{O}_X)
\]
是两个范畴的紧投射生成元。

\ob{4.7.5}{n:4.7.5}第 2 步。证明函子
\[
D\big(\Crys^r(X)^{\heartsuit}\big)^{+}\to\Crys^r(X)^{+}
\]
是等价。为此，只需证明每个对象 $M\in\Crys^r(X)^{\heartsuit}$ 可以嵌入一个内射对象，即一个 $I\in\Crys^r(X)^{\heartsuit}$，使得
\[
H^0(N)=0\Rightarrow\Hom_{\Crys^r(X)}(N,I)=0.
\]
取闭浸入 $i:X\hookrightarrow Z$，其中 $Z$ 是光滑经典概形。取嵌入 $i_{\dR,*}(M)\hookrightarrow J$，其中 $J$ 是 $\Crys^r(Z)$ 中的内射对象（按同样的意义）；由第 1 步，它存在。

由于函子 $\indf^r_Z$ t-正合，得到 $\oblv^r_Z(J)$ 是 $\QCoh(Z)^{\heartsuit}$ 的内射对象。这蕴含 $I:=i^{\dagger,r}(J)$ 属于 $\Crys^r(X)^{\heartsuit}$ 并具有所需性质。

\ob{4.7.6}{n:4.7.6}第 3 步。注意由推论 4.5.4，范畴 $\Crys^r(X)$ 与 $\Crys^r(X)^{+}$ 的左完备化等同。因此，只需证明典范嵌入
\[
D\big(\Crys^r(X)^{\heartsuit}\big)^{+}\hookrightarrow D\big(\Crys^r(X)^{\heartsuit}\big)
\]
把 $D(\Crys^r(X)^{\heartsuit})^{+}$ 与 $D(\Crys^r(X)^{\heartsuit})$ 的左完备化等同。

为此，只需给出 $\Crys^r(X)^{\heartsuit}$ 的一个有界 Ext 维数的生成元 $P$。

考虑对象
\[
P:=\indf^r_X(\mathcal{O}_X).
\]
它由命题 4.2.11(b) 具有所需性质。

\begin{othm}{注}{4.7.7}{n:4.7.7}
标准的论证可以把命题 4.7.3 的断言推广到对角仿射的拟紧 DG 概形 $X$。
\end{othm}

\begin{othm}{注}{4.7.8}{n:4.7.8}
一旦我们在光滑仿射经典概形上把晶体与 D-模等同起来，就会得到拟紧 DG 概形上 $\Crys^r(X)$ 的许多其他性质：例如，Abel 范畴 $\Crys^r(X)^{\heartsuit}$ 是局部 Noether 的\footnote{这里的意思是：$\Crys^r(X)^{\heartsuit}$ 由它的紧对象生成，且紧对象的子对象是紧的。}，且 $\Crys^r(X)$ 关于其 t-结构有有限上同调维数\footnote{这里的意思是：存在 $N\in\mathbb{N}$，使得对 $n>N$ 及 $\mathcal{M}_1,\mathcal{M}_2\in\Crys^r(X)^{\heartsuit}$，有 $\Hom_{\Crys^r(X)}(\mathcal{M}_1,\mathcal{M}_2[n])=0$。}。注意由命题 4.2.5，为建立这两条性质，只需对仿射光滑经典概形证明它们。
\end{othm}

\chapter{与 D-模的关系}
\markboth{第 5 章\ 与 D-模的关系}{第 5 章\ 与 D-模的关系}

本节把单子 $\oblv^r_X\circ\indf^r_X$ 与 $\oblv^l_X\circ\indf^l_X$ 同微分算子层联系起来。作为结果，我们把 DG 概形上的范畴 $\Crys^r$ 与 D-模的（导出）范畴联系起来。

\section{用积分变换表示晶体}

本小节中，设 $X$ 是局部几乎有限型的 DG ind-概形。

\ob{5.1.1}{n:5.1.1}回忆：对 $X\in\DGIndSch_{\mathrm{laft}}$，范畴 $\IndCoh(X)$ 可对偶化，且有典范的自对偶性，见 [GR1, Sect.~2.6]。

于是对 $X,Y\in\DGIndSch$，范畴 $\Funct_{\mathrm{cont}}(\IndCoh(X),\IndCoh(Y))$ 与
\[
\IndCoh(X)^{\vee}\otimes\IndCoh(Y)\simeq\IndCoh(X\times Y)
\]
等同。具体地说，对象 $Q\in\IndCoh(X\times Y)$ 通过
\begin{equation}
F\mapsto (p_2)_*^{\IndCoh}\circ(\Delta_X\times\mathrm{id}_Y)^!(F\boxtimes Q)\tag{5.1}
\end{equation}
定义函子 $F_Q:\IndCoh(X)\to\IndCoh(Y)$，其中 $p_2:X\times Y\to Y$ 是投影映射，$\Delta_X$ 是对角映射 $X\to X\times X$。

特别地，自函子 $\oblv^r_X\circ\indf^r_X$ 定义了一个对象，记为
\[
\mathcal{D}^r_X\in\IndCoh(X\times X).
\]
我们将确定这个对象。

\ob{5.1.2}{n:5.1.2}记 $\widehat{\Delta}_X$ 为映射
\[
X\times_{X_{\dR}}X\simeq(X\times X)^{\wedge}_X\to X\times X.
\]

\begin{othm}{命题}{5.1.3}{n:5.1.3}
在 $\IndCoh(X\times X)$ 中有典范同构
\[
\mathcal{D}^r_X\simeq(\widehat{\Delta}_X)_*^{\IndCoh}\big(\omega_{X\times_{X_{\dR}}X}\big).
\]
\end{othm}

\begin{proof}
我们从下述一般观察开始。

设函子 $F\in\Funct_{\mathrm{cont}}(\IndCoh(X),\IndCoh(Y))$ 由对应给出，即有 DG ind-概形的图表
\[
\begin{tikzcd}[column sep=large, row sep=large]
 & Z \arrow[dl, "q_1"'] \arrow[dr, "q_2"] & \\
 X & & Y
\end{tikzcd}
\]
并设 $F:=(q_2)_*^{\IndCoh}\circ q_1^!$。设
\[
i:Z\to X\times Y
\]
为诱导的乘积映射。

\begin{othm}{引理}{5.1.4}{n:5.1.4}
在上述情形，函子
\[
(q_2)_*^{\IndCoh}\circ q_1^!:\IndCoh(X)\to\IndCoh(Y)
\]
由核 $Q=i_*^{\IndCoh}(\omega_Z)$ 给出。
\end{othm}

\begin{proof}[引理 5.1.4 的证明]
有图表，其内方块是笛卡尔的：
\[
\begin{tikzcd}[column sep=large, row sep=large]
 Z \arrow[r, "q_1\times\mathrm{id}_Z"] \arrow[d, "i"'] & X\times Z \arrow[r] \arrow[d, "\mathrm{id}_X\times i"'] & X \\
 X\times Y \arrow[r, "\Delta_X\times\mathrm{id}_Y"'] & X\times X\times Y \arrow[ur, "p_2"'] &
\end{tikzcd}
\]
其内方块是笛卡尔的。
对 $F\in\IndCoh(X)$，有
\[
(q_2)_*^{\IndCoh}\circ q_1^!(F)\simeq(p_2)_*^{\IndCoh}\circ i_*^{\IndCoh}\circ(q_1\times\mathrm{id}_Z)^!(F\boxtimes\omega_Z).
\]
由 [GR1, Proposition 2.9.2]\footnote{严格地说，[GR1, Proposition 2.9.2] 中的基变换同构只对竖直箭头为 ind-固有的情形陈述，这翻译成 $i$ 固有。在命题 5.1.3 的证明中，我们正是在这样的情形应用它。}，
\[
i_*^{\IndCoh}\circ(q_1\times\mathrm{id}_Z)^!(F\boxtimes\omega_Z)\simeq(\Delta_X\times\mathrm{id}_Y)^!\big(F\boxtimes i_*^{\IndCoh}(\omega_Z)\big).
\]
\end{proof}

我们把这个引理用于命题 5.1.3 的证明：

由命题 3.3.2(b)，函子 $\oblv^r_X\circ\indf^r_X$ 由对应
\[
\begin{tikzcd}[column sep=large, row sep=large]
 & (X\times X)^{\wedge}_X \arrow[dl, "p_s"'] \arrow[dr, "p_t"] & \\
 X & & X
\end{tikzcd}
\]
给出。断言现在由引理 5.1.4 推出。
\end{proof}

\ob{5.1.5}{n:5.1.5}作为命题 5.1.3 的推论，得到：

\begin{othm}{推论}{5.1.6}{n:5.1.6}
存在典范同构 $\sigma(\mathcal{D}^r_X)\simeq\mathcal{D}^r_X$，其中 $\sigma$ 是交换 $X\times X$ 两因子的换位。
\end{othm}

\section{其他函子的显式公式}

本小节中，设 $X$ 是几乎有限型的最终余连接拟紧 DG 概形。

\ob{5.2.1}{n:5.2.1}现在回忆：范畴 $\QCoh(X)$ 也是紧生成且自对偶的。在同构
\[
\QCoh(X)^{\vee}\simeq\QCoh(X),\qquad \IndCoh(X)^{\vee}\simeq\IndCoh(X)
\]
下，函子 $\Upsilon_X$ 的对偶是 [IndCoh, Sect.~1.1.5] 中的函子 $\Psi_X$（对偶性陈述见 [IndCoh, Proposition 9.3.3]）。

特别地，对 $\mathcal{C}'\in\DGCat_{\mathrm{cont}}$，有
\[
\Funct_{\mathrm{cont}}(\QCoh(X),\mathcal{C}')\simeq\QCoh(X)\otimes\mathcal{C}',
\]
其公式与 (5.1) 类似。

\ob{5.2.2}{n:5.2.2}设 $\mathcal{C}$ 是下列范畴之一：
\[
\QCoh(X\times X)\simeq\QCoh(X)\otimes\QCoh(X),\quad \IndCoh(X\times X)\simeq\IndCoh(X)\otimes\IndCoh(X),
\]
\[
\QCoh(X)\otimes\IndCoh(X)\quad\text{或}\quad\IndCoh(X)\otimes\QCoh(X).
\]
则 $\mathcal{C}$ 是 $\QCoh(X\times X)$ 上的模，我们定义 $\mathcal{C}$ 的自函子
\[
F\mapsto F\{X\},
\]
为与对象
\[
\Cone\big(\mathcal{O}_{X\times X}\to j_*j^*(\mathcal{O}_{X\times X})\big)[-1]
\]
作张量积，其中 $j$ 是开浸入 $X\times X-X\hookrightarrow X\times X$。

注意由 [IndCoh, Proposition 4.1.7 and Corollary 4.4.3]，对 $\mathcal{C}=\IndCoh(X\times X)$，这个函子与
\[
(\widehat{\Delta}_X)_*^{\IndCoh}\circ(\widehat{\Delta}_X)^!
\]
等同，其中 $\widehat{\Delta}_X$ 表示映射
\[
X\times_{X_{\dR}}X\simeq(X\times X)^{\wedge}_X\to X\times X.
\]

\ob{5.2.3}{n:5.2.3}我们断言：

\begin{othm}{命题}{5.2.4}{n:5.2.4}
(a) 与 $\oblv^l_X\circ\indf^l_X$ 相对应、属于
\[
\QCoh(X)\otimes\QCoh(X)\simeq\Funct_{\mathrm{cont}}(\QCoh(X),\QCoh(X))
\]
的对象，典范地等同于
\[
(\Psi_X(\omega_X)\boxtimes\mathcal{O}_X)\{X\}.
\]

(b) 与 $\oblv^r_X\circ\Upsilon_{X_{\dR}}\circ\indf^l_X$ 相对应、属于
\[
\QCoh(X)\otimes\IndCoh(X)\simeq\Funct_{\mathrm{cont}}(\QCoh(X),\IndCoh(X))
\]
的对象，典范地等同于
\[
(\Psi_X(\omega_X)\boxtimes\omega_X)\{X\}.
\]

(c) 与 $\oblv^l_X\circ(\Upsilon_{X_{\dR}})^{-1}\circ\indf^r_X$ 相对应、属于
\[
\IndCoh(X)\otimes\QCoh(X)\simeq\Funct_{\mathrm{cont}}(\IndCoh(X),\QCoh(X))
\]
的对象，典范地等同于
\[
(\omega_X\boxtimes\mathcal{O}_X)\{X\}.
\]

(d) 与 $\oblv^r_X\circ{}'\indf^r_X$ 相对应、属于
\[
\QCoh(X)\boxtimes\IndCoh(X)\simeq\Funct_{\mathrm{cont}}(\QCoh(X),\IndCoh(X))
\]
的对象，典范地等同于
\[
(\mathcal{O}_X\boxtimes\omega_X)\{X\}.
\]

(e) 与 ${}'\oblv^r_X\circ\indf^r_X$ 相对应、属于
\[
\IndCoh(X)\boxtimes\QCoh(X)\simeq\Funct_{\mathrm{cont}}(\IndCoh(X),\QCoh(X))
\]
的对象，典范地等同于
\[
(\omega_X\boxtimes\Psi_X(\omega_X))\{X\}.
\]

(f) 与 ${}'\oblv^r_X\circ{}'\indf^r_X$ 相对应、属于
\[
\QCoh(X\times X)\simeq\Funct_{\mathrm{cont}}(\QCoh(X),\QCoh(X))
\]
的对象，典范地等同于
\[
(\mathcal{O}_X\boxtimes\Psi_X(\omega_X))\{X\}.
\]

(g) 与 $\oblv^r_X\circ\Upsilon_{X_{\dR}}\circ\indf^l_X$ 相对应、属于
\[
\QCoh(X\times X)\simeq\Funct_{\mathrm{cont}}(\QCoh(X),\QCoh(X))
\]
的对象，典范地等同于
\[
(\Psi_X(\omega_X)\boxtimes\Psi_X(\omega_X))\{X\}.
\]

(h) 与 $\oblv^l_X\circ(\Upsilon_{X_{\dR}})^{-1}\circ{}'\indf^r_X$ 相对应、属于
\[
\QCoh(X\times X)\simeq\Funct_{\mathrm{cont}}(\QCoh(X),\QCoh(X))
\]
的对象，典范地等同于
\[
(\mathcal{O}_X\boxtimes\mathcal{O}_X)\{X\}.
\]
\end{othm}

\begin{proof}
设 $\mathcal{C},\mathcal{D}$ 是 $\DGCat_{\mathrm{cont}}$ 的对象，$\mathcal{C}$ 可对偶化，于是
\[
\Funct_{\mathrm{cont}}(\mathcal{C},\mathcal{D})\simeq\mathcal{C}^{\vee}\otimes\mathcal{D}.
\]
设 $F:\mathcal{C}_1\to\mathcal{C}$ 与 $G:\mathcal{D}\to\mathcal{D}_1$ 是连续函子。则诱导的函子
\[
\Funct_{\mathrm{cont}}(\mathcal{C},\mathcal{D})\to\Funct_{\mathrm{cont}}(\mathcal{C}_1,\mathcal{D}_1)
\]
由
\[
(F^{\vee}\otimes G):\mathcal{C}^{\vee}\otimes\mathcal{D}\to\mathcal{C}_1^{\vee}\otimes\mathcal{D}_1
\]
给出。

记住这一点，有：

(a) 与 (c) 由命题 5.1.3、引理 3.4.9 以及下述断言结合推出：
\end{proof}

\begin{othm}{引理}{5.2.5}{n:5.2.5}
伴随 $\mathrm{Id}_{\QCoh(X)}\to\Xi^{\vee}_X\circ\Upsilon_X$ 的单位定义同构
\[
\mathcal{O}_X\to\Xi^{\vee}_X(\omega_X).
\]
\end{othm}

(b) 由命题 5.1.3，利用 $\indf^l_X\simeq(\Upsilon_{X_{\dR}})^{-1}\circ\indf^r_X\circ\Upsilon_X$ 推出。

(d) 由命题 5.1.3，利用同构 ${}'\indf^r_X\simeq\indf^r_X\circ\Xi^{\vee}_X$ 与引理 5.2.5 推出。

(e) 由命题 5.1.3 推出。(f) 由 (d) 推出。(g) 由 (b) 推出。

(h) 由 (d) 推出，利用引理 3.4.9 与引理 5.2.5。

\ob{5.2.6}{n:5.2.6}记 $\mathcal{D}^l_X,\mathcal{D}^{l\to r}_X,\mathcal{D}^{r\to l}_X,\mathcal{D}'^r_X$ 为
\[
\QCoh(X)\otimes\QCoh(X)\simeq\QCoh(X\times X)
\]
中分别对应于函子
\[
\oblv^l_X\circ\indf^l_X,\quad {}'\oblv^r_X\circ\Upsilon_{X_{\dR}}\circ\indf^l_X,\quad \oblv^l_X\circ(\Upsilon_{X_{\dR}})^{-1}\circ{}'\indf^r_X,\quad {}'\oblv^r_X\circ{}'\indf^r_X
\]
的对象。

有：

\begin{othm}{命题}{5.2.7}{n:5.2.7}
对象 $\mathcal{D}^r_X,\mathcal{D}^l_X,\mathcal{D}^{l\to r}_X,\mathcal{D}^{r\to l}_X$ 与 $\mathcal{D}'^r_X$ 之间成立下述关系：

(i) $\mathcal{D}^l_X\simeq(\Psi_X\boxtimes\Xi^{\vee}_X)(\mathcal{D}^r_X)\in\QCoh(X\times X)$；

(ii) $(\Psi_X\boxtimes\mathrm{Id}_{\IndCoh(X)})(\mathcal{D}^r_X)\simeq(\mathrm{Id}_{\QCoh(X)}\boxtimes\Upsilon_X)(\mathcal{D}^l_X)\in\QCoh(X)\otimes\IndCoh(X)$；

(iii) $\mathcal{D}'^{l\to r}_X\simeq(\Psi_X\boxtimes\Psi_X)(\mathcal{D}^r_X)\in\QCoh(X\times X)$；

(iii') $\mathcal{D}'^{l\to r}_X\simeq(\mathrm{Id}_{\QCoh(X)}\boxtimes\Psi_X\circ\Upsilon_X)(\mathcal{D}^l_X)\simeq(\mathcal{O}_X\boxtimes\Psi_X(\omega_X))\otimes_{\mathcal{O}_{X\times X}}\mathcal{D}^l_X\in\QCoh(X\times X)$；

(iv) $\mathcal{D}'^{r\to l}_X\simeq(\Xi^{\vee}_X\boxtimes\Xi^{\vee}_X)(\mathcal{D}^r_X)$；

(v) $\mathcal{D}'^r_X\simeq(\Xi^{\vee}_X\boxtimes\Psi_X)(\mathcal{D}^r_X)\in\QCoh(X\times X)$。
\end{othm}

\begin{proof}
(i) 由引理 3.4.9 推出。

(ii) 由函子的（同义反复的）同构
\[
\oblv^r_X\circ\indf^r_X\circ\Upsilon_X\simeq\oblv^r_X\circ\Upsilon_{X_{\dR}}\circ\indf^l_X\simeq\Upsilon_X\circ\oblv^l_X\circ\indf^l_X
\]
推出。

(iii) 是同义反复的。

(iii') 由 (ii) 与 (iii) 结合推出。

(iv) 由引理 3.4.9 推出。

(v) 是同义反复的。
\end{proof}

\section{关于 t-结构的行为}

继续假设 $X$ 是几乎有限型的拟紧 DG 概形。

\ob{5.3.1}{n:5.3.1}我们注意到：

\begin{othm}{引理}{5.3.2}{n:5.3.2}
对象 $\mathcal{D}^r_X$ 下有界，即属于 $\IndCoh(X\times X)^{+}$。
\end{othm}

\begin{proof}
由命题 5.1.3，利用 $\omega_X\in\IndCoh(X)^{+}$ 以及函子
\[
F\mapsto F\{X\},\qquad \IndCoh(X\times X)\to\IndCoh(X\times X)
\]
右 t-正合推出。
\end{proof}

\ob{5.3.3}{n:5.3.3}现在假设 $X$ 最终余连接。我们断言：

\begin{othm}{命题}{5.3.4}{n:5.3.4}
$\QCoh(X\times X)$ 的对象 $\mathcal{D}^l_X,\mathcal{D}'^{l\to r}_X,\mathcal{D}'^{r\to l}_X,\mathcal{D}'^r_X$ 都最终余连接，即属于 $\QCoh(X\times X)^{+}$。
\end{othm}

\begin{proof}
由命题 5.2.4，利用 $\Psi_X(\omega_X),\mathcal{O}_X\in\QCoh(X)^{+}$ 以及函子
\[
F\mapsto F\{X\},\qquad \QCoh(X\times X)\to\QCoh(X\times X)
\]
右 t-正合推出。
\end{proof}

\ob{5.3.5}{n:5.3.5}最后，假设 $X$ 是光滑经典概形。我们断言：

\begin{othm}{命题}{5.3.6}{n:5.3.6}
对象 $\mathcal{D}^l_X\in\QCoh(X\times X)$ 落在 t-结构的心中。
\end{othm}

\begin{proof}
断言是 Zariski 局部的，因此不妨设 $X$ 仿射。只需证明
\[
(p_2)_*(\mathcal{D}^l_X)\in\QCoh(X)
\]
落在 t-结构的心中。有
\[
(p_2)_*(\mathcal{D}^l_X)\simeq\oblv^l_X\circ\indf^l_X(\mathcal{O}_X)\simeq\Xi^{\vee}_X\circ(\oblv^r_X\circ\indf^r_X)\circ\Upsilon_X(\mathcal{O}_X).
\]
现在，函子 $\oblv^r_X\circ\indf^r_X$ t-正合（见命题 4.2.11），函子 $\Upsilon_X$ 是平移次数 $[n]$ 的等价，而 $\Xi^{\vee}_X$ 是 $\Upsilon_X$ 的逆。
\end{proof}

\section{与微分算子层的关系}

本小节中，取 $X$ 为光滑经典概形。我们将把 $\mathcal{D}^l_X$ 与经典微分算子层 $\mathrm{Diff}_X$ 底层的 $\QCoh(X\times X)$ 中的对象等同起来。

\ob{5.4.1}{n:5.4.1}对任意 $Q\in\QCoh(X\times X)^{\heartsuit}$（它在集合论意义下支撑在对角线上）以及 $F_1,F_2\in\QCoh(X)^{\heartsuit}$，一个映射
\[
(p_2)_*(p_1^*(F_1)\otimes Q)\to F_2
\]
的数据，等价于一个映射
\[
Q\to\mathrm{Diff}_X(F_1,F_2)
\]
的数据。

此外，这个对应与 $\QCoh(X\times X)^{\heartsuit}$ 上的单子结构相容，该结构由卷积与微分算子的复合给出。

\ob{5.4.2}{n:5.4.2}取 $Q=\mathcal{D}^l_X$，$F_1=F_2=\mathcal{O}_X$，由单子 $\oblv^l_X\circ\indf^l_X$ 在 $\mathcal{O}_X$ 上的作用，得到所需的映射
\begin{equation}
\mathcal{D}^l_X\to\mathrm{Diff}_X,\tag{5.2}
\end{equation}
它与代数结构相容。

我们断言：

\begin{othm}{引理}{5.4.3}{n:5.4.3}
映射 (5.2) 是代数的同构。
\end{othm}

\begin{proof}
只需证明 (5.2) 在 $\QCoh(X\times X)$ 的底层对象层面是同构。后者例如由命题 5.2.4 给出的 $\mathcal{D}^l_X$ 作为拟凝聚层的描述推出。
\end{proof}

\section{晶体与 D-模的关系}

设 $X$ 是有限型经典概形。我们将证明：范畴 $\Crys^r(X)$ 可以典范地等同于 $X$ 上右 D-模的（导出）范畴 $\mathrm{D\text{-}mod}^r(X)$。

\begin{othm}{注}{5.5.1}{n:5.5.1}
范畴 $\mathrm{D\text{-}mod}^r(X)$ 满足 Zariski 下降。因此，由命题 3.2.2，下文只需对仿射概形建立典范等价。
\end{othm}

\ob{5.5.2}{n:5.5.2}设 $Z$ 是光滑经典仿射概形，$i:X\hookrightarrow Z$ 是闭浸入。由经典 Kashiwara 引理与命题 2.5.6，为构造等价
\[
\Crys^r(X)\simeq\mathrm{D\text{-}mod}^r(X),
\]
只需对 $Z$ 构造。

于是可以假设 $X$ 本身是光滑经典仿射概形。我们将构造所求的等价，连同函子的交换图
\[
\begin{tikzcd}[column sep=large, row sep=large]
 \Crys^r(X) \arrow[r] \arrow[d, "\oblv^r_X"'] & \mathrm{D\text{-}mod}^r(X) \arrow[d] \\
 \IndCoh(X) \arrow[r, "\Psi_X"'] & \QCoh(X),
\end{tikzcd}
\]
其中右竖直箭头是自然的遗忘函子，函子 $\Psi_X$ 是 [IndCoh, Lemma 1.1.6] 中的等价。

由命题 2.4.4，按上述方式构造等价 $\Crys^r(X)\simeq\mathrm{D\text{-}mod}^r(X)$，等同于构造左晶体与左 D-模之间的等价，连同函子的交换图
\begin{equation}
\begin{tikzcd}[column sep=large, row sep=large]
 \Crys^l(X) \arrow[rr] \arrow[dr, "\oblv^l_X"'] & & \mathrm{D\text{-}mod}^l(X) \arrow[dl] \\
 & \QCoh(X) &
\end{tikzcd}\tag{5.3}
\end{equation}

\ob{5.5.3}{n:5.5.3}由命题 4.7.3、2.4.4 与 4.4.2，范畴 $\Crys^l(X)$ 与其 t-结构心的导出范畴等同。范畴 $\mathrm{D\text{-}mod}^l(X)$ 按定义是 $\mathrm{D\text{-}mod}^l(X)^{\heartsuit}$ 的导出范畴。此外，图 (5.3) 中的竖直箭头 t-正合。

因此，只需在相应的 Abel 范畴层面构造所求的等价
\begin{equation}
\begin{tikzcd}[column sep=large, row sep=large]
 \Crys^l(X)^{\heartsuit} \arrow[r] \arrow[d, "\oblv^l_X"'] & \mathrm{D\text{-}mod}^l(X)^{\heartsuit} \arrow[d, "\mathrm{Id}"] \\
 \QCoh(X)^{\heartsuit} \arrow[r] & \QCoh(X)^{\heartsuit}.
\end{tikzcd}\tag{5.4}
\end{equation}

\ob{5.5.4}{n:5.5.4}后者是 Grothendieck 的经典计算：

具体地说，把 $\Crys^l(X)^{\heartsuit}$ 解释为截断单纯对象
\[
(X\times X\times X)^{\wedge}_X\substack{p_{12}\\[-2pt]p_{13}\\[-2pt]p_{23}}\rightrightarrows(X\times X)^{\wedge}_X\substack{p_1\\[-2pt]p_2}\rightrightarrows X
\]
上的拟凝聚层范畴的心。也就是说，具体地，$\Crys^l(X)^{\heartsuit}$ 的对象是拟凝聚层 $F\in\QCoh(X)^{\heartsuit}$ 连同同构
\[
\varphi:p_2^*(F)\xrightarrow{\sim}p_1^*(F),
\]
它在对角线上限制为恒等，并满足余循环条件
\[
p_{13}^*(\varphi)=p_{12}^*(\varphi)\circ p_{23}^*(\varphi).
\]
下面给出建立 (5.4) 中等价的另一种途径。

\ob{5.5.5}{n:5.5.5}Abel 范畴 $\Crys^l(X)^{\heartsuit}$ 与 $\mathrm{D\text{-}mod}^l(X)^{\heartsuit}$ 分别实现为作用在 $\QCoh(X)^{\heartsuit}$ 上的单子 $M_{\Crys^l(X)}$ 与 $M_{\mathrm{D\text{-}mod}(X)}$ 的模。

按定义，$M_{\mathrm{D\text{-}mod}(X)}$ 由微分算子代数 $\mathrm{Diff}_X$ 给出。单子 $M_{\Crys^l(X)}$ 由 $\oblv^l_X\circ\indf^l_X$ 给出。于是所求等价由引理 5.4.3 推出。

\begin{othm}{注}{5.5.6}{n:5.5.6}
由构造可知，等价
\[
\Crys^l(X)\to\mathrm{D\text{-}mod}^l(X)
\]
与光滑经典概形间映射 $f:Y\to X$ 的拉回相容。
\end{othm}

\begin{annote}
这个等价说明晶体确实推广了 D-模：Kashiwara 引理在晶体的定义中是自动的，而在经典 D-模理论中是需要证明的定理；此外晶体的定义不依赖 $X$ 的光滑性，也不依赖 $X$ 的经典性。\hfill【来源：原书 §0.1.1、§2.5、§5.5（A）】
\end{annote}


\chapter{扭化}
\markboth{第 6 章\ 扭化}{第 6 章\ 扭化}

本节中，我们不假设所考虑的预叠与 DG 概形局部几乎有限型。我们将在 §6.7 恢复这个假设。

\section{Gerbe}

\ob{6.1.1}{n:6.1.1}记 $\mathrm{pt}/G_m$ 为群 $G_m$ 的分类叠。换言之，$\mathrm{pt}/G_m$ 是表示下述函子的代数叠：它把仿射 DG 概形 $S$ 对应到 $S$ 上线丛的 $\infty$-群胚。

事实上，由于 $G_m$ 是 Abel 群，叠 $\mathrm{pt}/G_m$ 有自然的 Abel 群结构。$\mathrm{pt}/G_m$ 上的乘法映射表示线丛的张量积。这个结构把 $\mathrm{pt}/G_m$ 升级为从仿射 DG 概形到 $\infty$-Picard 群胚（即连接谱）的函子。

对我们的目的而言，$G_m$-gerbe 是 $\mathrm{pt}/G_m$-torsor 的预层 $G$，它满足下述三个（互不等价的）条件之一：

(i) $G$ 在平展拓扑下局部非空\footnote{由 Toën 定理，这等价于在 fppf 拓扑下局部非空。}。

(ii) $G$ 在 Zariski 拓扑下局部非空。

(iii) $G$ 整体非空。

具体地说，记 $B^{\mathrm{naive}}(\mathrm{pt}/G_m)$ 为 $\mathrm{pt}/G_m$ 的分类预叠。它由单纯预叠
\[
B^{\mathrm{naive}}(\mathrm{pt}/G_m):=\cdots\mathrm{pt}/G_m\times\mathrm{pt}/G_m\rightrightarrows\mathrm{pt}/G_m\to\mathrm{pt}
\]
的几何实现给出。记 $B^{\mathrm{Zar}}(\mathrm{pt}/G_m)$（相应地，$B^{\mathrm{et}}(\mathrm{pt}/G_m)$）为预叠 $B^{\mathrm{naive}}(\mathrm{pt}/G_m)$ 的 Zariski（相应地，平展）层化。

预叠 $B^{\mathrm{et}}(\mathrm{pt}/G_m)$、$B^{\mathrm{Zar}}(\mathrm{pt}/G_m)$ 与 $B^{\mathrm{naive}}(\mathrm{pt}/G_m)$ 分别表示满足上述条件 (i)、(ii)、(iii) 的 $G_m$-gerbe。

记 $(\mathbf{Ge}_{G_m})_{\DGSch_{\mathrm{aff}}}$ 为函子
\[
(\DGSch_{\mathrm{aff}})^{\mathrm{op}}\to\infty\text{-}\mathrm{PicGrpd},
\]
它把仿射 DG 概形 $S$ 对应到 $G_m$-gerbe 的群胚，其中 gerbe 取上述三种概念中的任何一种。

\begin{othm}{注}{6.1.2}{n:6.1.2}
虽然这三个版本并不给出等价的 $G_m$-gerbe 概念，但我们很快会看到，它们导致相同的扭化定义，因为相关的 gerbe 将是那些限制到 $\clred S$ 上平凡化的 gerbe。
\end{othm}

\ob{6.1.3}{n:6.1.3}定义函子
\[
(\mathbf{Ge}_{G_m})_{\PreStk}:(\PreStk)^{\mathrm{op}}\to\infty\text{-}\mathrm{PicGrpd}
\]
为 $(\mathbf{Ge}_{G_m})_{\DGSch_{\mathrm{aff}}}$ 沿
\[
(\DGSch_{\mathrm{aff}})^{\mathrm{op}}\hookrightarrow(\PreStk)^{\mathrm{op}}
\]
的右 Kan 延拓。也就是说，对 $Y\in\PreStk$，
\[
\mathbf{Ge}_{G_m}(Y):=\varprojlim_{S\in(\DGSch_{\mathrm{aff}}/Y)^{\mathrm{op}}}\mathbf{Ge}_{G_m}(S).
\]
等价地，
\[
\mathbf{Ge}_{G_m}(Y)=\Maps_{\PreStk}\big(Y,B^{?}(\mathrm{pt}/G_m)\big)
\]
其中 $?=\mathrm{naive},\mathrm{Zar}$ 或 $\mathrm{et}$。

因此，粗略地说，$Y$ 上的 $G_m$-gerbe 是对每个映射到 $Y$ 的 $S\in\DGSch_{\mathrm{aff}}$ 指定 $S$ 上的一个 $G_m$-gerbe，并对 $S$ 函子性地相容。

对子范畴 $\mathcal{C}\subset\PreStk$，记 $(\mathbf{Ge}_{G_m})_{\mathcal{C}}$ 为限制 $(\mathbf{Ge}_{G_m})_{\PreStk}|_{\mathcal{C}}$。

\section{扭化的概念}

\ob{6.2.1}{n:6.2.1}设 $Y$ 是预叠。$Y$ 上扭化的 Picard 群胚定义为
\[
\Tw(Y):=\ker\big(p^*_{\dR,Y}:\mathbf{Ge}_{G_m}(Y_{\dR})\to\mathbf{Ge}_{G_m}(Y)\big),
\]
其中 $\mathbf{Ge}_{G_m}$ 按 naive、Zar 或 et 三种版本中的任何一种理解。我们很快会看到（见 §6.4），三个版本都等价。

粗略地说，$Y$ 上的扭化 $T$ 是如下数据：对每个带有映射 $\clred S\to Y$ 的 $S\in\DGSch_{\mathrm{aff}}$，指定一个对象 $G_S\in\mathbf{Ge}_{G_m}(S)$，它在映射 $S_1\to S_2$ 下相容地变化。此外，对上述映射到映射 $S\to Y$ 的每个延拓，指定 $G_S$ 的一个平凡化，它同样对映射 $S_1\to S_2$ 函子性地变化。

\begin{othm}{注}{6.2.2}{n:6.2.2}
当我们写 $\ker(A_1\to A_2)$（其中 $A_1\to A_2$ 是 $\infty$-PicGrpd 中的映射）时，指的是
\[
A_1\times_{A_2}\{*\},
\]
其中纤维积在 $\infty$-PicGrpd 中取。也就是说，这等于在所有（即不必连接的）谱的范畴中取的纤维积的连接截断。
\end{othm}

\ob{6.2.3}{n:6.2.3}例。设 $L$ 是 $Y$ 上的线丛。如下定义 $Y$ 上的扭化 $T(L)$：它对每个带有映射 $\clred S\to Y$ 的 $S\in\DGSch_{\mathrm{aff}}$ 指定平凡的 $G_m$-gerbe。对映射 $S\to Y$，把上述 gerbe 的平凡化乘以 $L$ 来平凡化它。

\ob{6.2.4}{n:6.2.4}显然扭化构成函子
\[
\Tw_{\PreStk}:\PreStk^{\mathrm{op}}\to\infty\text{-}\mathrm{PicGrpd}.
\]
对态射 $f:Y_1\to Y_2$，记 $f^*$ 为相应的函子
\[
\Tw(Y_2)\to\Tw(Y_1).
\]
若 $\mathcal{C}$ 是 $\PreStk$ 的子范畴（例如 $\mathcal{C}=\DGSch_{\mathrm{aff}}$ 或 $\DGSch$），记 $\Tw_{\mathcal{C}}$ 为 $\Tw_{\PreStk}$ 到 $\mathcal{C}^{\mathrm{op}}$ 的限制。

\ob{6.2.5}{n:6.2.5}按构造，函子 $\Tw_{\PreStk}$ 把 $\PreStk$ 中的余极限送到 $\infty$-PicGrpd 中的极限。于是由推论 1.1.5 得到：

\begin{othm}{引理}{6.2.6}{n:6.2.6}
函子 $\Tw_{\PreStk}$ 同构于 $\Tw_{\mathcal{C}}$ 沿
\[
\mathcal{C}^{\mathrm{op}}\hookrightarrow\PreStk^{\mathrm{op}}
\]
的右 Kan 延拓，其中 $\mathcal{C}$ 是下列范畴之一：
\[
\DGSch_{\mathrm{aff}},\qquad \DGSch_{\mathrm{qs\text{-}qc}},\qquad \DGSch.
\]
\end{othm}

具体地说，这个引理说明映射
\[
\Tw(Y)\to\varprojlim_{S\in(\DGSch_{\mathrm{aff}}/Y)^{\mathrm{op}}}\Tw(S)
\]
是同构（并且 $\DGSch_{\mathrm{aff}}$ 可以换成 $\DGSch_{\mathrm{qs\text{-}qc}}$ 或 $\DGSch$）。

粗略地说，这意味着：在预叠 $Y$ 上指定一个扭化，等价于对每个映射到 $Y$ 的仿射 DG 概形 $S$ 指定一族相容的扭化。

\section{变体：其他结构群}

\ob{6.3.1}{n:6.3.1}设 $S$ 是仿射 DG 概形。考虑 Picard 群胚
\[
\mathbf{Ge}^{/\mathrm{red}}_{G_m}(S):=\ker\big(\mathbf{Ge}_{G_m}(S)\to\mathbf{Ge}_{G_m}(\clred S)\big).
\]
记 $(\mathbf{Ge}^{/\mathrm{red}}_{G_m})_{\DGSch_{\mathrm{aff}}}$ 为所得的函子
\[
(\DGSch)^{\mathrm{op}}\to\infty\text{-}\mathrm{PicGrpd}.
\]

\ob{6.3.2}{n:6.3.2}按定义，可以把 $\mathbf{Ge}^{/\mathrm{red}}_{G_m}(S)$ 看作关于 Abel 群预层
\[
(\mathcal{O}^{\times})^{/\mathrm{red}}_S:\ker(\mathcal{O}^{\times}_S\to\mathcal{O}^{\times}_{\clred S})
\]
的 gerbe（取 §6.1.1 三种版本中的任何一种）。

\ob{6.3.3}{n:6.3.3}除了 $G_m$-gerbe，还可以考虑 $G_a$-gerbe。有函子
\[
(\mathbf{Ge}_{G_a})_{\DGSch_{\mathrm{aff}}}:(\DGSch_{\mathrm{aff}})^{\mathrm{op}}\to\infty\text{-}\mathrm{PicGrpd},
\]
它把仿射 DG 概形 $S$ 对应到 $S$ 上 $G_a$-gerbe 的群胚。

注意与 $G_m$-gerbe 的情形不同，§6.1.1 中讨论的三种 gerbe 概念对 $G_a$-gerbe 是等价的。这是由于对仿射 DG 概形 $S$，
\[
H^2_{\mathrm{Zar}}(S,G_a)=H^2_{\mathrm{et}}(S,G_a)=0.
\]
于是 $(\mathbf{Ge}_{G_a})_{\DGSch_{\mathrm{aff}}}$ 由 $B^2(G_a)$ 表示，后者是相应单纯预叠的几何实现。

\ob{6.3.4}{n:6.3.4}按定义，对仿射 DG 概形 $S$，有
\[
\mathbf{Ge}_{G_a}(S)=B^2\big(\Maps(S,G_a)\big)\simeq B^2\big(\Gamma(S,\mathcal{O}_S)\big).
\]
特别地，看作连接谱时，$\mathbf{Ge}_{G_a}(S)$ 有基域 $k$ 上模的自然结构。这把 $(\mathbf{Ge}_{G_a})_{\DGSch_{\mathrm{aff}}}$ 升级为函子
\[
(\DGSch)^{\mathrm{op}}\to\infty\text{-}\mathrm{PicGrpd}_k,
\]
其中 $\infty\text{-}\mathrm{PicGrpd}_k$ 表示连接谱中 $k$-模范畴。注意由 Dold--Kan 对应，
\[
\infty\text{-}\mathrm{PicGrpd}_k\simeq\Vect^{\leq 0}.
\]
定义函子
\[
(\mathbf{Ge}_{G_a})_{\PreStk}:\PreStk^{\mathrm{op}}\to\infty\text{-}\mathrm{PicGrpd}_k
\]
为函子 $(\mathbf{Ge}_{G_a})_{\DGSch_{\mathrm{aff}}}$ 沿
\[
(\DGSch_{\mathrm{aff}})^{\mathrm{op}}\hookrightarrow\PreStk^{\mathrm{op}}
\]
的右 Kan 延拓。

\ob{6.3.5}{n:6.3.5}如同 $G_m$-gerbe 的情形，可以考虑 Picard 群胚
\[
\mathbf{Ge}^{/\mathrm{red}}_{G_a}(S):=\ker\big(\mathbf{Ge}_{G_a}(S)\to\mathbf{Ge}_{G_a}(\clred S)\big),
\]
并记 $(\mathbf{Ge}^{/\mathrm{red}}_{G_a})_{\DGSch_{\mathrm{aff}}}$ 为所得的函子
\[
(\DGSch)^{\mathrm{op}}\to\infty\text{-}\mathrm{PicGrpd}_k.
\]
按定义，对仿射 DG 概形 $S$，$\mathbf{Ge}^{/\mathrm{red}}_{G_a}(S)$ 是关于连接谱预层
\[
\mathcal{O}^{/\mathrm{red}}_S:=\ker(\mathcal{O}_S\to\mathcal{O}_{\clred S})
\]
的 gerbe。具体地，
\[
\mathbf{Ge}^{/\mathrm{red}}_{G_a}(S)\simeq B^2\big(\Gamma(S,\mathcal{O}^{/\mathrm{red}}_S)\big).
\]

\ob{6.3.6}{n:6.3.6}回忆 [GR1, Sect.~6.8.8]：指数映射定义同构
\[
\exp:\mathcal{O}^{/\mathrm{red}}_S\to(\mathcal{O}^{\times})^{/\mathrm{red}}_S.
\]
于是得到：

\begin{othm}{推论}{6.3.7}{n:6.3.7}
指数映射定义函子的同构
\begin{equation}
\exp:(\mathbf{Ge}_{G_a})^{/\mathrm{red}}_{\DGSch_{\mathrm{aff}}}\to(\mathbf{Ge}_{G_m})^{/\mathrm{red}}_{\DGSch_{\mathrm{aff}}}\tag{6.1}
\end{equation}
对 $(\mathbf{Ge}_{G_m})^{/\mathrm{red}}_{\DGSch_{\mathrm{aff}}}$ 的三种版本（naive、Zar 或 et）中的任何一种都成立。
\end{othm}

于是，若把 $\mathbf{Ge}^{/\mathrm{red}}_{G_m}(S)$ 实现为平展或 Zariski 拓扑中的 gerbe，则这个范畴有平凡的 $\pi_0$ 与 $\pi_1$。换言之，仿射 DG 概形上任何这样的 gerbe 都整体非空，且任何自同构都非典范地同构于恒等。

\ob{6.3.8}{n:6.3.8}同构 (6.1) 赋予 $\mathbf{Ge}^{/\mathrm{red}}_{G_m}(S)$（看作连接谱）基域 $k$ 上模的结构。这把 $(\mathbf{Ge}^{/\mathrm{red}}_{G_m})_{\DGSch_{\mathrm{aff}}}$ 升级为函子
\[
(\DGSch)^{\mathrm{op}}\to\infty\text{-}\mathrm{PicGrpd}_k.
\]
定义函子
\[
(\mathbf{Ge}^{/\mathrm{red}}_{G_m})_{\PreStk}:\PreStk^{\mathrm{op}}\to\infty\text{-}\mathrm{PicGrpd}_k
\]
为函子 $(\mathbf{Ge}^{/\mathrm{red}}_{G_m})_{\DGSch_{\mathrm{aff}}}$ 沿
\[
(\DGSch_{\mathrm{aff}})^{\mathrm{op}}\hookrightarrow\PreStk^{\mathrm{op}}
\]
的右 Kan 延拓。

\ob{6.3.9}{n:6.3.9}按定义，对 $Y\in\PreStk$，
\[
\mathbf{Ge}^{/\mathrm{red}}_{G_m}(Y):=\varprojlim_{S\in(\DGSch_{\mathrm{aff}}/Y)^{\mathrm{op}}}\mathbf{Ge}^{/\mathrm{red}}_{G_m}(S).
\]
粗略地说，对 $Y\in\PreStk$，对象 $G\in\mathbf{Ge}^{/\mathrm{red}}_{G_m}(Y)$ 是对每个 $S\in\DGSch_{\mathrm{aff}}/Y$ 指定一个对象 $G_S\in\mathbf{Ge}^{/\mathrm{red}}_{G_m}(S)$，并对每个 $S'\to S$ 指定同构
\[
G_{S'}\simeq f^*(G_S).
\]
下述结果由定义推出：

\begin{othm}{引理}{6.3.10}{n:6.3.10}
对 $Y\in\PreStk$，自然映射
\[
\mathbf{Ge}^{/\mathrm{red}}_{G_m}(Y)\to\ker\big(\mathbf{Ge}_{G_m}(Y)\to\mathbf{Ge}_{G_m}(\clred Y)\big)
\]
是同构，其中
\[
\clred Y:=\mathrm{LKE}_{(\mathrm{red}\,\Sch_{\mathrm{aff}})^{\mathrm{op}}\hookrightarrow(\PreStk)^{\mathrm{op}}}(Y|_{\mathrm{red}\,\Sch_{\mathrm{aff}}}).
\]
\end{othm}

\section{扭化：重新表述}

我们将证明：扭化的概念可以用
\[
(\mathbf{Ge}^{/\mathrm{red}}_{G_m})_{\PreStk},\qquad(\mathbf{Ge}_{G_a})_{\PreStk},\qquad(\mathbf{Ge}^{/\mathrm{red}}_{G_a})_{\PreStk}
\]
代替 $(\mathbf{Ge}_{G_m})_{\PreStk}$ 来表述。

\ob{6.4.1}{n:6.4.1}考虑函子
\[
\Tw^{/\mathrm{red}},\ \Tw_a,\ \Tw^{/\mathrm{red}}_a:\PreStk^{\mathrm{op}}\to\infty\text{-}\mathrm{PicGrpd},
\]
它们分别由
\[
\Tw^{/\mathrm{red}}(Y):=\ker\big(p^*_{\dR,Y}:\mathbf{Ge}^{/\mathrm{red}}_{G_m}(Y_{\dR})\to\mathbf{Ge}^{/\mathrm{red}}_{G_m}(Y)\big),
\]
\[
\Tw_a(Y):=\ker\big(p^*_{\dR,Y}:\mathbf{Ge}_{G_a}(Y_{\dR})\to\mathbf{Ge}_{G_a}(Y)\big)
\]
以及
\[
\Tw^{/\mathrm{red}}_a(Y):=\ker\big(p^*_{\dR,Y}:\mathbf{Ge}^{/\mathrm{red}}_{G_a}(Y_{\dR})\to\mathbf{Ge}^{/\mathrm{red}}_{G_a}(Y)\big)
\]
给出。

有下述由指数映射与显然的遗忘函子给出的函子图：
\begin{equation}
\begin{tikzcd}[column sep=large, row sep=large]
 & \Tw^{/\mathrm{red}}_a \arrow[r, "\exp"] \arrow[d] & \Tw^{/\mathrm{red}} \arrow[d] \\
 & \Tw_a \arrow[r] & \Tw
\end{tikzcd}\tag{6.2}
\end{equation}

\begin{othm}{命题}{6.4.2}{n:6.4.2}
(6.2) 中的函子都是等价。
\end{othm}

\begin{proof}
由指数映射给出的函子是等价，见 §6.3.6。

来证明 (6.2) 中右竖直映射是等价。这其实是同义反复：

两个函子都是沿
\[
(\DGSch_{\mathrm{aff}})^{\mathrm{op}}\to(\PreStk)^{\mathrm{op}}
\]
的右 Kan 延拓，因此只需证明所讨论的映射在对象 $S\in\DGSch_{\mathrm{aff}}$ 上取值时为同构。

有：
\[
\Tw^{/\mathrm{red}}(S)=\mathbf{Ge}^{/\mathrm{red}}_{G_m}(S_{\dR})\times_{\mathbf{Ge}^{/\mathrm{red}}_{G_m}(S)}\{*\}\overset{\text{引理 6.3.10}}{\simeq}
\]
\[
\simeq\ker\big(\mathbf{Ge}_{G_m}(S_{\dR})\to\mathbf{Ge}_{G_m}(\clred(S_{\dR}))\big)\times_{\ker(\mathbf{Ge}_{G_m}(S)\to\mathbf{Ge}_{G_m}(\clred S))}\{*\}=
\]
\[
=\ker\big(\mathbf{Ge}_{G_m}(S_{\dR})\to\mathbf{Ge}_{G_m}(\clred S)\big)\times_{\ker(\mathbf{Ge}_{G_m}(S)\to\mathbf{Ge}_{G_m}(\clred S))}\{*\}\simeq
\]
\[
\simeq\mathbf{Ge}_{G_m}(S_{\dR})\times_{\mathbf{Ge}_{G_m}(S)}\{*\}=\Tw(S).
\]
(6.2) 中左竖直箭头是等价，证明类似。
\end{proof}

\ob{6.4.3}{n:6.4.3}作为命题 6.4.2 的推论，得到：

\begin{othm}{推论}{6.4.4}{n:6.4.4}
扭化的三种版本——naive、Zar 与 et——等价。
\end{othm}

此外：

\begin{othm}{推论}{6.4.5}{n:6.4.5}
函子 $\Tw:(\PreStk)^{\mathrm{op}}\to\infty\text{-}\Grpd$ 典范地升级为函子
\[
(\PreStk)^{\mathrm{op}}\to\infty\text{-}\mathrm{PicGrpd}_k.
\]
\end{othm}

\ob{6.4.6}{n:6.4.6}例。可以利用 $\Tw$ 上自然的 $k$-模结构来产生更多的扭化例子。设 $L$ 是 $Y$ 上的线丛，$T(L)$ 是 §6.2.3 中的扭化。现在，对 $a\in k$，$\Tw(Y)$ 上的 $k$-模结构给出新的扭化 $T(L^{\otimes a})$。

\begin{othm}{注}{6.4.7}{n:6.4.7}
注意：仿射 DG 概形 $X$ 上的扭化 $T$ 未必平凡，即使在 Zariski 或平展拓扑局部也未必。对任意带有映射 $S\to X_{\dR}$ 的 $S\in\DGSch_{\mathrm{aff}}$，相应的 $S$ 上 $G_m$-gerbe 可以非典范地平凡化；但这种平凡化不一定能对不同选择的 $S$ 相容地选取。

对光滑经典 $X$，这样的 gerbe 的例子可以由一个非平展局部正合的闭 $2$-形式给出（见 §6.5.4）。

然而注意，例 6.4.6 中描述的 gerbe 是 Zariski 局部平凡的，因为线丛具有相应的性质。
\end{othm}

\ob{6.4.8}{n:6.4.8}收敛性。现在断言：

\begin{othm}{命题}{6.4.9}{n:6.4.9}
函子 $\Tw:(\DGSch_{\mathrm{aff}})^{\mathrm{op}}\to\infty\text{-}\mathrm{PicGrpd}$ 收敛\footnote{收敛的概念见 §1.3.1。}。
\end{othm}

\begin{proof}
我们将证明函子 $\Tw^{/\mathrm{red}}$ 收敛。为此，只需证明函子
\[
S\mapsto\mathbf{Ge}^{/\mathrm{red}}_{G_m}(S_{\dR})\quad\text{与}\quad\mathbf{Ge}^{/\mathrm{red}}_{G_m}(S)
\]
收敛。

$\mathbf{Ge}^{/\mathrm{red}}_{G_m}((-)_{\dR})$ 的收敛性是显然的，因为这个函子只依赖于底层的约化经典概形。于是只需证明 $\mathbf{Ge}^{/\mathrm{red}}_{G_m}(-)$ 的收敛性。

有：
\[
\mathbf{Ge}^{/\mathrm{red}}_{G_m}(S)=\mathbf{Ge}_{G_m}(S)\times_{\mathbf{Ge}_{G_m}(\clred S)}\{*\}.
\]
因此只需证明函子 $\mathbf{Ge}_{G_m}(-)$ 收敛。后者由
\[
\mathbf{Ge}_{G_m}(-)=B^2\big(\Maps(S,G_m)\big)
\]
以及 $G_m$ 作为 DG 概形收敛推出。
\end{proof}

可以把命题 6.4.9 同义反复地重新表述为：

\begin{othm}{推论}{6.4.10}{n:6.4.10}
函子 $\Tw_{\PreStk}$ 同构于
\[
\Tw_{{}^{<\infty}\DGSch_{\mathrm{aff}}}:=\Tw_{\DGSch_{\mathrm{aff}}}|_{{}^{<\infty}\DGSch_{\mathrm{aff}}}
\]
沿
\[
({}^{<\infty}\DGSch_{\mathrm{aff}})^{\mathrm{op}}\hookrightarrow(\DGSch_{\mathrm{aff}})^{\mathrm{op}}\hookrightarrow(\PreStk)^{\mathrm{op}}
\]
的右 Kan 延拓。
\end{othm}

\begin{othm}{注}{6.4.11}{n:6.4.11}
可以利用命题 6.4.9 证明函子 $\mathbf{Ge}_{G_m}$（三种版本中任何一种）也收敛。
\end{othm}

\ob{6.4.12}{n:6.4.12}局部几乎有限型情形下的扭化。推论 6.4.10 蕴含：为知道 $Y\in\PreStk$ 上的扭化（当 $Y$ 局部几乎有限型时），我们「不需要知道」非局部几乎有限型的 DG 概形。

\begin{othm}{推论}{6.4.13}{n:6.4.13}
(a) 对 $Y\in\PreStk_{\mathrm{laft}}$，自然定义的映射
\[
\Tw(Y)\to\varprojlim_{S\in(({}^{<\infty}\DGSch_{\mathrm{aff}}^{\mathrm{aft}})/Y)^{\mathrm{op}}}\Tw(S)
\]
是等价。

(b) 函子 $\Tw_{\PreStk_{\mathrm{laft}}}$ 同构于 $\Tw_{{}^{<\infty}\DGSch_{\mathrm{aff}}^{\mathrm{aft}}}$ 沿包含
\[
({}^{<\infty}\DGSch_{\mathrm{aff}}^{\mathrm{aft}})^{\mathrm{op}}\hookrightarrow(\DGSch_{\mathrm{aff}}^{\mathrm{aft}})^{\mathrm{op}}\hookrightarrow(\PreStk_{\mathrm{laft}})^{\mathrm{op}}
\]
的右 Kan 延拓。
\end{othm}

\begin{proof}
把 $\Tw$ 换成任意收敛预叠 $(\DGSch_{\mathrm{aff}})^{\mathrm{op}}\to\infty\text{-}\Grpd$ 时，这都是对的。
\end{proof}

\begin{othm}{注}{6.4.14}{n:6.4.14}
由注 6.4.11 推出：函子 $\mathbf{Ge}_{G_m}$（三种版本中任何一种）看作预层时属于 $\PreStk_{\mathrm{laft}}$。事实上，在 naive 版本中这是显然的，因为 $\mathrm{pt}/G_m$ 属于 $\PreStk_{\mathrm{laft}}$。对 Zar 与 et 版本，这由 [GL:Stacks, Corollary 2.5.7] 推出：该结论说，在 $n$-连接的预叠语境中，「局部有限型」这一条件在层化后仍保持，只要我们限制到截断预叠。
\end{othm}

\section{扭化 Picard 群胚的确定}

可以利用用 $G_a$-gerbe 对扭化的描述，给出扭化群胚的上同调描述。

\ob{6.5.1}{n:6.5.1}de Rham 上同调。设 $Y$ 是预叠。回忆 $Y$ 的凝聚上同调定义为
\[
H(Y):=\Gamma(Y,\mathcal{O}_Y)=\Maps_{\QCoh(Y)}(\mathcal{O}_Y,\mathcal{O}_Y).
\]
定义 $Y$ 的 de Rham 上同调为 $Y_{\dR}$ 的凝聚上同调，即
\[
H_{\dR}(Y):=H(Y_{\dR})=\Maps_{\QCoh(Y_{\dR})}(\mathcal{O}_{Y_{\dR}},\mathcal{O}_{Y_{\dR}}).
\]
注意由于 $\QCoh(Y_{\dR})$ 是稳定 $\infty$-范畴，上述 $\Maps$ 给出一个（不必连接的）谱。

设 $X$ 是光滑经典概形。此时由 §5.5，有
\[
H_{\dR}(X)=\Maps_{\mathrm{D\text{-}mod}^l(X)}(\mathcal{O}_X,\mathcal{O}_X).
\]
特别地，我们对 de Rham 上同调的定义与光滑经典概形的通常定义一致。

\ob{6.5.2}{n:6.5.2}考虑表示 $G_a$-gerbe 的函子 $B^2(G_a)$。按定义，对预叠 $Y$，有连接谱的同构：
\[
\Maps(Y,B^2(G_a))\simeq\tau^{\leq 0}\Maps_{\QCoh(Y)}(\mathcal{O}_Y,\mathcal{O}_Y)[2]\simeq\tau^{\leq 0}(H(Y)[2]).
\]
于是由命题 6.4.2 得到：

\begin{othm}{推论}{6.5.3}{n:6.5.3}
对预叠 $Y$，扭化的群胚由
\[
\Tw(Y)\simeq\tau^{\leq 2}\Big(H_{\dR}(Y)\times_{H(Y)}\{*\}\Big)[2]
\]
给出。
\end{othm}

\ob{6.5.4}{n:6.5.4}现在设 $X$ 是光滑经典概形。此时有
\[
H_{\dR}(X)\simeq\Gamma(X,\Omega^\bullet),
\]
其中 $\Omega^\bullet$ 是 $X$ 上的 de Rham 微分复形。自然映射
\[
H_{\dR}(X)\to H(X)
\]
由投影映射 $\Omega^\bullet\to\mathcal{O}_X$ 的整体截面给出。因此有
\[
\Tw(X)\simeq\tau^{\leq 2}\Big(H_{\dR}(X)\times_{H(X)}\{*\}\Big)[2]\simeq\tau^{\leq 2}\Gamma\big(X,\Omega^\bullet\times_{\mathcal{O}_X}0\big)[2]\simeq
\]
\[
\simeq\tau^{\leq 2}\Gamma\big(X,\tau^{\leq 2}(\Omega^\bullet\times_{\mathcal{O}_X}0)\big)[2].
\]
复形 $\tau^{\leq 2}(\Omega^\bullet\times_{\mathcal{O}_X}0)$ 与复形
\[
\Omega^1\to\Omega^{2,\mathrm{cl}}
\]
等同，其中 $\Omega^{2,\mathrm{cl}}$ 是闭 $2$-形式的层（放在上同调次数 $2$），映射是 de Rham 微分。

于是得到：$X$ 上扭化的 Picard 群胚由
\[
\Tw(X)\simeq\tau^{\leq 2}\Gamma(X,\Omega^1\to\Omega^{2,\mathrm{cl}})[2]
\]
给出。特别地，对光滑经典概形，我们的扭化定义与 [BB] 中 TDO 的概念一致。

\begin{annote}
这与 [BB] 的经典结论一致：光滑经典 $X$ 上扭化的 Picard 群胚是 $\tau^{\leq 2}\Gamma(X,\Omega^1\to\Omega^{2,\mathrm{cl}})[2]$，即扭化的对象由闭 $2$-形式给出、同构由 $1$-形式给出；相应的扭 D-模范畴就是 TDO 上的模范畴（§7.5.1）。\hfill【来源：原书 §0.4.2、§6.5.4；[BB]（A）】
\end{annote}


\section{扭化与无穷小群胚}

\ob{6.6.1}{n:6.6.1}设
\[
Y_1\rightrightarrows Y_0
\]
是 $\PreStk$ 中的群胚对象，$Y_\bullet$ 为相应的单纯对象。回忆这个群胚对象被 $G_m$ 中心扩张的概念。（这里 $G_m$ 可以换成任意交换群对象 $H\in\PreStk$。）

按定义，$Y_1\rightrightarrows Y_0$ 被 $G_m$ 的中心扩张是 $\mathbf{Ge}_{G_m}(|Y_\bullet|)$ 的一个对象，连同一个其沿
\[
Y_0\to|Y_\bullet|
\]
的限制的平凡化。

\ob{6.6.2}{n:6.6.2}粗略地说，这样的中心扩张的数据是：$Y_1$ 上的线丛 $L$，其沿退化映射 $Y_0\to Y_1$ 的拉回被平凡化，并且对三个映射
\[
p_{1,2},\ p_{2,3},\ p_{1,3}:Y_2\to Y_1,
\]
给定同构
\[
p_{1,2}^*(L)\otimes p_{2,3}^*(L)\simeq p_{1,3}^*(L),
\]
并满足进一步的融贯条件。

\ob{6.6.3}{n:6.6.3}对 $Y\in\PreStk$，考虑其无穷小群胚
\[
Y\times_{Y_{\dR}}Y\rightrightarrows Y.
\]
按定义，$Y$ 上的扭化给出其无穷小群胚被 $G_m$ 的中心扩张。

反之，由引理 1.2.4 得到：

\begin{othm}{推论}{6.6.4}{n:6.6.4}
假设 $Y$ 经典形式光滑。则上述函子
\[
\Tw(Y)\to\{Y\text{ 的无穷小群胚被 }G_m\text{ 的中心扩张}\}
\]
是等价。
\end{othm}

\section{ind-概形上的扭化}

\ob{6.7.1}{n:6.7.1}设 $X$ 是 $\DGIndSch_{\mathrm{laft}}$ 的对象。我们将证明推论 6.6.4 的断言对 $X$ 成立：

\begin{othm}{命题}{6.7.2}{n:6.7.2}
函子
\[
\Tw(X)\to\{X\text{ 的无穷小群胚被 }G_m\text{ 的中心扩张}\}
\]
是等价。
\end{othm}

本小节余下部分专门证明命题 6.7.2。

\ob{6.7.3}{n:6.7.3}第 1 步。由推论 6.4.13(b)，需要证明下述事情：

对每个 $S\in(\DGSch_{\mathrm{aff}}^{\mathrm{aft}})/X_{\dR}$，在 $S$ 上给定一个 $G_m$-gerbe 并给定其拉回到 $S\times_{X_{\dR}}X$ 的平凡化，等价于在单纯预叠
\[
S^\bullet:=S\times_{X_{\dR}}(X^\bullet/X_{\dR})
\]
上给定一个 $G_m$-gerbe 并给定其在 $0$-单形上（即 $S\times_{X_{\dR}}X$ 上）的平凡化。

\ob{6.7.4}{n:6.7.4}第 2 步。注意单纯预叠 $\clred(S^\bullet)$ 是常值的，值为 $\clred S$。于是由引理 6.3.7，可以用 $G_a$-gerbe 代替 $G_m$-gerbe。

因此只需证明映射
\[
\Maps_{\QCoh(S)}(\mathcal{O}_S,\mathcal{O}_S)\to\Tot\Big(\Maps_{\QCoh(S^\bullet)}(\mathcal{O}_{S^\bullet},\mathcal{O}_{S^\bullet})\Big)
\]
在 $\Vect$ 中是同构。

\ob{6.7.5}{n:6.7.5}第 3 步。注意对任意 $X'\in\DGIndSch_{\mathrm{laft}}$，典范映射
\[
\Maps_{\QCoh(X')}(\mathcal{O}_{X'},\mathcal{O}_{X'})\to\Maps_{\IndCoh(X')}(\omega_{X'},\omega_{X'})
\]
是同构。这例如由 DG 概形的相应断言，即引理 5.2.5 推出。

因此只需证明映射
\[
\Maps_{\IndCoh(S)}(\omega_S,\omega_S)\to\Tot\Big(\Maps_{\IndCoh(S^\bullet)}(\omega_{S^\bullet},\omega_{S^\bullet})\Big)
\]
是同构。

\ob{6.7.6}{n:6.7.6}第 4 步。注意 $S^\bullet$ 与映射
\begin{equation}
S\times_{X_{\dR}}X\to S\tag{6.3}
\end{equation}
的 Čech 神经等同。

如同命题 3.1.3 的证明，所有 $S^i$ 都属于 $\DGIndSch$，且态射 (6.3) 是 ind-固有的且满。

于是所求断言由 $\IndCoh$ 在 DG ind-概形的 ind-固有满映射下的下降推出，见 [GR1, Lemma 2.10.3]。

\chapter{扭晶体}
\markboth{第 7 章\ 扭晶体}{第 7 章\ 扭晶体}

本节将说明扭化的数据如何用来修改左晶体与右晶体的范畴。主要结论是：「其实没有太大变化」。

\section{扭左晶体}

本小节中，我们不假设 DG 概形与预栈局部几乎有限型。

\ob{7.1.1}{n:7.1.1}设 $Y$ 是预栈。考虑范畴 $\PreStk_{/Y}$ 以及函子
\[
\QCoh_{\DGSch_{\mathrm{aff}}/Y}:(\DGSch_{\mathrm{aff}}/Y)^{\mathrm{op}}\to\DGCat_{\mathrm{cont}}.
\]
群栈 $\mathrm{pt}/G_m$ 通过用线丛作张量积作用于 $\QCoh$。

设 $G$ 是 $Y$ 上的 $G_m$-gerbe。则 $G$ 通过 $\mathrm{pt}/G_m$ 对 $\QCoh$ 的作用给出函子 $\QCoh_{\DGSch_{\mathrm{aff}}/Y}$ 的一个扭版本。这定义了函子
\[
\QCoh^G_{\DGSch_{\mathrm{aff}}/Y}:(\DGSch_{\mathrm{aff}}/Y)^{\mathrm{op}}\to\DGCat_{\mathrm{cont}}.
\]

\ob{7.1.2}{n:7.1.2}特别地，若 $T$ 是 $Y$ 上的扭化，则得到函子
\[
\QCoh^T_{\DGSch_{\mathrm{aff}}/Y_{\dR}}:(\DGSch_{\mathrm{aff}}/Y_{\dR})^{\mathrm{op}}\to\DGCat_{\mathrm{cont}}.
\]
记 $\QCoh^T_{\DGSch_{\mathrm{aff}}/Y}$ 为它沿映射
\[
(\DGSch_{\mathrm{aff}}/Y)^{\mathrm{op}}\to(\DGSch_{\mathrm{aff}}/Y_{\dR})^{\mathrm{op}}
\]
的限制。按构造，$\QCoh^T_{\DGSch_{\mathrm{aff}}/Y}$ 典范同构于 $\QCoh_{\DGSch_{\mathrm{aff}}/Y}$。

\ob{7.1.3}{n:7.1.3}更一般地，可以考虑函子
\[
\QCoh^T_{\PreStk_{/Y_{\dR}}}:(\PreStk_{/Y_{\dR}})^{\mathrm{op}}\to\DGCat_{\mathrm{cont}},
\]
它是 $\QCoh^T_{\DGSch_{\mathrm{aff}}/Y_{\dR}}$ 沿
\[
(\DGSch_{\mathrm{aff}}/Y_{\dR})^{\mathrm{op}}\hookrightarrow(\PreStk_{/Y_{\dR}})^{\mathrm{op}}
\]
的右 Kan 延拓。

$\QCoh^T_{\PreStk_{/Y_{\dR}}}$ 沿
\[
\PreStk_{/Y}\to\PreStk_{/Y_{\dR}}
\]
的限制典范同构于 $\QCoh_{\PreStk_{/Y}}$。

\ob{7.1.4}{n:7.1.4}对预栈 $Y$ 上的扭化 $T$，$Y$ 上 $T$-扭左晶体的范畴定义为
\[
\Crys^{T,l}(Y):=\QCoh^T(Y_{\dR}).
\]
具体地，有
\[
\Crys^{T,l}(Y)=\varprojlim_{S\in(\DGSch_{\mathrm{aff}}/Y_{\dR})^{\mathrm{op}}}\QCoh^T(S).
\]

\ob{7.1.5}{n:7.1.5}更一般地，定义函子
\[
\Crys^{T,l}_{\PreStk_{/Y_{\dR}}}:(\PreStk_{/Y_{\dR}})^{\mathrm{op}}\to\DGCat_{\mathrm{cont}}
\]
为复合 $\QCoh^T_{\PreStk_{/Y_{\dR}}}\circ\dR$。推论 2.1.4 的类比对这个函子成立。

\ob{7.1.6}{n:7.1.6}有典范的自然变换
\[
\oblv^{T,l}:\Crys^{T,l}_{\PreStk_{/Y_{\dR}}}\to\QCoh^T_{\PreStk_{/Y_{\dR}}}.
\]
对单个 $Y'\in\PreStk_{/Y_{\dR}}$，记所得的函子
\[
\Crys^{T,l}(Y')\to\QCoh^T(Y')
\]
为 $\oblv^{T,l}(Y')$。

\ob{7.1.7}{n:7.1.7}记 $\Crys^{T,l}_{\PreStk_{/Y}}$ 为 $\Crys^{T,l}_{\PreStk_{/Y_{\dR}}}$ 沿 $\PreStk_{/Y}\to\PreStk_{/Y_{\dR}}$ 的限制。

按轻微的记号滥用，我们用同一个符号 $\oblv^{T,l}$ 表示所得的自然变换
\[
\Crys^{T,l}_{\PreStk_{/Y}}\to\QCoh_{\PreStk_{/Y}}.
\]

\section{扭右晶体}

从此刻起，我们恢复所有 DG 概形与预栈局部几乎有限型的假设，并保持到全文结束。

\ob{7.2.1}{n:7.2.1}设 $Y$ 是 $\PreStk_{\mathrm{laft}}$ 的对象，$G$ 是 $Y$ 上的 $G_m$-gerbe。

$\QCoh_{\PreStk_{\mathrm{laft}}}$ 对 $\IndCoh_{\PreStk_{\mathrm{laft}}}$ 的作用（见 [IndCoh, Sect.~10.3]）可以定义函子
\[
\IndCoh^G_{(\PreStk_{\mathrm{laft}})_{/Y}}:((\PreStk_{\mathrm{laft}})_{/Y})^{\mathrm{op}}\to\DGCat_{\mathrm{cont}},
\]
其性质与
\[
\IndCoh_{(\PreStk_{\mathrm{laft}})_{/Y}}:=\IndCoh_{\PreStk_{\mathrm{laft}}}|_{(\PreStk_{\mathrm{laft}})_{/Y}}
\]
的性质类似。

\ob{7.2.2}{n:7.2.2}特别地，对 $T\in\Tw(Y)$，有函子
\[
\Crys^{T,r}_{(\PreStk_{\mathrm{laft}})_{/Y_{\dR}}}:((\PreStk_{\mathrm{laft}})_{/Y_{\dR}})^{\mathrm{op}}\to\DGCat_{\mathrm{cont}},
\]
以及自然变换
\[
\Crys^{T,r}_{(\PreStk_{\mathrm{laft}})_{/Y_{\dR}}}\to\IndCoh^T_{(\PreStk_{\mathrm{laft}})_{/Y_{\dR}}}\quad\text{与}\quad\Crys^{T,r}_{(\PreStk_{\mathrm{laft}})_{/Y_{\dR}}}\to\IndCoh_{(\PreStk_{\mathrm{laft}})_{/Y_{\dR}}}.
\]
推论 2.3.7 与 2.3.9 以及引理 2.3.11 与 2.3.12 的类比对 $Y'\in(\PreStk_{\mathrm{laft}})_{/Y_{\dR}}$ 成立，证明相同。

\section{扭晶体的性质}

如上所述，所有 DG 概形与预栈都假定局部几乎有限型。

固定 $Y\in\PreStk_{\mathrm{laft}}$ 与 $T\in\Tw(Y)$。

\begin{othm}{注}{7.3.1}{n:7.3.1}
一般说来，关于晶体的结论不会自动对扭晶体成立。在我们的某些证明中，需要把给定的仿射 DG 概形 $X$ 嵌入光滑经典概形 $Z$。在扭晶体的情形，问题在于我们可能找不到这样的 $Z$，使它同时映射到 $Y$（乃至 $Y_{\dR}$）。

然而，有一大类例子（涵盖了迄今为止应用中出现的所有情形），其中结论的推广是自动的：即当 $T$ 限制到任意 $S\in(\DGSch_{\mathrm{aff}}^{\mathrm{aft}})_{/Y}$ 上都在 Zariski 或平展拓扑下局部平凡时（另见注 6.4.7）。形如 $L^{\otimes a}$（$L\in\mathrm{Pic}(Y)$，$a\in k$）的扭化及其张量积就是这种情形。
\end{othm}

\ob{7.3.2}{n:7.3.2}推论 2.2.2 与 2.2.4 以及引理 2.1.7 的类比对扭左晶体成立，证明相同。

此外，Kashiwara 引理对左、右扭晶体都成立，证明也相同。

最后注意，存在典范的自然变换
\begin{equation}
\Upsilon:\Crys^{T,l}_{(\PreStk_{\mathrm{laft}})_{/Y_{\dR}}}\to\Crys^{T,r}_{(\PreStk_{\mathrm{laft}})_{/Y_{\dR}}}.\tag{7.1}
\end{equation}

\begin{othm}{命题}{7.3.3}{n:7.3.3}
自然变换 (7.1) 是等价。
\end{othm}

\begin{proof}
论证与命题 2.4.4 相同：

我们不需要光滑经典概形 $Z$ 映射到 $Y$。相反，我们利用下述事实：若 $Y$ 是光滑经典概形 $Z$ 沿 Zariski 闭子集的形式完备化，$G$ 是 $Y$ 上的 $G_m$-gerbe 且在 $\clred Y$ 上平凡，则函子
\[
\Upsilon_Y:\QCoh^G(Y)\to\IndCoh^G(Y)
\]
是等价。后者由非扭情形中的相应事实（在命题 2.4.4 的证明过程中已证）推出，因为 $G$ 是（非典范地）平凡的。
\end{proof}

作为推论，我们得到：引理 2.2.6 的类比在扭情形也成立。

\ob{7.3.4}{n:7.3.4}于是，对 $Y'\in(\PreStk_{\mathrm{laft}})_{/Y_{\dR}}$，可以把 $Y'$ 上的晶体看作一个范畴 $\Crys^T(Y')$，它带有两个遗忘函子
\begin{equation}
\begin{tikzcd}[column sep=6em, row sep=4em]
 & \Crys^T(Y') \arrow[dl, "\oblv^{T,l}_{Y'}"'{above left}] \arrow[dr, "\oblv^{T,r}_{Y'}"{above right}] & \\
 \QCoh^T(Y') \arrow[rr, "\Upsilon_{Y'}"'] & & \IndCoh^T(Y')
\end{tikzcd}\tag{7.2}
\end{equation}

对 $Y'\in(\PreStk_{\mathrm{laft}})_{/Y}$，上述遗忘函子落到非扭层上：
\begin{equation}
\begin{tikzcd}[column sep=6em, row sep=4em]
 & \Crys^T(Y') \arrow[dl, "\oblv^{T,l}_{Y'}"'{above left}] \arrow[dr, "\oblv^{T,r}_{Y'}"{above right}] & \\
 \QCoh(Y') \arrow[rr, "\Upsilon_{Y'}"'] & & \IndCoh(Y')
\end{tikzcd}\tag{7.3}
\end{equation}

\ob{7.3.5}{n:7.3.5}设 $X\in(\DGIndSch_{\mathrm{laft}})_{/Y_{\dR}}$。命题 3.1.3 的类比原封不动地成立。特别地，得到函子
\[
\indf^{T,r}_X:\IndCoh^T(X)\to\Crys^{T,r}(X),
\]
它是 $\oblv^{T,r}_X$ 的左伴随。

类似地，命题 3.4.3 的类比在当前语境下也成立。

\ob{7.3.6}{n:7.3.6}下述观察在后续中有用：

设 $X$ 是 $Y_{\dR}$ 上的仿射 DG 概形（或 ind-仿射 DG ind-概形）。选取 $X$ 上相应 $G_m$-gerbe 的一个平凡化。这个选择给出等同
\[
\alpha:\IndCoh^T(X)\simeq\IndCoh(X).
\]

\begin{othm}{引理}{7.3.7}{n:7.3.7}
单子 $\oblv^{T,r}_X\circ\indf^{T,r}_X$，看作函子（忘掉单子结构）
\[
\IndCoh(X)\overset{\alpha^{-1}}{\simeq}\IndCoh^T(X)\to\IndCoh^T(X)\overset{\alpha}{\simeq}\IndCoh(X),
\]
非典范同构于 $\oblv^r_X\circ\indf^r_X$。
\end{othm}

\begin{proof}
首先注意，命题 5.1.3 的类比成立；即，定义函子 $\oblv^{T,r}_X\circ\indf^{T,r}_X$ 的 $\IndCoh(X\times X)$ 中对象由
\[
(\widehat{\Delta}_X)_*^{\IndCoh}\big(L\otimes\omega_{X\times_{X_{\dR}}X}\big)
\]
给出，其中 $L$ 是 $X\times_{X_{\dR}}X$ 上对应于 $T$ 与 $\alpha$ 的线丛，如 §6.6.1 所述。

按构造，$L$ 限制到 $X\hookrightarrow X\times_{X_{\dR}}X$ 上时平凡。于是，由于 $X$ 仿射，$L$ 可以在整个 $X\times_{X_{\dR}}X$ 上平凡化。
\end{proof}

\section{扭晶体上的 t-结构}

与上一小节一样，固定 $Y\in\PreStk_{\mathrm{laft}}$ 与 $T\in\Tw(Y)$。

\ob{7.4.1}{n:7.4.1}若 $X$ 是 DG 概形，$G$ 是其上的 $G_m$-gerbe，则扭范畴 $\QCoh^G(X)$ 与 $\IndCoh^G(X)$ 带有自然的 t-结构，其性质与通常的 $\QCoh(X)$、$\IndCoh(X)$ 的相应性质类似。

特别地，对任意 $Y'\in(\PreStk_{\mathrm{laft}})_{/Y_{\dR}}$，$\Crys^{T,l}(Y')$ 上有「左」t-结构。（这个 t-结构的定义不需要 $Y$ 或 $Y'$ 局部几乎有限型的假设。）

DG 概形上扭 $\IndCoh$ 的 t-结构允许我们为 DG ind-概形 $X$ 定义 $\IndCoh^G(X)$ 上的 t-结构。于是可以定义范畴 $\Crys^{T,r}(X)$ 上的「右」t-结构。

\ob{7.4.2}{n:7.4.2}注意命题 4.2.5 原封不动地移植到扭情形。现在断言：

\begin{othm}{命题}{7.4.3}{n:7.4.3}
设 $X$ 是映射到 $Y_{\dR}$ 的拟紧 DG 概形。

(a) 函子 $\indf^{T,r}_X$ 是 t-正合的。

(b) 对拟紧概形 $X$，函子 $\oblv^{T,r}_X$ 有界上同调幅度。
\end{othm}

\begin{proof}
函子 $\indf^{T,r}_X$ 右 t-正合，因为其右伴随 $\oblv^{T,r}_X$ 左 t-正合。按「右」t-结构的定义，$\indf^{T,r}_X$ 的左 t-正合性等价于复合 $\oblv^{T,r}_X\circ\indf^{T,r}_X$ 的同一性质。

断言是 Zariski 局部的，故可设 $X$ 仿射。此时 $\oblv^{T,r}_X\circ\indf^{T,r}_X$ 左 t-正合由引理 7.3.7 以及非扭情形中相应断言成立推出。

由于 $\IndCoh^T(X)^{\leq 0}$ 经函子 $\indf^{T,r}_X$ 生成 $\Crys^{T,l}(X)^{\leq 0}$，为证 $\oblv^{T,r}_X$ 的上同调幅度上有界，只需对 $\oblv^{T,r}_X\circ\indf^{T,r}_X$ 证明同样的事。这又由引理 7.3.7 推出。
\end{proof}

\ob{7.4.4}{n:7.4.4}现在断言：

\begin{othm}{命题}{7.4.5}{n:7.4.5}
设 $X$ 是映射到 $Y_{\dR}$ 的拟紧 DG 概形。

(a) $\Crys^T(X)$ 上的「左」t-结构与「右」t-结构相差有限上同调幅度。

(b) 函子 $\oblv^{T,l}_X:\Crys^T(X)\to\QCoh(X)$ 有界上同调幅度。\footnote{由 (a)，这一断言不依赖于我们在 $\Crys^T(X)$ 上考虑两个 t-结构中的哪一个。}
\end{othm}

\ob{7.4.6}{n:7.4.6}先证明下述结论：

设 $i:X\to Z$ 是闭浸入，$Z$ 是光滑经典概形。设 $Y$ 是 $Z$ 沿 $X$ 的形式完备化。

\begin{othm}{引理}{7.4.7}{n:7.4.7}
函子
\[
\oblv^{T,r}_Y:\Crys^{T,r}(Y)\to\IndCoh^T(Y)
\]
是 t-正合的。
\end{othm}

\begin{proof}
与命题 7.4.3 的证明一样，只需证明函子
\[
\oblv^{T,r}_Y\circ\indf^{T,r}_Y:\IndCoh^T(Y)\to\IndCoh^T(Y)
\]
是 t-正合的。断言是 Zariski 局部的，故可设 $X$ 仿射。此时与命题 7.4.3 的证明一样，所论函子非典范同构于非扭版本 $\oblv^r_Y\circ\indf^r_Y$，而后者由命题 4.2.11(a) 是 t-正合的。
\end{proof}

\ob{7.4.8}{n:7.4.8}命题 7.4.5 的证明。断言是 Zariski 局部的，故可设 $X$ 仿射，并把它嵌入光滑经典概形 $Z$。记 $Y$ 为 $X$ 在 $Z$ 中的形式完备化。按定义，$T$ 定义 $Y$ 上的一个 $G_m$-gerbe $G$。记 $i':X\to Y$ 为相应的映射。

为证 (a)，由引理 4.3.7（其第二个证明在扭情形同样适用），可以把 $X$ 换成 $Y$，于是只需证明 $\Crys^T(Y)$ 上两个 t-结构之间的差异是有限的。由命题 4.1.3（在扭情形应用）与引理 7.4.7，只需证明函子
\[
\Psi_Y:\QCoh^T(Y)\to\IndCoh^T(Y)
\]
有界上同调幅度。这等价于
\[
\Psi_Y:\QCoh(Y)\to\IndCoh(Y)
\]
的相应事实，而后者又由 $Z$ 的相应事实推出。

(b) 由函子
\[
i'^*:\QCoh^T(Y)\to\QCoh^T(X)
\]
有界幅度推出，这又等价于
\[
i'^*:\QCoh(Y)\to\QCoh(X)
\]
的相应事实，而后者由 $Z$ 的相应事实推出。\qed

\ob{7.4.9}{n:7.4.9}§4.6 中为非扭晶体建立的关于「粗糙的」遗忘函子与归纳函子的结论，自动移植到扭情形。

\ob{7.4.10}{n:7.4.10}当前的目标是证明：

\begin{othm}{命题}{7.4.11}{n:7.4.11}
设 $X$ 是映射到 $Y_{\dR}$ 的拟紧 DG 概形。

(a) $\Crys^{T,r}(X)$ 上的「右」t-结构是左完备的。

(b) 对仿射 $X$，自然函子 $D(\Crys^{T,r}(X)^\heartsuit)^+\to\Crys^{T,r}(X)^+$（心取关于「右」t-结构）唯一延拓为等价
\[
D(\Crys^{T,r}(X)^\heartsuit)\to\Crys^{T,r}(X).
\]
\end{othm}

\ob{7.4.12}{n:7.4.12}命题 7.4.11(a) 的证明。同样，断言是 Zariski 局部的，我们沿用命题 7.4.5 证明中的设置。

只需给出一族对象
\[
P_\alpha\in\Crys^{T,r}(X),
\]
它们生成 $\Crys^{T,r}(X)$ 且有界 Ext 维数，即：对每个 $\alpha$，存在整数 $k_\alpha$，使得
\[
\Hom_{\Crys^T(X)}(P_\alpha,M)=0\quad\text{若 }M\in\Crys^T(X)^{<-k_\alpha}.
\]

把 $\Crys^{T,r}(X)$ 实现为 $\Crys^{T,r}(Y)$。由引理 4.3.7 与 7.4.7，$\Crys^{T,r}(X)\simeq\Crys^{T,r}(Y)$ 上的 t-结构由下述性质刻画：
\[
M\in\Crys^{T,r}(Y)^{\geq 0}\iff\oblv^{T,r}_Y(M)\in\IndCoh^T(Y)^{\geq 0}.
\]

取 $P_\alpha$ 形如 $\indf^{T,r}(F)$，其中 $F\in\Coh^T(Y)^\heartsuit$。为证所需的 Ext 消没，需要证明：对 $M\in\Crys^{T,r}(Y)^{\ll 0}$，
\[
\Hom_{\IndCoh^T(Y)}\big(F,\oblv^{T,r}_Y(M)\big)=0.
\]
然而这由范畴 $\IndCoh^T(Y)$ 关于其 t-结构有有限上同调维数推出\footnote{关于「有限上同调维数」的含义，见注 4.7.8 的脚注。}：事实上，所论范畴非典范等价于 $\IndCoh(Y)$，而后者的上同调维数被 $\IndCoh(Z)$ 的上同调维数界住。

\ob{7.4.13}{n:7.4.13}命题 7.4.11(b) 的证明。沿用 (a) 的证明中的记号。

与命题 4.7.3 的证明一样，在 (a) 已证的基础上，只需验证：对 $M_1,M_2\in\Crys^T(X)^\heartsuit$ 与任意 $k\geq 0$，映射
\[
\Ext^k_{\Crys^T(X)^\heartsuit}(M_1,M_2)\to\Hom_{\Crys^T(X)}(M_1,M_2[k])
\]
是同构。

为此只需证明：范畴 $\Crys^{T,r}(X)^\heartsuit$ 含有 $\Crys^{T,r}(X)$ 的一个 pro-投射生成元，即存在一族带满射的滤过逆向族 $P_\alpha\in\Crys^{T,r}(X)^\heartsuit$，使得函子
\[
\varinjlim_\alpha\Maps_{\Crys^{T,r}(X)}(P_\alpha,-)
\]
在 $\Crys^{T,r}(X)$ 上 t-正合且保守。

取
\[
P_\alpha=\indf^{T,r}_Y(\mathcal{O}_{X_n})\in\Crys^{T,r}(Y)^\heartsuit\simeq\Crys^{T,r}(X)^\heartsuit,
\]
其中 $X_n$ 是 $\clred X$ 在 $Z$ 中的第 $n$ 个无穷小邻域。\qed

\section{其他结果}

\ob{7.5.1}{n:7.5.1}扭晶体与扭 D-模。设 $X$ 是光滑经典概形。§6.5.4 已说明：$X$ 上扭化的 Picard 范畴等价于 $X$ 上 TDO 的 Picard 范畴。

给定扭化 $T$ 与相应的 TDO（记作 $\mathrm{Diff}^T_X$），存在典范等价
\[
\Crys^{T,l}(X)\simeq\mathrm{D}\text{-}\mathrm{mod}^{T,l}(X),
\]
它与到 $\QCoh(X)$ 的遗忘函子交换；扭右晶体有类似结论。证明或是对 §5.5.4 中所述策略的细化，或是利用 §5.4。

\ob{7.5.2}{n:7.5.2}扭 D-模与 TDO 上的模之间的关系可以推广到下述情形：用光滑经典概形 $Z$ 中 DG 概形 $X$ 的形式完备化 $Y$ 代替光滑经典概形 $X$。

这允许证明：

\begin{othm}{命题}{7.5.3}{n:7.5.3}
设 $X$ 是 $Y_{\dR}$ 上的拟紧 DG 概形。

(a) Abel 范畴 $\Crys^{T,r}(X)^\heartsuit$ 是局部 Noether 的。

(b) $\Crys^{T,r}(X)$ 关于其 t-结构有有限上同调维数。
\end{othm}

（关于命题 (a)、(b) 中断言的性质的含义，见注 4.7.8 的脚注。）


\backmatter
\chapter*{参考文献}
\addcontentsline{toc}{chapter}{参考文献}
\markboth{参考文献}{参考文献}

\begin{thebibliography}{[Si]}
\bibitem[BB]{BB} A.~Beilinson and J.~Bernstein, \emph{A proof of Jantzen conjectures}, I.~M.~Gelfand Seminar, 1--50, Adv. Sov. Math.~16, Part 1, AMS, 1993.

\bibitem[GL:DG]{GL:DG} \emph{Notes on Geometric Langlands, Generalities on DG categories}, available at \url{http://www.math.harvard.edu/~gaitsgde/GL/}.

\bibitem[GL:Stacks]{GL:Stacks} \emph{Notes on Geometric Langlands, Stacks}, available at \url{http://www.math.harvard.edu/~gaitsgde/GL/}.

\bibitem[GL:QCoh]{GL:QCoh} \emph{Notes on Geometric Langlands, Quasi-coherent sheaves on stacks}, available at \url{http://www.math.harvard.edu/~gaitsgde/GL/}.

\bibitem[IndCoh]{IndCoh} D.~Gaitsgory, \emph{Ind-coherent sheaves}, arXiv:1105.4857.

\bibitem[GR1]{GR1} D.~Gaitsgory and N.~Rozenblyum, \emph{DG indschemes}, arXiv:1108.1738.

\bibitem[GR2]{GR2} D.~Gaitsgory and N.~Rozenblyum, \emph{A study in derived algebraic geometry}, in preparation, preliminary version will gradually become available at \url{http://www.math.harvard.edu/~gaitsgde/GL/}.

\bibitem[Gr]{Gr} A.~Grothendieck, \emph{Crystals and the de Rham cohomology of schemes}, Dix expos\'es sur la cohomologie des sch\'emas, 306--358.

\bibitem[Lu0]{Lu0} J.~Lurie, \emph{Higher Topos Theory}, Princeton Univ. Press (2009).

\bibitem[Lu1]{Lu1} J.~Lurie, \emph{DAG-VIII}, available at \url{http://www.math.harvard.edu/~lurie}.

\bibitem[Lu2]{Lu2} J.~Lurie, \emph{Higher algebra}, available at \url{http://www.math.harvard.edu/~lurie}.

\bibitem[Si]{Si} C.~Simpson, \emph{Homotopy over the complex numbers and generalized cohomology theory}, Moduli of vector bundles (Taniguchi Symposium, December 1994), M.~Maruyama ed., Dekker Publ, 229--263, 1996.
\end{thebibliography}


\end{document}
